\documentclass[../main.tex]{subfiles}
\begin{document}
%==========================================TIÊU ĐỀ TẬP========================================
\begin{table}[h]
  \begin{adjustwidth}{-1cm}{}
    \begin{tabular}{>{\centering\arraybackslash}p{6cm}|>{\raggedright\arraybackslash}p{12.5cm}}
      \multirow{5}{6cm}
      % Cột bên trái: ÉPISODE và số tập
      {\\[-40pt]
        \flaregothic\fontsize{20pt}{20pt}\selectfont \centering
        {\contour{errortwo}{\textcolor{black}{ÉPISODE}}}\color{black}\\
        \burbank\fontsize{72pt}{72pt}\selectfont
        \includegraphics[height=3cm]{../../../../Image/Book?/number/num6.png}
        \\[18pt]
        % \flaregothic\fontsize{14pt}{14pt}\selectfont
        % \color{epattr}BIENVENUE, \uline{\color{Banpire} Mel} !
        % \flaregothic\fontsize{18pt}{18pt}\selectfont
        % \color{epattr}ÉPISODE PERDU
        \color{black}
      }
       &
      % Dòng thứ nhất: tên tập tiếng Nhật
      {\fotskip\fontsize{18pt}{18pt}\selectfont \ruby{憧}{あこが}れの\ruby{彼女}{かのじょ}
      }    \\[6pt]
      % Dòng thứ hai: tên tập tiếng Anh
       & {\cafeteria\fontsize{18pt}{18pt}\selectfont The Role Model}             \\[4pt]
      % Dòng thứ ba: tên tập tiếng Việt
       & {\vnmsans\fontsize{18pt}{18pt}\selectfont Kẻ noi gương}                 \\[8pt]
      % Dòng thứ tư: Nhân vật xuất hiện
       & {\caviar{\fontsize{14pt}{14pt}\selectfont Track used: The Autumn}}
      \\[6pt]
      % Dòng cuối cùng: Thời gian diễn ra của tập (thường sẽ là ngày phát hành)
       & {\continuum\fontsize{16pt}{16pt}\selectfont Ngày phát hành: 30/03/2022} \\[8pt]
    \end{tabular}
  \end{adjustwidth}
\end{table}
%===========================================NỘI DUNG==========================================
\color{black}

\begin{center}
  \uline{Một ngày tháng Bảy, năm 1946\\Ngôi đền Aogami.}
\end{center}

Cơn mưa tháng Bảy trút xuống ào ạt, dội lên mái ngói cong cong của ngôi đền ẩn mình trong rừng.

Tiếng mưa rơi hòa lẫn tiếng lá cây xào xạc, tạo thành một bức màn xám xịt bao phủ cả không gian. Ngôi đền nhỏ bé đứng lặng lẽ giữa màn mưa ấy, trầm mặc như một chứng nhân đã quen với sự tĩnh mịch của núi rừng. Những chiếc đèn lồng treo dưới hiên hắt ra ánh sáng mờ nhạt, lay động theo từng đợt gió.

Trên bậc thềm ướt lạnh, Misora Shino đứng nghiêm trang. Cô cúi đầu, hai tay chắp trước ngực, rồi đưa lên trán thành kính. Theo đúng nghi thức cầu phúc, cô kéo sợi dây thừng lớn treo ở gian chính, lắc mạnh để tiếng chuông ngân vang khắp không gian ẩm ướt. Sau đó, cô bỏ vài đồng xu vào hòm công đức, khẽ cúi chào hai lần, vỗ tay hai cái, rồi nhắm mắt cầu nguyện.

\img{e1-6b}

''Con\dots\ mong được kết thân với nhiều bạn bè hơn, và được sống quãng thời gian tuổi học trò với nhiều kỷ niệm đáng nhớ.''

Lời nguyện tưởng chừng đơn giản ấy nhưng lại nhỏ nhẹ đến lạ. Câu nói dần hòa tan trong tiếng mưa rả rích, chẳng biết liệu có chạm đến tai thần linh không, hay sẽ lại chỉ tan biến giữa không trung, như một lần nữa phụ lòng cô nàng.

Khi mở mắt ra, Shino khẽ thở dài. Một tay cô giương chiếc ô, tay kia xách cặp, chậm rãi bước xuống những bậc đá đã loang lổ nước mưa.

Không biết cô đã cầu nguyện bao nhiêu lần rồi. Chỉ biết rằng, số lần cô lên ngôi đền cầu phước cho bản thân nhiều đến độ mà cô còn chẳng muốn để tâm, và cũng là từng ấy lần chỉ đem lại cho cô một lần thất vọng.

''Giá như\dots\ kết bạn có thể dễ dàng như thế\dots'' cô lẩm bẩm. Khóe miệng nhếch lên một nụ cười gượng gạo, rồi lại tắt lịm. Tự trong lòng, Shino hiểu rằng muốn thay đổi, bản thân phải bước ra khỏi cái vỏ nhút nhát của mình. Nhưng cô chẳng có gì đặc biệt cả. Không tài năng, không nổi bật, thậm chí còn hay lúng túng mỗi khi nói chuyện. Một đứa con gái như thế, liệu ai sẽ muốn làm bạn?

''Mình có cầu nguyện nhiều đến đâu chăng nữa\dots\ cũng chẳng thể có bạn bè thật sự.''

Lời thì thầm ấy như hòa tan trong từng hồi thở dài nặng trĩu. Shino cảm thấy tủi thân, cảm thấy cô độc với chính mình. Cô tự hỏi liệu mình có đáng để ai đó dành thời gian làm quen, hay chỉ là một cái bóng mờ nhạt trong mắt mọi người. Những lúc như thế này, cô chỉ muốn thu mình lại, trốn vào cái vỏ an toàn mà cô đã dựng lên từ bấy lâu nay.

\img{e1-6c}

Tiếng mưa vẫn rơi, từng giọt đập vào mái hiên và nền đá, càng khiến bước chân Shino nặng nề hơn. Mỗi bước đi như mang theo cả gánh nặng của sự cô đơn và tự ti, những cảm xúc đã quá quen thuộc với cô gái trẻ đơn độc.

Nhưng rồi\dots

Khi cô đi gần tới bậc cuối cùng của con dốc đá, một bóng dáng bất chợt lướt qua tầm mắt. Shino ngừng bước, ngước nhìn lên bóng người xuất hiện ấy.

Đó là một cô gái có đôi mắt lạ lùng: một bên xanh biếc như đại dương, một bên đỏ rực như máu. Nước mưa chảy dài trên gương mặt vội vã của cô, càng khiến đôi mắt ấy trở nên nổi bật đến mức gần như siêu thực. Cô gái ấy có tóc đen dài đến ngang vai, được cột đều hai bên. Trên người cô gái là bộ đồng phục đen giống Shino: áo dài tay, váy xếp ly đồng màu, điểm thêm những sọc trắng ở cổ tay và cổ áo, cùng chiếc khăn quàng xanh thẫm. Trên cổ tay trái cô có đeo một chiếc vòng đỏ có đính ba chiếc chuông vàng nhỏ chiếu xuống.

\img{e1-6d}

Shino khựng lại. Trong thoáng chốc, cô không thể rời mắt. Vẻ đẹp lạnh lùng mà hối hả ấy khiến Shino cảm thấy vừa ngưỡng mộ, vừa chạnh lòng. ''\textit{Giá như mình có thể giống cô ấy\dots\ thì có lẽ mình đã đủ tự tin để mở lời với người khác rồi.}''

Cô đứng lặng hồi lâu, để mặc cho cơn mưa che giấu sự bối rối trong lòng. Trong thoáng chốc, cô thấy như mình đang đối diện với một tấm gương méo mó: một bên là dáng hình tự tin, sắc sảo và nổi bật đến mức chẳng thể không nhìn; một bên là bản thân mình, là một sự nhạt nhòa, lặng lẽ và chẳng có gì đáng để người khác dừng mắt lại.

Từ nhỏ, Shino vốn đã chẳng phải kiểu người thu hút ánh nhìn. Cô không học giỏi hơn ai, cũng chẳng có năng khiếu gì đặc biệt. Ở trường cũ, bạn bè xung quanh dần có những nhóm nhỏ để chơi cùng nhau, còn Shino thì thường ngồi im lặng ở góc lớp, lắng nghe mà ít khi chen vào câu chuyện. Không ai ghét bỏ cô, nhưng cũng chẳng ai thực sự nhớ đến cô. Và chính cái ''vô hình'' ấy, theo thời gian, đã trở thành một bức tường vô hình trong lòng cô.

Đôi khi, chỉ một lời trêu đùa vô tình cũng đủ để làm cô rút vào vỏ ốc của mình. ''Misora, cậu thật chả có gì nổi bật cả.'' Một câu nói bâng quơ của bạn cùng lớp ngày trước, đến giờ vẫn vang vọng trong ký ức Shino. Không hẳn là ác ý, nhưng nó như một con dao nhỏ, lặng lẽ khắc sâu thêm vào niềm tin rằng mình thật sự chỉ là một kẻ tầm thường.

Thậm chí, có những lúc Shino cố gắng chủ động bắt chuyện, mong muốn được kết bạn mãnh liệt đến mức khiến một số bạn học cảm thấy phiền hà. Sự nhiệt tình ấy đôi khi lại bị đáp lại bằng những lời châm chọc, giễu cợt, hoặc ánh mắt khó chịu. Có bạn còn nói thẳng: ''Cậu cứ bám riết thế này, thật phiền quá đi!'' Những câu nói ấy khiến Shino càng thêm thu mình lại, vừa xấu hổ vừa tủi thân, và nỗi cô đơn lại càng lớn dần lên trong lòng cô.

Lòng cô nặng trĩu. Trong đầu lại văng vẳng lời cầu nguyện ban nãy: ''Con mong được kết thân với nhiều bạn bè hơn.'' Chỉ là một mong ước giản đơn, vậy mà sao nó xa vời đến thế?

Shino nhìn xuống đôi giày lười ướt sũng vì mưa, rồi khẽ lắc đầu. ''\textit{Mình chẳng bao giờ dám chủ động bước tới gần người khác. Cứ mãi thế này, có lẽ mình sẽ chẳng bao giờ thoát khỏi cái bóng của chính mình.}''

Nhưng rồi, lý trí lại nhanh chóng kéo cô trở lại. Một sự quen thuộc lẫn bất an len lỏi. Trong ký ức của mình, Shino chưa từng thấy gương mặt ấy bao giờ. Ở trường, chắc chắn không có ai như vậy. ''\textit{Nếu thật sự chưa từng gặp\dots\ tại sao mình lại có cảm giác như đã nhìn thấy cô ấy ở đâu đó rồi?}''

\pov{???}

Ngày thứ ba liên tiếp trời đổ mưa.

Tôi kéo sát chiếc ô màu đen trên đầu, thở ra một hơi dài. Âm thanh rả rích hòa lẫn tiếng gió khiến con đường về nhà càng thêm ảm đạm. Nước mưa theo nhịp gió tạt vào ống tay áo, lạnh lẽo đến rùng mình.

Thật tình\dots\ thời tiết này làm tôi chẳng yên tâm chút nào. Đống quần áo tôi treo ngoài hiên chắc đã ướt sũng mất rồi. Mà nhà tôi\dots\ vốn chỉ có mình tôi xoay xở. Bố mẹ thì bận công việc ở nơi xa, thỉnh thoảng mới gửi thư về, nên việc chăm lo từng bữa cơm, từng bộ đồ, tất cả đều là một tay tôi. Nhiều khi, nghĩ đến cảnh chỉ có một mình trong căn nhà vắng, tôi lại tự nhắc nhở bản thân: phải trưởng thành hơn.

Ngày mai, tôi sẽ chính thức nhập học ở Trường Trung học Aogami. Thực ra hôm nay tôi đã đến trường để làm thủ tục, nhận lớp và bộ đồng phục. Mọi thứ diễn ra nhanh chóng, thầy cô chỉ điểm qua vài lời dặn dò, còn chúng tôi thì xếp hàng để nhận từng bộ đồ. Tôi nhớ rõ khoảnh khắc khi mặc thử nó trong phòng thay đồ cũ kỹ, ngắm mình trong gương: một chiếc áo dài tay màu đen, váy xếp ly đồng màu, khăn quàng xanh thẫm nơi cổ. Trông tôi không đến nỗi nào.

Nhưng tôi vẫn cảm thấy có gì đó xa lạ, như thể bộ đồng phục này vốn thuộc về ai khác chứ không phải tôi.

\dots Thôi thì, ít ra tôi đã có làm quen phần nào môi trường học tập ở nơi này. Tôi cũng chẳng thể nào nhớ nổi đây đã là lần thứ mấy trong cuộc đời học sinh tôi đã chuyển trường rồi. Trộm vía là tôi thích ứng với môi trường mới rất nhanh, nên tôi không quá bị ngợp khi phải thay đổi liên tục như vậy.

Tôi cũng đã gặp vài người bạn cùng lớp mới. Một vài cô gái tươi cười chào hỏi, vài cậu con trai tò mò liếc nhìn học sinh chuyển trường. Tôi đáp lại bằng những nụ cười nhẹ, vài lời chào đơn giản. Có lẽ, trước kia tôi sẽ chỉ im lặng, cúi đầu và rụt rè tránh đi. Nhưng bây giờ, tôi muốn thử mạnh dạn hơn một chút, muốn để lại một dấu ấn nào đó dù nhỏ nhoi. Tôi tự nhủ, sang ngày mai, mình sẽ không chỉ đơn thuần ''tồn tại'' trong lớp nữa.

Nghĩ đến đó, tôi siết chặt quai cặp, bước nhanh hơn trong cơn mưa nặng hạt. Thế rồi\dots

''\textbf{Aaaaahhh----!!}''

Một tiếng hét từ trên trời cao.

Tôi khựng lại. Cả cơ thể bất giác căng cứng như vì tiếng kêu bất ngờ. Chưa kịp ngẩng mặt lên để xem nó là cái gì, thì *RẦM!*, mặt đất ngay trước mắt tôi bắn tung bùn nước, một bóng đen nặng nề rơi xuống như thiên thạch. Tôi giật bắn, suýt tuột cả cặp lẫn ô khỏi tay.

Khói bụi và nước mưa hòa lẫn thành một lớp mờ đục. Tôi run rẩy lùi lại một bước. Khi tầm nhìn dần rõ, hai bóng người nằm sóng soài hiện ra trước mắt.

Một cô gái. Mái tóc xoăn dài được buộc đuôi ngựa bằng một dải nơ đỏ, đôi mắt xanh nước biển ánh lên giữa màn mưa. Cô mặc bộ đồng phục y hệt tôi đang mặc: áo đen, váy xếp ly, khăn quàng xanh thẫm.

Bên cạnh là một chàng trai. Mái tóc cũng xoăn, hơi rối, có gì đó giống hệt cô gái kia. Đôi mắt đỏ sẫm lóe lên thứ ánh sáng kỳ lạ. Trên người cậu là áo sơ mi dài tay đen, quần dài cùng màu, khoác thêm chiếc áo choàng dài khiến vóc dáng càng thêm nổi bật.

Tôi đứng sững lại. Trong lồng ngực, một cảm giác vừa quen vừa lạ xộc đến. Như thể tôi đã từng biết họ, từng ở cạnh họ\dots\ nhưng không sao nhớ ra được. Một nỗi nghẹn ngào mơ hồ dâng lên, làm tim tôi thắt lại.

''N-Này\dots\ Hai người ổn chứ?'' tôi cất giọng, hơi run rẩy vì bất ngờ, nhưng cố giữ bình tĩnh.

Cô gái chậm rãi ngồi dậy, tay giữ lấy đầu. Cậu con trai khẽ rên một tiếng, chống tay ngẩng lên nhìn quanh.

''Đau thật đấy\dots'' cậu con trai lầm bầm.

Cô gái thì nhăn mặt, khẽ lắc đầu: ''Chúng ta\dots\ rơi từ trên trời xuống thật à?''

Tôi không kiềm được, buột miệng thốt ra: ''Ừ thì\dots\ Hai người\dots\ đúng là từ trên trời rơi xuống. Mình\dots\ không nhìn nhầm đâu.''

Họ quay sang nhìn tôi. Trong khoảnh khắc ánh mắt giao nhau, cảm giác quen thuộc ấy lại siết lấy trái tim tôi mạnh hơn. Rõ ràng là tôi đã nhìn thấy họ ở đâu rồi, thậm chí tôi với hai người ấy đã từng chu du ở một nơi nào đó\dots

\dots Tức thật, tôi chẳng tài nào mà nhớ nổi.

Tôi vội vàng hít một hơi sâu, lắc đầu nhẹ để xua tan cơn hoang mang.

''Dù thế nào thì\dots\ đứng dậy đi, kẻo cảm lạnh đấy. Trời đang mưa to lắm.''

Tôi đưa tay ra. Cô gái ngập ngừng, rồi nắm lấy, dùng sức kéo dậy, kéo theo cả người bạn đồng hành - tôi cho là vậy - đứng lên. Chàng trai cũng gượng đứng lên nhờ vào ô còn sót lại chút che chắn của tôi.

''Cảm ơn\dots\ Tôi là Saikyou,'' cô gái giới thiệu, giọng vẫn còn run nhẹ.

''Tôi là Kyuen,'' chàng trai nối lời, đôi mắt đỏ vẫn không rời khỏi tôi.

Tôi nuốt khan một cái, rồi đáp với giọng dè chừng, dù có một chút nhẹ nhàng: ''Mình là Reiko. Nhà của mình cũng gần đây thôi. Nếu hai người không phiền\dots\ cứ về đó nghỉ tạm. Trông hai người thế này mà cứ đứng mưa, chắc chắn sẽ bệnh mất.''

Họ nhìn nhau, thoáng lưỡng lự. Nhưng rồi, Saikyou-san mỉm cười nhẹ: ''Nếu vậy thì\dots\ cảm ơn cậu, Reiko-san.''

''Nếu cậu không cảm thấy khó chịu thì\dots\ cho bọn tôi đi chung nhé?'' Kyuen-san hỏi, khép sát hơn vào tôi và cô gái kia. Tôi dám cá là hai người đó đi chung với nhau.

''Mình không ngại đâu. Ô của mình đủ lớn mà.''

Thế là, cả ba chúng tôi cùng chen chúc dưới một chiếc ô nhỏ, bước đi giữa màn mưa trắng xóa. Vai kề vai, thỉnh thoảng chạm nhau, nhưng không ai mở lời thêm một tiếng nào khác.

Chỉ có tiếng mưa rơi rả rích, và nhịp tim tôi cứ dồn dập như thể vừa đánh mất, vừa tìm lại một thứ gì đó vô cùng thân thuộc.

\ct

Mưa vẫn chưa dứt khi tôi đưa họ bước vào căn nhà nhỏ của mình. Cả ba chúng tôi đều ướt sũng, sàn gỗ kêu cót két dưới những vết chân loang nước. Tôi vội vàng đặt chiếc ô sang một bên, chạy vào trong lấy hai cái khăn sạch rồi đưa cho họ.

''Nhà mình đơn sơ lắm, hai cậu cứ tạm dùng vậy nhé,'' tôi gượng cười, tự dưng thấy ngượng vì nơi ở quá bình thường của mình.

''Không sao, tôi thấy như vậy cũng ổn mà,'' Kyuen-san nói, nhận lấy chiếc khăn rồi lau qua đầu của mình cùng bộ đồng phục sũng nước ấy.

''Bọn tôi cũng không ngại lắm đâu. Cảm ơn cậu nha, Reiko-san,'' Saikyou-san tiếp lời, lau kỹ càng hơn cậu chàng.

Tôi tất tả đun nhanh một ít nước nóng, rót ra chén cho họ cầm ấm tay. Căn phòng trở nên yên tĩnh, chỉ nghe tiếng mưa gõ nhịp ngoài mái ngói. Trong lòng tôi lúc đó đầy mâu thuẫn: một mặt thấy cảnh tượng ''hai người rơi từ trên trời xuống'' phi lý đến mức buồn cười, mặt khác lại không nỡ đẩy họ ra ngoài.

Nhìn thấy hai người họ đã ngồi ấm chỗ với cốc nước nóng trên tay, tôi cất tiếng hỏi, cố giữ giọng điềm tĩnh: ''Cho mình hỏi\dots\ vì sao hai cậu lại từ trên trời rơi xuống như vậy?''

Hai người đưa mắt nhìn nhau, rồi cùng đáp gần như đồng thanh: ''Bọn tôi\dots\ nói sao nhỉ, không thể nhớ nổi.''

Câu trả lời làm tôi thoáng nhíu mày. Rõ ràng không hợp lý, nhưng cách họ nói, giọng vừa bối rối vừa chân thành, lại khiến tôi khó ép mình nghi ngờ thêm. Tôi hít một hơi thật sâu, tự trấn an bản thân rằng: có lẽ họ thực sự bị mất trí nhớ sau cú ngã. Có lẽ họ sẽ nhớ lại sớm thôi.

\dots Mà kể cũng lạ. Bằng cách nào mà họ rơi xuống như thế được nhỉ? Liệu có phải họ như những người con từ trên trời được Thần linh cử xuống giống như những gì mà người dân ở đây truyền nhau không?

Thôi thì, tôi sẽ để việc đó sang một bên vậy.

Để phá vỡ sự im lặng, tôi nảy ra ý nghĩ khác:
''Vậy\dots\ mối quan hệ của hai người là thế nào vậy?''

Saikyou-san ôm lấy vai của Kyuen-san, mạnh dạn đáp: ''Tôi với anh ấy là hai anh em song sinh của nhau. Tôi là em gái, còn Nii-chan là\dots\ ừm, cậu cũng đoán được rồi đấy.''

Tôi liếc mắt nhìn qua nhìn lại hai người. Quả thật, hai người có nhiều điểm tương đồng thật: từ mái tóc xoăn, cho tới khuôn mặt nữa. Tôi cảm tưởng rằng\dots\ cô em gái lại là người hoạt bát hơn hẳn, hơn cả những gì tôi tưởng, còn người anh thì ngược lại, có phần trầm tính hơn.

''\dots Tên đầy đủ của hai người là gì?'' tôi bất chợt hỏi một câu khác.

Người con trai lên tiếng trước: ''Haragen Kyuen.''

Rồi đến lượt cô gái: ''Haragen Saikyou.''

Tôi khựng lại. Trái tim tôi lỡ một nhịp, tay vô thức siết chặt lấy tách trà. ''\textit{Haragen\dots?}'' Tôi mấp máy môi, cuối cùng cũng giới thiệu: ''Mình là\dots\ Haragen Reiko (\ruby{原}{はら}\ruby{阮}{げん}\ruby{鈴}{れい}\ruby{子}{こ}, \ruby{Ha-ra}{Nguyên}-\ruby{ghen}{Nguyễn}\ \ruby{Rê-i}{Linh}-\ruby{cô}{Tử}).''

''Cậu cũng có họ\dots\ `Nguyên Nguyễn' giống như bọn tôi sao?''

''\dots Ừ. Chính là Hán tự đó.''

Cả ba chúng tôi bỗng nhìn nhau, ánh mắt thoáng ngập ngừng. Tôi nghe rõ nhịp tim mình dội trong tai.  Cùng một họ, lại cùng một cách viết hiếm gặp. Chẳng lẽ chỉ là trùng hợp? Nhưng trong ánh mắt họ, tôi lại thấy một tia quen thuộc khó tả, giống như đã từng gắn bó ở đâu đó. Cảm giác ấy khiến tôi bất giác rùng mình.

Kyuen phá vỡ sự im lặng, hỏi khẽ: ''Cậu sống một mình ở đây sao?''

Tôi gật đầu, cố giấu nụ cười có phần méo xệch, dường như vì vẫn còn sốc: ''Ừ\dots\ Bố mẹ mình thường xuyên đi làm xa, có khi cả tháng mới gửi thư về. Nên gần như một mình phải tự lo hết. Riết cũng quen rồi.''

Tôi đưa tay trái vuốt lại mái tóc ướt sũng chưa kịp lau, và chợt nhận ra Saikyou-san đang chăm chú nhìn vào cổ tay tôi.

''Cái vòng này đẹp quá\dots\ ba chiếc chuông nhỏ, nghe vui tai thật,'' cô ấy khẽ nói.

Tôi ngập ngừng giây lát rồi chìa cổ tay ra cho họ thấy rõ hơn. Chiếc vòng tay đỏ, với ba chiếc chuông vàng khẽ rung lên theo nhịp cử động. Tôi đáp: ''À\dots\ cái này là quà của mẹ tôi. Bà bảo nó sẽ bảo vệ tôi khi ở một mình. Tôi cũng không biết có thật vậy không\dots\ nhưng nó là kỷ vật, nên tôi luôn mang theo.''

Nói tới đây, lòng tôi chợt se lại. Từ khi sống một mình, tôi đã học cách tỏ ra bình thản, nhưng mỗi khi nhắc tới món quà này, tôi lại cảm thấy mình nhỏ bé như đứa trẻ được mẹ che chở. Ánh mắt Kyuen-san và Saikyou-san nhìn chiếc vòng không đơn thuần là tò mò, mà còn phảng phất một nỗi bâng khuâng khó gọi thành tên.

Trong thoáng chốc, giữa tiếng mưa nặng hạt của một buổi chiều tháng Bảy, tôi lại thấy rõ rệt cái cảm giác quen thuộc ấy\dots\ như thể ba chữ Haragen vừa vô tình kéo chúng tôi lại gần hơn, trong một mối dây liên kết vô hình mà tôi chưa thể nào lý giải.

Tôi khẽ xoay nhẹ cổ tay, ba chiếc chuông leng keng như muốn xua đi cái im lặng nặng nề ấy. Nhưng thật tình, tôi cũng chằng biết phải tiếp tục cuộc hội thoại này như thế nào nữa.

Có rất nhiều điều vẫn còn đang lửng lơ trên đầu. Nhưng câu hỏi lớn nhất mà tôi không biết có nên thổ lộ không, đó là: họ - Kyuen-san và Saikyou-san - thực sự là ai? Và tại sao họ lại mang tới cho tôi một cái gì đó thân thuộc như thế?

Cái cảm giác lấp lửng này thực sự khiến tôi khó chịu. Nhưng tôi thực sự cũng không biết phải tìm ra lời giải như thế nào.

Đúng lúc đó, cậu chàng nhấp một ngụm trà, rồi bất ngờ cười nhẹ: ''Nói gì thì nói, được ngồi trong một căn nhà ấm áp thế này, giữa lúc trời mưa tầm tã ngoài kia, cũng là may mắn lắm rồi.''

Lời nói ấy khiến tôi bất giác thở phào. cpp em gái của cậu cũng gật đầu, thêm vào: ''Ừ, mưa kiểu này mà còn phải lang thang bên ngoài thì chắc ốm mất. Tôi lo cho Nii-chan của tôi lắm chứ.''

''Sao em cứ phải lo xa cho anh làm gì nhỉ\dots''

''Ủa? Chứ không phải lúc nào anh cũng dầm mưa về nhà để rồi bị cảm lạnh mấy ngày liền à?'' cô ấy đáp lại, hai má phụng phịu. ''Có cô em gái lo cho anh thế mà không biết điều!''

Tôi bật cười theo, bầu không khí vì thế cũng dịu đi. Không thể ngờ hai anh em nhà họ có thể tình cảm như thế đấy.

Chả bù cho tôi, giờ bỗng dưng tôi cảm thấy mình tủi thân. Là con một của một gia đình, tôi đôi lúc cũng muốn ước có thêm một người anh chị hoặc em để cho vui cửa vui nhà, làm cho cuộc sống xán lạn và vui vẻ hơn\dots\ Nhưng tôi chỉ có thể giữ điều đó trong lòng, không dám nói với bố mẹ.

\dots Mà ngẫm lại, tôi có cảm giác tôi và hai người họ đã từng ở trong một gia đình rồi\dots\ Nhất là từ Saikyou-san.

Kyuen-san bất chợt lên tiếng, ngoảnh mặt về phía tôi, rồi hỏi tôi một điều, như thể là một lời cầu xin: ''Ngày mai cậu nhập học ở Trường Trung học Aogami đúng không? Nếu được, bọn tôi cũng muốn học cùng lớp với cậu.''

Saikyou-san gật đầu, ánh mắt đầy hy vọng: ''Ừ, nếu trường cho phép, chúng tôi muốn nhập học cùng ngày với Reiko-san luôn. Có bạn đồng hành chắc sẽ dễ hòa nhập hơn.''

Tôi ngẫm nghĩ một lát rồi mỉm cười: ''Mình không thấy có vấn đề gì đâu. Trường cũng khá linh động với học sinh chuyển đến mà. Nếu hai cậu muốn, mình sẽ hỏi thầy cô giúp. Bộ ba chúng ta sẽ là những học sinh mới cùng lớp.''

Cả Kyuen-san và Saikyou-san đều tỏ vẻ nhẹ nhõm, không khí trong phòng trở nên gần gũi hơn. Tôi cảm nhận rõ sự đồng cảm giữa ba người, như thể chúng tôi vừa lập thành một nhóm nhỏ giữa thế giới rộng lớn này.

\ct

Từ đó, câu chuyện rẽ sang những điều bình thường hơn: về ngôi trường mà tôi sẽ nhập học ngày mai, về thời tiết dạo này, về những thói quen vụn vặt trong cuộc sống hằng ngày.

Tôi nhận ra, sự xuất hiện của hai vị khách xa lạ bất ngờ ấy đã khiến căn phòng vốn yên ắng và trầm lặng trở nên sống động hơn. Tiếng nói cười, những lời trò chuyện giản dị giữa ba người làm không khí ấm áp hẳn lên, xua tan đi cảm giác cô đơn thường trực trong lòng tôi. Đã lâu lắm rồi tôi mới có dịp được trò chuyện cùng ai đó như thế này, và bất giác, tôi thấy lòng mình nhẹ nhõm hơn đôi chút.

Nhưng mà\dots\ một cảm giác bất an mơ hồ bất chợt nảy lên trong lòng tôi. Tôi hướng mặt ra cửa sổ, nhìn những giọt mưa vẫn đang lộp bộp qua từng ô cửa sổ.

''\textit{Ba ngày nay mưa không dứt\dots\ Mình chưa từng thấy mưa dai dẳng thế này,}'' tôi lẩm bẩm, mắt liếc ra ngoài cửa sổ. Nước mưa đập vào mái ngói nghe như trống trận, từng vệt nước trườn dài trên khung kính. Trong lòng tôi thoáng rùng mình. Không hiểu sao, tôi có linh cảm rằng trận mưa này sẽ còn mang đến điều gì đó chẳng lành.

Kyuen-san chống cằm, đáp bằng giọng nửa đùa nửa thật, nhưng không phải là không chứa một sự lo lắng nhẹ: ''Ở đâu đó chắc chắn sẽ có đất đá bị cuốn trôi thôi. Mưa dày cộp quá mức rồi.''

Saikyou liếc cậu ta: ''Đừng hù người ta vậy, Nii-chan.''

Tôi khẽ cười, nhưng nụ cười tắt đi nhanh chóng. Có lẽ vì chính tôi cũng nghĩ vậy.

Chúng tôi tiếp tục chuyện phiếm thêm một lúc nữa. Cậu bạn hỏi tôi thích làm gì khi rảnh, còn cô em gái kể vài điều vụn vặt về bản thân. Tôi cảm thấy thật kỳ lạ: mới chỉ gặp vài giờ trước, nhưng ở cạnh họ, tôi lại chẳng thấy xa lạ. Như thể đã từng quen biết từ rất lâu, rất lâu rồi. Một cảm giác vừa dễ chịu, vừa khiến tôi bối rối.

Cuộc nói chuyện kéo dài đến trời xẩm tối, nơi mà bầu trời dần chuyển về ban đêm. Mấy ngày nay trời mưa rào nặng hạt, nên ngày tắt cũng sớm hơn so với thường lệ.

Saikyou-san bất chợt, mở lời: ''Reiko-san này\dots\ Bọn tôi có thể ở lại đây một đêm không? Chỉ một đêm thôi. Trời vẫn mưa, mà hai anh em tôi cũng không có chỗ nào khác để đi.''

Tôi thoáng khựng lại, cốc trà trong tay ấm lên như để che đi sự bối rối. Thông thường, tôi sẽ chẳng bao giờ đồng ý để hai người lạ mặt ở lại nhà mình. Nhưng với họ\dots\ tôi không hiểu vì sao, sự cảnh giác thường ngày bỗng mềm đi. Với lại, tôi cũng thấy họ không phải là người xấu, cũng không có dã tâm gì tâm độc với tôi cả, kể cả là từ Kyuen-san.

Tôi ngẩng lên, nhìn vào đôi mắt xanh biển của cô gái, rồi ánh đỏ thẳm của anh trai cô. Một thoáng xao động chạy dọc tim tôi. Cảm giác này\dots\ giống như gia đình. Một thứ gì đó mà tôi đã không trải qua từ một thời gian rất dài, tới độ mà có lẽ bắt đầu phai mờ trong tâm trí.

''Ừ,'' tôi đáp khẽ, mỉm cười để che đi sự ngượng ngùng. ''Cứ coi như đây là nhà của hai cậu, ít nhất cho tới khi mưa tạnh.''

Khi hai người đồng loạt cúi đầu cảm ơn, trong lòng tôi chợt dấy lên một nỗi bồi hồi. Như thể sợi dây vô hình mang tên Haragen đang quấn chặt cả ba chúng tôi lại với nhau. Tôi không thể lý giải, nhưng cũng không thể chối bỏ sự ấm áp ấy.

Ngoài kia, tiếng mưa vẫn rơi triền miên, dai dẳng\dots\ như một điềm báo cho điều gì đó sắp ập đến.

\pov{-}

Bữa tối hôm ấy trôi qua trong sự giản đơn và ấm cúng. Ba người cùng nhau chuẩn bị cơm canh, chẳng có món gì đặc biệt ngoài một ít rau, cá kho và súp miso, nhưng bầu không khí lại khiến mọi thứ trở nên ngon miệng lạ thường. Tiếng mưa rơi ngoài hiên hòa cùng tiếng cười trò chuyện vang trong căn bếp nhỏ, khiến không ai còn để ý đến sự nặng nề ban đầu.

Đêm xuống. Sau khi dọn dẹp xong, Reiko sớm chìm vào giấc ngủ trong căn phòng nhỏ của mình. Hơi thở cô đều đặn, khuôn mặt khi ngủ toát lên một vẻ bình yên hiếm có.

Kyuen và Saikyou ngồi bên chiếu trải tạm cho mình, chưa vội nằm xuống. Ánh đèn dầu hắt bóng cả hai trên vách, lung linh theo từng cơn gió nhẹ len qua khe cửa. Cả hai đưa mắt nhìn cô chủ nhà say ngủ, trong lòng cùng dâng lên một cảm xúc khó tả: một sự gần gũi như đã từng quen biết, và một cảm giác khó nắm bắt, cứ lẩn khuất mãi mà không thể lý giải. Như thể có một sợi dây vô hình nào đó đang ràng buộc ba người họ lại với nhau.

Saikyou khẽ lên tiếng, giọng thấp để không làm Reiko tỉnh giấc: ''\hl{Anh thấy không\dots{} Cô ấy thật giống như đã ở cạnh ta từ lâu rồi.}''

Kyuen mỉm cười nhạt, đôi mắt vẫn dõi theo bàn tay khẽ nắm lại của Reiko khi ngủ: ''\hl{Ừ. Anh cũng nghĩ vậy. Cảm giác này\dots{} anh từng trải qua rồi. Em cũng cảm nhận thấy nó chứ?}''

Cô em gái gật nhẹ. Một sự đồng điệu khó giải thích đang dần hiện rõ trong tim họ.

Một thoáng im lặng trôi qua, rồi Saikyou thì thầm: ''\hl{Còn Shino thì sao? Em không biết nên tin ả ta đến đâu\dots{} Nhưng những lời cô nói về hoa bỉ ngạn, về cái tên `Hikawa' kia\dots{} Em thấy nó không thể chỉ là ngẫu nhiên.}''

Cậu khẽ chau mày. Cậu nhớ lại nụ cười ngờ nghệch nhưng ánh mắt vô hồn của Shino, cùng những lời nói mơ hồ cô để lại.

''\hl{Anh cũng không chắc. Nhưng rõ ràng\dots{} chúng ta đã bị cuốn vào câu chuyện của cô ấy. Và có lẽ, câu trả lời sẽ nằm trong ngôi trường Aogami kia. Có lẽ\dots{} ta nên tin vào Shino một chút.}''

''\hl{\dots Em cũng thấy vậy. Dường như ả ta đang muốn tin vào câu chuyện mà ả kể. Dù gì thì ả cũng không có ý định gì xấu\dots}''

Saikyou kéo chiếc chăn mỏng phủ lên người, khẽ thở dài: ''\hl{Vậy thì ngày mai, chúng ta sẽ bắt đầu tìm hiểu. Dù có ra sao đi nữa.}''

Kyuen khẽ gật, ánh mắt kiên định, nhưng sâu trong lòng anh vẫn gợn lên một nỗi dày vò. Tiếng mưa ngoài hiên vẫn không ngừng rơi, dai dẳng như tiếng gõ nhịp vô hình báo hiệu cho những biến cố đang chờ phía trước.

\begin{center}
  \uline{Sáng hôm sau.}
  \pov{Haragen Saikyou}
\end{center}

Tiếng mưa vẫn rơi không dứt, gõ lộp độp lên mặt đất lầy lội. Tôi, Nii-chan, và Reiko-san cùng rời khỏi nhà sớm, vẫn chỉ với một chiếc ô duy nhất. Ba người chúng tôi chen chúc dưới mái che nhỏ bé ấy, vai kề vai, hơi ấm lan tỏa giữa màn mưa lạnh buốt.

Thị trấn Aogami hiện ra trong tầm mắt tôi hoàn toàn trái ngược với những năm đầu thế kỷ XXI, với những tòa nhà bê tông san sát, những quán cà phê ấm áp ánh đèn neon hay tiếng xe điện lướt qua nhịp nhàng.

Thay vào đó, nơi đây mang một dáng vẻ của một thị trấn vừa bước ra từ chiến tranh: những con đường đất còn lấm bùn, vài đoạn được lát bằng đá đã sứt mẻ; những ngôi nhà gỗ lợp mái ngói xám xịt, nhiều chỗ còn in dấu vết của những năm tháng khó khăn. Các cửa tiệm tạp hóa nhỏ nép mình trong góc phố, những tấm biển gỗ đã bạc màu, nhưng vẫn có người đứng che ô chào hỏi khách qua đường.

Nii-chan khẽ ngẩng nhìn bầu trời mù mịt, rồi thở dài. Tôi hiểu, trong lòng anh ấy hẳn đang so sánh khung cảnh này với nơi mà chúng tôi rời đi. Thời hiện đại đầy tiện nghi, ồn ã và ánh sáng rực rỡ, còn đây, là một Nhật Bản vừa thua trận, vẫn đang loay hoay tìm lại hơi thở của sự sống. Thế nhưng, giữa cơn mưa xám xịt, thị trấn này vẫn có nét thanh bình riêng: tiếng chuông gió khe khẽ nơi hiên nhà, mùi khói củi từ bếp lò len lỏi theo gió, và tiếng bước chân lộp bộp của những người dân trên đường đi chợ buổi sớm. Một sự thanh thản, trầm lắng - thứ mà ở những năm 2018 vẫn còn lưu giữ.

Đang mải nhìn ngắm, tôi bỗng nhận ra một bóng dáng quen thuộc. Cách đó không xa, một cô gái tóc nâu, đôi mắt xanh thẳm, với cài tóc ngôi sao lệch bên phải cũng đang cầm ô, bước trên con đường dẫn đến trường. Trái tim tôi khựng lại.

Không thể nào nhầm lẫn: Misora Shino. Chính ả ta, người đã dẫn chúng tôi vượt thời gian, cho thấy những mảnh ký ức đau thương ở hành lang kia.

Tôi thoáng nhìn sang Nii-chan. Anh cũng đã nhận ra. Nhưng có gì đó\dots\ khác lạ. Shino ở đây dường như không hề để ý đến chúng tôi, ánh mắt ả tập trung vào con đường phía trước, đau đáu một chút trầm buồn, không mang theo chút sát khí hay bí ẩn nào. Chỉ đơn thuần là một nữ sinh bình thường, như bao học sinh khác như chúng tôi.

''\dots Nii-chan,'' tôi khẽ thì thầm.

Anh đáp lại bằng ánh mắt, rồi cả hai cùng im lặng. Cảm giác khó nói xiết len vào trong tim, nhưng tôi biết, đây chưa phải lúc để vội vàng hành động. Reiko-san thì dường như không để tâm, vẫn tiếp tục đi về phía trước, miệng lẩm bẩm: ''Hôm nay trời lại mưa nhỉ\dots\ Đến bao giờ trời ráo đây\dots''

Cả nhóm tiếp tục bước đi, dần dần lướt qua Shino, hướng tới ngôi trường duy nhất của thị trấn. Càng tiến lại gần hơn, tôi càng thấy có nhiều học sinh khác cũng đang điềm tĩnh tới trường, với những chiếc ô dù mang màu sắc đơn điệu, mặc những bộ đồng phục như ba chúng tôi.

Ngôi trường Aogami hiện ra với dáng dấp hoài cổ, nổi bật giữa màn mưa xám. Toàn bộ công trình được xây bằng gỗ nâu, các thanh xà và cột trụ lộ rõ vết thời gian, khác hẳn với những ngôi trường hiện đại thường xây bằng bê tông và sơn trắng. Mái ngói xám phủ rêu, các cửa sổ gỗ mở rộng đón gió, tạo nên một không gian vừa mộc mạc vừa ấm áp. Thứ duy nhất khiến nơi này gợi nhớ đến trường học thời hiện đại là chiếc đồng hồ lớn dựng ở mặt tiền, kim giây vẫn đều đặn nhích từng nhịp, như một điểm giao thoa giữa quá khứ và hiện tại.

Dưới sự sắp xếp của nhà trường, cả ba chúng tôi được xếp vào cùng lớp 1-B. Không biết có phải vì không còn lớp trống, hay là một sự trùng hợp đầy ngẫu nhiên.

Lớp học ở ngôi trường những năm 1946 cũng khác lạ hơn so với thời hiện đại. Một lớp học rộng, với những dãy bàn ghế gỗ xếp thẳng hàng, trên bục còn có một cái bàn giáo viên đã sờn màu nâu của gỗ, và một tấm bảng xanh viết phấn. Hai bên lớp học là những ô cửa sổ, được mở ra để hứng lấy gió và tiếng mưa rơi ngoài hiên.

Rất giống với lớp học mà hai anh em chúng tôi đã nhìn thấy trong giấc mơ đêm hôm ấy. Nhưng khác ở chỗ, nó sáng sủa hơn và sạch sẽ gọn gàng hơn, vì lúc này mới là ban ngày.

Tôi và Nii-chan ngồi cạnh nhau ở cuối lớp, trong khi Reiko-san ngồi bên cạnh tôi, bên cửa kéo ra vào.

''Đây là ba học sinh mới sẽ nhập học từ hôm nay: Haragen Reiko-san, Haragen Kyuen-san, và Haragen Saikyou-san. Mọi người hãy cùng nhau giúp đỡ các bạn nhé.''

Giọng giáo viên vừa dứt, cả lớp lập tức xôn xao. Tôi và Nii-chan còn chưa kịp ấm chỗ trên hàng ghế gỗ thì mấy nữ sinh đã ào tới, vây kín lấy ba đứa chúng tôi.

Một cô gái với mái tóc xanh cột đuôi ngựa gọn gàng, gương mặt sáng sủa nhưng mang nét kiêu kỳ, ánh mắt xanh lam tỏa sáng ngôi sao liền bước thẳng lại. ''Haragen-san, đúng không? Lạ ghê, ba người cùng họ à?'' cô hỏi, giọng có chút sắc bén.

Tôi khẽ gật đầu, còn Reiko-san lễ phép đáp: ''Ừ, chỉ là trùng hợp thôi.''

Cô gái đó hơi nheo mắt. ''Tôi là Furukawa Nanase. Tôi ghét sự giả dối, nên hy vọng mấy cậu đừng có giấu chuyện gì kì lạ.''

Tôi cau mày, nhưng Nii-chan chỉ mỉm cười lịch sự: ''Tôi cũng chẳng ưa gì bọn nói dối đâu, nên cậu cứ yên tâm là tôi sẽ không có gì để giấu cả.''

''Vậy thì tốt. Mong chúng ta có thể giúp đỡ nhau,'' Nanase-san chìa tay ra và lịch sự bắt tay với anh trai tôi, như thể cô ấy đang đi gặp một đối tác lớn.

\chrint

Furukawa Nanase(\ruby{古}{ふる}\ruby{河}{かわ}\ruby{七}{なな}\ruby{星}{せ}, \ruby{Phư-rư}{Cổ}-\ruby{ca-oa}{Hà}\ \ruby{Na-na}{Thất}-\ruby{sê}{Tinh}) là một người rất đặc biệt trong mắt tôi. Cô ấy là con gái của một doanh nhân giàu có, và chính công ty của bố cô đã góp phần xây dựng ngôi trường trung học này. Có lẽ vì xuất thân như vậy, cô luôn mang theo mình một khí chất tự tin, thậm chí có phần kiêu kỳ. Nhưng tôi biết, điều đó không phải là sự tự cao rỗng tuếch, mà là kết quả của việc cô ấy luôn phải chứng minh bản thân giữa những ánh nhìn soi mói về địa vị gia đình.

\begin{wrapfigure}[14]{r}{5cm}
  \includegraphics[height=8cm]{../../../../Image/Book?/nanase.png}
\end{wrapfigure}

Nanase-san cực kỳ chính trực và trung thực, ghét sự bất công và gian dối. Cô ấy không ngại nói thẳng, đôi khi khiến người khác cảm thấy khó chịu vì sự bộc trực ấy. Nhưng tôi hiểu, đó là cách cô bảo vệ những giá trị mà cô tin tưởng, nhất là khi xung quanh luôn có những lời đồn đoán hoặc sự nghi ngờ về động cơ của cô. Đối với tôi, cô ấy là kiểu người mà nếu đã tin tưởng thì sẽ không bao giờ phản bội, và nếu đã ghét sự giả dối thì sẽ không ngần ngại vạch trần nó trước mặt mọi người.

\chrint

Một nữ sinh khác có mái tóc xám với phần tóc tết ở đỉnh đầu và cơ thể rắn rỏi liền chen ngang. ''Thôi nào, Nanase, đừng có khô khốc chứ. Tên mình là Shiromi Yae. Nói thật nhé, nhìn cậu kia,''  cô chỉ thẳng vào tôi ''trông cũng khỏe phết. Có chơi thể thao không? Hay chỉ giỏi cầm sách vở?''

Tôi cười nhạt. ''Tôi thì môn nào cũng chơi được. Muốn thử không?''

Cô nàng bật cười sảng khoái. ''Khá đấy, tôi thích người thẳng thắn!''

\chrint

\begin{wrapfigure}[14]{l}{7.5cm}
  \includegraphics[height=8cm]{../../../../Image/Book?/yae.png}
\end{wrapfigure}

Shiromi Yae (\ruby{銀}{しろ}\ruby{鏡}{み}\ruby{八}{や}\ruby{重}{え}, \ruby{Si-rô}{Bạc}-\ruby{mi}{Cảnh}\ \ruby{Gia}{Bát}-\ruby{ê}{Trùng}) là một cô gái cực kỳ năng động, khả năng thể thao thì khỏi phải bàn. Xem chừng đối thủ tiềm năng của tôi đã có rồi đây.

Tôi để ý thấy cô ấy luôn nói về lý tưởng của mình về việc một cô gái nên sống ra sao, ''chuẩn''  truyền thống như thế nào, và tôi cảm thấy cô nàng đang theo đúng những gì mà cô theo đuổi. Đến mức sự nhẹ nhàng và hiền hậu của cô có phần hơi ngược với khả năng thể thao xuất chúng của mình.

Yae-san là bạn thân thuở nhỏ của Nanase-san, và hai người họ gắn bó với nhau như hình với bóng. Tôi cảm nhận được cô ấy luôn quan tâm đến cô bạn của mình, dù ngoài mặt thì Nanase-san lúc nào cũng tỏ ra cứng đầu. Chính cô là người duy nhất hiểu được những cảm xúc tinh tế của Nanase-san, có thể nhìn thấu được những mặt khác mà không phải ai cũng để ý.

Có lẽ chính tình bạn ấy đã giúp cô bạn của mình giữ được sự chân thành và thẳng thắn, dù cho môi trường xung quanh có phức tạp đến đâu.

\chrint

Ngay sau đó, một cô gái tóc vàng rực rỡ, nụ cười lanh lợi, chồm tới, ánh mắt lấp lánh. ''Ồ ồ! Học sinh mới, lại ba người liền cơ à! Quá ngầu luôn! Này, Reiko-chan, Saikyou-chan, Kyuen-kun, ba người có phải họ hàng thật không? Hay là định lập ban nhạc gia đình hả?''

Cô nàng có mái tóc tết buộc lệch bên phải, với một chỏm tóc hồng đính ở giữa đỉnh đầu trông như một cái chong chóng, và có một cặp\dots\ tai cáo? Tai mèo? Tôi cũng chẳng rõ nữa. Mà, gọi chúng tôi thẳng bằng tên một cách tự nhiên như thế ngay từ lần gặp đầu tiên sao?

Reiko-san hơi đỏ mặt. ''Không\dots\ không phải như vậy đâu. Bọn mình cũng chỉ mới gặp nhau thôi.'''

Cô gái đó cười khanh khách. ''Tôi là Omori Honoka! Rất vui được gặp nhé! Nếu có tin tức gì thú vị, nhớ kể tôi nghe!''

\chrint

\begin{wrapfigure}[14]{r}{7.5cm}
  \includegraphics[height=8cm]{../../../../Image/Book?/honoka_human.png}
\end{wrapfigure}

Omori Honoka (\ruby{尾}{お}\ruby{森}{もり}\ruby{歩}{ほの}\ruby{香}{か}, \ruby{Ô}{Vĩ}-\ruby{mô-ri}{Sâm}\ \ruby{Hô-nô}{Bộ}-\ruby{ca}{Hương}) là kiểu người mà tôi chỉ cần nhìn thoáng qua đã biết sẽ chẳng bao giờ chịu ngồi yên một chỗ. Từ chuyện trên trời dưới biển, từ chuyện nhỏ tới chuyện lớn, cô ấy là kiểu người mà chỉ cần một ai đó mở lời là lập tức sẽ sán vào ngay lập tức.

Honoka-san là một cô gái hoạt bát, nhí nhảnh, luôn cố gắng khuấy động bầu không khí trong lớp. Cô ấy lúc nào cũng lạc quan, chẳng bao giờ để tâm đến những chuyện buồn quá lâu. Thậm chí đôi lúc có phần lạc quan thái quá. Y như cô bạn Suzune-san hay Hanabi-san mà tôi đã gặp.

Cô thích buôn chuyện, cực kỳ tò mò, và luôn sẵn sàng nhúng tay vào những chuyện rắc rối - ít nhất đó là lời kể của các bạn học trong lớp. Ở cạnh Honoka-san, tôi cảm nhận được một luồng năng lượng tích cực toát ra từ cô ấy, khiến cho tôi cảm thấy an tâm phần nào.

\chrint

''À mà, cho mình hỏi chút, Omori-san\dots'' Nii-chan liền cất lời, đôi mắt nhìn lên cặp tai thú trên đỉnh đầu. ''Cái tai thú trên đỉnh kia\dots\ là thật hả?''

Honoka-san chỉ vào hai cái ''tai'' của mình, nghiêng đầu: ''À, cặp tai cáo sa mạc ấy á? Tai giả thôi! Đây này!''

Để chứng minh, cô nhấc cái cặp ''tai'' của mình lên, và quả như tôi nghĩ, chúng chỉ là đồ để trang trí. Cô nàng là một con người bình thường, giống như chúng tôi. ''He he, đáng yêu không\~{}?''

''\dots Ừ thì, đúng là đáng yêu thật,'' Reiko-san cười ngượng.

''À, mọi người có thể gọi tôi là Honoka được rồi. Không cần khệ nệ gì nha!'' cô bạn tóc vàng nháy mắt, nở một nụ cười nhe răng đầy vui vẻ.

Bên cạnh, một cô gái tóc hồng trong bộ đồng phục chỉnh tề bước tới. Cô có chỏm tóc ngố ở đỉnh và để xõa, với một cột tóc ngắn buộc lệch bên trái được đeo buộc tóc với cái chuông. Vẻ mặt cô ấy trông lơ đãng, đôi mắt như đang ấp ủ điều gì.

''Tôi là Shinomiya Sakura. Nhà tôi làm ở đền. Ba người mới tới\dots\ hừm, quanh người có chút khí tức khác thường đấy.'''

''Khí tức\dots?'' tôi nhướng mày.

Cô nghiêng đầu. ''Ừ, tôi cảm nhận được. Có gì đó hơi giống như\dots\ dấu vết của một lời nguyền. Tôi sẽ phải tìm hiểu thôi.''

Nii-chan cười gượng. ''Chúng tôi chỉ là học sinh bình thường thôi mà. Có phải là ma quỷ đâu mà lời nguyền gì chứ\dots''

Sakura mím môi, ánh mắt lấp lánh sự hứng thú.

\chrint

\begin{wrapfigure}[14]{l}{5cm}
  \includegraphics[height=8cm]{../../../../Image/Book?/sakura.png}
\end{wrapfigure}

Shinomiya Sakura (\ruby{四}{しの}\ruby{宮}{みや}\ruby{桜}{さくら}, \ruby{Si-nô}{Tứ}-\ruby{ni-gia}{Cung}\ \ruby{Xa-cư-ra}{Anh}) là một người khác biệt hẳn với những bạn nữ khác trong lớp. Cô ấy xuất thân từ một gia đình gắn bó với ngôi đền ở thị trấn Aogami từ thời xa xưa - chính cô ấy vừa nói với tôi như vậy với giọng đầy tự hào.

Sakura-san được mọi người trong lớp nói là có kiểu tính cách rất hay tranh luận, đôi khi nói ra những điều chẳng ai hỏi tới, cứ như một nhà bác học nhỏ tuổi luôn tò mò về mọi thứ xung quanh - hay như cách mà tôi hay đùa là ''bà cụ non.'' Dù không cố ý gây phiền hà, nhưng cô ấy thường buột miệng nói những điều không cần thiết, khiến không khí trong lớp đôi lúc trở nên kỳ lạ, đôi lúc còn khiến mọi người khó xử.

Điều khiến tôi chú ý nhất là niềm hứng thú của Sakura-san với những câu chuyện huyền bí. Cô ấy luôn hỏi han về các truyền thuyết, những hiện tượng lạ, và đặc biệt bị ám ảnh bởi việc thu thập thông tin. Có lẽ vì thế mà Sakura thường xuyên trốn học, chỉ để đi tìm hiểu những bí mật quanh thị trấn hoặc ngôi đền của gia đình mình.

Tôi cảm tưởng, ở cạnh Sakura, mình sẽ không bao giờ thiếu chuyện để nghe, dù đôi khi phải tự hỏi liệu những điều cô ấy kể có thật hay không. Và tôi nơm nớp nghĩ, có lẽ \emph{cô ấy sẽ đóng vai trò nào đó} về những hiện tượng kỳ quặc mà hai anh em tôi đã phải trải qua.

\chrint

Cuối cùng, từ cuối lớp, một tiếng ngáp vang lên. Một nữ sinh tóc bạc cột hai bên, tay đút túi váy, bước chậm rãi lại gần.

''Ồn ào quá đấy.'' - giọng cô ta trầm, có chút ngái ngủ nhưng đầy sức nặng. ''Fukami Shiina. Học sinh mới à?''

Cả ba chúng tôi gật đầu. ''Ừm. Bọn mình là những học sinh mới chuyển tới.''

Cô nhìn tôi từ trên xuống dưới, rồi nhìn sang Nii-chan, cuối cùng dừng lại ở Reiko-san. Cô nhíu mày, cất giọng hơi khàn như một gã đồ tể: ''Được đấy. Ít ra trông không yếu ớt. Nếu bọn côn đồ ngoài phố dám động tới, cứ gọi tôi.''

Tôi bất giác bật cười. ''Chắc tôi sẽ ăn nhập với cậu nhất đấy, Fukami-san. Tôi cũng là một kẻ chuyên trừng trị bọn bắt nạt Nii-chan mà.''

Cô nhếch môi, một nụ cười khẽ, rồi im lặng quan sát như thể chẳng có gì quan trọng.

\chrint

\begin{wrapfigure}[14]{r}{8cm}
  \includegraphics[height=8cm]{../../../../Image/Book?/shiina_human.png}
\end{wrapfigure}

Fukami Shiina (\ruby{深}{ふか}\ruby{見}{み}\ruby{獅}{し}\ruby{那}{いな}, \ruby{Phư-ca}{Thâm}-\ruby{mi}{Kiến}-\ruby{Si}{Sư}-\ruby{na}{Na}) là kiểu người mà tôi có thể đoán được trong cái chớp mắt: cô là một học sinh cá biệt, mạnh nhất thị trấn Aogami, hay dùng vũ lực để giải quyết vấn đề.

Cô ấy tự tin vào khả năng chiến đấu của mình, liên tục gây gổ với đám côn đồ từ các trường khác và đều thắng, chẳng bao giờ chịu thua ai. Tôi cũng là người hay động chân động tay, nhưng chỉ khi cần bảo vệ Nii-chan hiền hòa của mình, còn Shiina thì sẵn sàng ra tay bất cứ lúc nào nếu thấy ai đó dám thách thức cô.

Dù nóng tính, Shiina-san lại có một mặt điềm tĩnh, luôn quan sát xung quanh trước khi hành động. Sự điềm tĩnh đó chính là minh chứng cho sức mạnh thật sự của cô ấy. Tôi cảm thấy vừa nể vừa có chút cảnh giác, bởi cô ấy không chỉ mạnh mẽ mà còn rất khó đoán.

Tiện thể thì, đôi tai trên đỉnh đầu của cô ấy cũng là tai giả. Honoka-san nằng nặc muốn cô đeo chúng, và cô cũng chỉ có thể nghe lời. Ban đầu cô còn cảm thấy khó chịu vì thấy vướng, nhưng đeo nhiều rồi cũng quen.

\chrint

Không khí giữa chúng tôi dần náo nhiệt. Cả năm cô gái bạn học của lớp, mỗi người một vẻ: Nanase-san vẫn giữ vẻ lạnh lùng, thỉnh thoảng xen vào câu hỏi trực diện; Yae-san thì cười nói ầm ĩ, hứng khởi muốn so tài thể thao; Honoka-san buôn chuyện không ngừng; Sakura-san lẩm bẩm về điềm lạ; còn Shiina-san thì khoanh tay dựa bàn, thỉnh thoảng liếc sang, như một con sói im lặng.

Reiko-san có vẻ hơi lúng túng trước sự chú ý bất ngờ, nhưng Nii-chan thì khéo léo đáp lại từng câu hỏi, giữ cho bầu không khí nhẹ nhàng. Tôi ngồi cạnh, nghe mà vừa mệt vừa buồn cười.

Đúng lúc cuộc trò chuyện đang rôm rả, cánh cửa lớp bật mở. Một bóng dáng bước vào, chiếc ô còn nhỏ nước. Mọi người không hẹn mà cùng ngẩng lên.

Misora Shino.

Nhưng kỳ lạ thay, không một ai trong lớp để tâm. Tiếng cười nói vẫn tiếp tục, và chỉ có tôi, Nii-chan, cùng cô bạn mắt hai màu khựng lại khi thấy cô bước vào.


Tôi khẽ nhíu mày. ''\textit{Chẳng lẽ Shino là kiểu người nhạt nhòa đến mức chẳng ai buồn nhìn? Vậy mà\dots}'' trong đầu tôi còn đang cố ráp lại mảnh ghép, so sánh cô với Shino mà chúng tôi từng gặp ở hành lang lúc đó.

Ngay lúc ấy, một giọng nói tươi rói cắt ngang mạch suy nghĩ.

Honoka-san gần như dí sát mặt tôi, đôi mắt sáng long lanh: ''Saikyou-chan, nghe này, nghe này!''

Tôi thở dài. ''Gì thế, Honoka-san?''

Cô cười rạng rỡ, vẻ mặt háo hức: ''Kể chuyện gì hay hay đi, cho mọi người cùng nghe với! Ba người tới đây bằng cách nào? Có phải vừa chuyển từ Tokyo không? Hay từ Hokkaido?!''

Thế là tôi lại bị cuốn trở lại vòng xoáy trò chuyện, bỏ mặc bóng dáng Shino đang lặng lẽ ngồi xuống cuối lớp.

Honoka-san tiến sát tới chúng tôi, cất giọng hào hứng tới lạ: ''Này này, thật sự ba người có họ giống nhau mà không phải anh em ruột hả? Kỳ lạ ghê nha!''

Reiko-san rụt rè đáp, giọng nhỏ nhẹ: ''Ừm\dots\ chỉ là trùng hợp thôi, không có quan hệ họ hàng. Kyuen-san với Saikyou-san là hai anh em ruột thịt.''

Nanase-san khoanh tay trước ngực, nhíu mày: ''Hừ, nghe hơi khó tin đấy. Dám chắc là các cậu không định lừa phỉnh bọn này để tạo bất ngờ đấy chứ? Ý tôi là, ba người có cùng cách đọc và cách viết Hán tự như thế mà.''

Anh tôi mỉm cười nhã nhặn, giọng điềm đạm: ''Chúng tôi cũng không muốn khiến mọi người hiểu lầm đâu. Reiko-san chỉ tình cờ cùng họ với chúng tôi thôi.''

Yae-san hào hứng chen ngang, mắt sáng lên: ''Thế còn thể thao thì sao? Saikyou-chan này, nhìn cơ bắp cậu chắc khỏe lắm. Chạy 100m được bao nhiêu giây?''

Tôi cười nhạt: ''Chưa từng bấm đồng hồ, nhưng đánh nhau thì nhanh hơn cả chạy.''

Shiina-san tựa cằm lên tay, liếc nhìn tôi với ánh mắt thích thú: ''Hừm. Nghe giống tôi đấy. Có máu chiến thì giữ vững đi, đừng làm kẻ yếu.''

''Tôi chẳng định yếu đâu,'' tôi đáp gọn lỏn. ''Nếu yếu thì ai bảo vệ cho Nii-chan của tôi, đúng không?''

''Hai người chỉ nghĩ tới mấy chuyện đó thôi à\dots'' cô bạn tóc bạc thở dài.

Honoka-san vỗ tay cái đét, reo lên: ''Ôi trời, hai người này hợp cạ rồi nha!''

Sakura-san ngả đầu, thì thầm với giọng bí ẩn: ''Haragen\dots\ cả ba người lại có cùng họ một cách trùng hợp thế kia\dots\ thú vị thật.''

''Cậu nghĩ vậy thật sao?'' Reiko-san hoảng hốt.

Cô nàng tóc hồng mỉm cười khó hiểu: ''Ừ. Biết đâu mấy cậu mang theo điềm lạ nào đó.''

Nanase-san nhăn mặt, lắc đầu: ''Lại mê tín nữa rồi, Sakura. Đừng dọa người khác.''

''Nhưng mà nghe cũng hồi hộp ghê!'' Honoka-san cười khúc khích. ''Có khi Saikyou-chan và Kyuen-kun là anh hùng bí mật đấy!''

Nii-chan liền lắc đầu nhẹ: ''Anh hùng thì không đâu. Chúng tôi chỉ là học sinh bình thường thôi.''

Yae-san giơ nắm tay lên đầy phấn khích: ''Bình thường mà mạnh thì càng tốt! Sau này nhớ đọ sức tay đôi với tôi nhé.''

''Sẵn sàng thôi,'' tôi đáp.

Honoka-san quay sang Reiko-san: ''Reiko-chan, còn cậu thì sao? Có năng khiếu gì đặc biệt không? Đừng nói là biết nấu ăn nhé, tôi sẽ mừng lắm luôn!''

''\dots Ừm\dots\ mình nấu ăn được một chút,'' Reiko-san ngập ngừng đáp. ''Mình không thực sự tự tin đến thế.''

Honoka-san mắt sáng rực lên: ''Thế là được rồi! Lần tới nhất định phải mời cả lớp ăn cơm hộp nha!''

''Cơm hộp thì tôi ăn,'' Shiina-san cắt lời, giọng lười nhác. ''Nhưng đừng có làm nhạt nhẽo quá.''

''Đúng là cậu chỉ biết ăn với đấm thôi,'' cô bạn tóc xanh lắc đầu.

Cô nàng bờm tai sư tử nhún vai: ''Có sao đâu.''

Cuộc nói chuyện càng lúc càng ồn ào, giống như một cái chợ nhỏ ngay trong lớp học. Tôi ngồi giữa đám con gái ấy, vừa thấy ngợp, vừa thấy thú vị.

\pov{Misora Shino}

Tôi bước vào lớp, như mọi ngày, lặng lẽ và không một tiếng động. Tôi đặt chiếc ô vẫn còn đang sũng nước để bên ngoài cửa ra vào, rồi nhanh chóng thay giày trong nhà, nhanh bước chạy về lớp 1-B.

Cửa vừa trượt mở ra, không khí náo nhiệt bên trong ùa vào, ồn ào hơn hẳn thường lệ. Nhưng cái sự ồn ào đó lại chẳng liên quan gì đến tôi.

Chẳng một ánh mắt nào hướng về phía tôi. Tất cả đều đang đổ dồn vào ba gương mặt mới mẻ ở dãy bàn gần cửa sổ. Tôi dừng chân trong thoáng chốc, rồi ánh mắt bất giác dán chặt vào một cô gái có mái tóc đen dài buộc hai bên, đôi mắt mang hai màu xanh và đỏ lạ lẫm.

\dots Đó chẳng phải là người tôi đã thoáng gặp hôm qua sao? Hóa ra là học sinh mới\dots\ không lạ gì khi cô ấy trở thành tâm điểm như vậy. Cũng dễ hiểu mà vì sao tôi không nhận ra.

\img{e1-6e}

Một thoáng nhói lên nơi ngực, tôi nhận ra mình đang chạnh lòng. Cái vị trí kia - nơi mà mọi ánh nhìn đều tập trung vào - chưa bao giờ thuộc về tôi.

Rồi tôi nhận ra còn hai người khác, một nam một nữ, cũng đang bị các bạn trong lớp vây lấy. Cả hai đều mang họ Haragen, ít nhất là những gì mà tôi nghe lỏm được từ các bạn trong lớp. Tôi không biết họ là ai, nhưng nghe qua vài câu hỏi han, hóa ra họ là anh em ruột.

Hai người ấy hướng mắt nhìn về phía tôi, một cách thoáng qua, có lẽ ngay từ lúc tiếng cửa mở. Đôi mắt ấy, với một cặp mắt xanh và một cặp mắt đỏ, giống như màu mắt của cô gái kia, có gì đó lạ lùng, như thể đang dò xét, nghi hoặc. Tôi thoáng rùng mình, không hiểu nổi tại sao lại như vậy. Trông tôi giống như một kẻ ngoại lai đến vậy sao?

Nhưng ngay sau đó, một bạn gái tóc vàng nổi bật, mà tôi nhớ hình như tên là Omori-san, đã kéo họ trở lại vào vòng trò chuyện, khiến ánh nhìn kia cũng biến mất. Những câu hỏi và câu đáp cứ thế trôi tuồn tuột đi, đưa tôi quay trở về thực tại đầy tàn nhẫn và phũ phàng trước mắt.

Không muốn dừng lại thêm, tôi lủi thủi bước nhanh về phía cuối lớp, ngồi xuống chỗ quen thuộc của mình, tréo ngoe thế nào lại ngồi ngay cạnh cô bạn có mắt hai màu kia. Từ chỗ này, tôi có thể thấy rõ cảnh tượng kia: ba học sinh mới ngập trong vòng vây bạn bè, tiếng cười nói rộn ràng, hoàn toàn trái ngược với tôi.

\img{e1-6f}

Tôi chống cằm, mắt dõi theo họ, trong lòng pha trộn giữa chán nản và một chút ghen tị. Cả ba con người ấy khi nhập học, lập tức có thể khiến cả lớp xôn xao, mọi sự chú ý đều đổ dồn vào họ. Còn tôi thì\dots\ mãi chỉ là một cái bóng mờ nhạt, mãi không thể bắt chuyện được với bất cứ ai.

Tôi cố tự an ủi bản thân: ''\textit{Dù có như thế nào, thì nó cũng chẳng liên quan gì đến cuộc đời mình cả.}'' Tôi sẽ tiếp tục ngày này qua ngày khác, chui rúc trong cái kén an toàn của riêng mình. Nhưng cái vị đắng của nỗi buồn bị bỏ rơi vẫn len vào, không tài nào giấu được.

Có lẽ tất cả đều là lỗi của tôi. Sự rụt rè và thiếu tự tin đã trở thành xiềng xích trói buộc tôi trong suốt bao năm qua. Mỗi khi có cơ hội kết bạn, tôi lại tự thu mình vào vỏ ốc, để rồi nhìn cơ hội ấy tuột khỏi tay. Năm tháng cứ thế trôi qua, và tôi vẫn không thể có được một người bạn thật sự. Nhìn những học sinh mới chuyển đến, tự tin và rạng rỡ kia, tôi càng thêm nhận ra khoảng cách giữa họ và bản thân mình xa vời đến nhường nào.

Cả đời tôi\dots\ dường như đã định đoạt như vậy. Cô độc. Một mình. Không có một ai có thể chia sẻ cùng. Và tất cả, đều bắt nguồn từ chính sự yếu đuối trong tôi.

\dots

\dots

''Phải bắt đầu bằng việc tự thay đổi bản thân\dots\ buộc mình phải bước ra khỏi cái vỏ nhút nhát của mình.''

Tôi thầm nhớ lại câu tự nhủ của chính mình hôm qua khi rời ngôi đền. Tôi đã lặp đi lặp lại câu đó nhiều lần, như một lá bùa hộ mệnh. Nhưng hôm nay, khi nhìn thấy sự náo nhiệt ở dãy bàn phía trước, lòng tôi chùng xuống. Tôi nhớ lại những lần mình không thể mở lời chào, những khoảnh khắc mà ai cũng đã có nhóm bạn riêng, riêng tôi thì lặng lẽ đứng ngoài cuộc. Từng kỷ niệm như một nhát dao cứa sâu, khiến tôi bứt rứt.

Tôi không thể để nó tiếp diễn nữa.

Tôi đứng bật dậy. Trái tim đập dồn dập đến mức tai ù đi. Cả lớp thoáng sững sờ. Chắc hẳn mọi người đang bất ngờ, bởi ''con bé Shino lặng thinh'' lại tự dưng lao đến thế này.

Nhưng tôi không quan tâm nữa. Nếu cơ hội này tôi không nắm bắt được, liệu tôi sẽ còn có những cơ hội nào khác nữa?

Tôi bước tới gần cô gái mới chuyển trường - mái tóc đen dài cột hai bên, đôi mắt xanh và đỏ như chứa một thứ ánh sáng kỳ lạ. Tôi hít sâu, nhưng tay vẫn run rẩy. Tôi đưa tay ra, chạm lấy bàn tay cô ấy, giọng lắp bắp:

''Cậu\dots\ ờm\dots\ cậu có muốn\dots\ làm bạn với mình\dots\ không?''

\img{e1-6g}

Tôi nghe rõ cả tiếng tim mình đập khi thốt ra câu ấy. Một phần tôi muốn buông tay chạy về chỗ vì xấu hổ, nhưng một phần khác lại muốn bám chặt để không bị nuốt chửng bởi nỗi sợ.

Đôi mắt tôi dường như ánh lên chút quyết tâm. Nhưng tôi biết rõ, đó là một sự quyết tâm mong manh, chực chờ bị một cơn gió nhẹ cuốn đi.

Và rồi, cô gái ấy khẽ nở một nụ cười hiền hâu và đáp lại: ''Được chứ. Rất sẵn lòng.''

Khoảnh khắc đó, tôi như vừa bước sang một trang sách mới. Tôi không còn hoàn toàn đơn độc nữa. Từ đó, tôi dần dần bắt chuyện với nhiều bạn hơn. Nhưng tôi cũng nhận ra một điều: không phải tất cả đều vui vẻ. Có một ánh mắt nào đó - sắc lạnh, khó chịu - đang nhìn tôi với sự bất bình. Tôi không kịp nhớ tên người ấy, chắc tại vì quá hấp tấp.

Và còn một điều nữa khiến tôi thấy là lạ: hai người còn lại - cậu con trai với đôi mắt đỏ và cô gái đi cùng cậu - cứ nhìn tôi bằng ánh mắt khó hiểu. Không hẳn ghét bỏ, cũng không phải khinh thường, mà là một sự cảnh giác, như thể tôi mang trong mình một điều gì đó mà chính tôi cũng chẳng hiểu nổi.

\pov{Haragen Kyuen}

Ngay khi Reiko-san gật đầu đồng ý, cô nàng tóc nâu với mắt xanh nước đậm ấy - Shino - liền nhìn về phía tôi và Saikyou. Đôi mắt long lanh ấy khiến tôi thoáng khựng lại.

Tôi đã từng thấy đôi mắt đó-- không, chính xác là cảm giác đó. Sự cầu xin tuyệt vọng. Hệt như cách mà Sora-san đã cầu xin chúng tôi trong giấc mơ kia. Trước khi tôi và Saikyou bừng tỉnh và trở về thế giới thực.

Một sự trùng hợp\dots\ quá mức đáng sợ.

Tôi khẽ nghiêng người, thì thầm với em gái:
''Saikyou, em thấy không? Ánh mắt đó\dots\ giống hệt lúc trong giấc mơ.''

Saikyou nhíu mày, nhưng rồi khẽ gật.

Tôi thở dài, lẩm bẩm trong lòng: Nếu Shino không phải là một linh hồn cực đoan chưa thể rời đi, thì có lẽ cô ấy chỉ là một nữ sinh bình thường. Một cô gái nhút nhát, e thẹn, và có chút đáng thương.

Trước đôi mắt ấy, tôi không cảm thấy sát khí nào cả. Khác xa với Shino mà chúng tôi từng đối diện ở hành lang: lạnh lẽo, u ám, đầy ác ý. Cô gái đứng đây chỉ là một Shino. Một nữ sinh bình thường.

''\hl{Nii-chan,}'' Saikyou khẽ kéo tay áo tôi, ''\hl{hay là ta\dots{} đồng ý đi?}''

Tôi nhìn sang. Shino vẫn chưa rời mắt khỏi chúng tôi, ánh nhìn run rẩy nhưng tràn đầy hy vọng. Tôi hít một hơi rồi gật đầu.

''\hl{Được. Nhưng phải cẩn trọng. Ta không biết cô ấy sẽ lại giở trò gì với chúng ta đâu.}''

Thế là, cả hai anh em chúng tôi đều đồng ý kết bạn.

''Được chứ, mình không ngại đâu,'' tôi nói.

''Mình cũng thoải mái về chuyện này mà,'' cô em tôi theo sau, ''nên là đừng có căng như thế chứ.''

Tôi thấy ánh mắt Shino sáng lên, như một ngọn đèn vừa được thắp. Cảm giác ấy\dots\ thật lạ lùng.

Tôi và Saikyou liếc nhìn nhau. Cả hai đều nghĩ đến cùng một điều: ít nhất, Shino ở đây không có ý định làm hại chúng tôi. Nhưng cho dù vậy\dots\ chúng tôi vẫn phải đề phòng.

\pov{Haragen Reiko}

Tôi vẫn còn đang loay hoay với những câu hỏi dồn dập từ các bạn trong lớp. Tất cả mọi người xung quanh dường như đều tò mò, muốn biết tôi là ai, đến từ đâu, thích gì\dots\ khiến tôi vừa thấy bối rối, vừa cảm thấy được chào đón.

Và rồi, bất ngờ, có một bóng người đứng dậy.

Tôi ngẩng lên, và thấy cô gái đó, cô gái có mái tóc nâu để xõa và đôi mắt xanh. Tôi nghe đồn rằng cô ấy luôn thù lù một mình, hầu như không bao giờ bắt chuyện với ai.

Thế mà bây giờ, cô ấy lại tiến về phía tôi, đôi mắt sáng lên như đang gắng gượng một điều gì đó. Cả lớp thoáng chững lại, vài tiếng xì xào nổi lên. Tôi hơi giật mình khi cô chìa tay ra, nắm lấy tay tôi.

''Cậu\dots\ ờm\dots\ Xin lỗi! Cậu có muốn làm bạn với mình không?''

Câu nói ấy run rẩy, ngập ngừng. Tôi nhận ra bàn tay cô lạnh và hơi run, nhưng lại có một sự quyết tâm kỳ lạ ẩn trong đó.

Không hiểu sao, tôi khựng mất một nhịp. Tôi cảm nhận được sự thay đổi trong bầu không khí của lớp. Như thể mọi người đang nhìn vào chúng tôi.

Đặc biệt là ánh mắt của một cô gái ngồi gần đó\dots\ Nanase-san, thì phải? Ánh mắt lạnh, sắc, như muốn xuyên thủng cả hai chúng tôi.

Tôi thấy nhói trong lòng ngực. Tôi không biết tại sao, nhưng tôi cảm thấy rõ ràng: cô ấy không hề vui khi thấy tôi nắm tay cô gái này. Hai người họ không lẽ có mối thù hằn gì sao?

Tôi mỉm cười, cố giữ giọng tự nhiên nhất có thể. ''Được chứ. Rất sẵn lòng.''

Ngay lập tức, đôi mắt cô ấy ánh lên, gương mặt cô như bừng sáng. Ánh mắt cô long lanh, đôi môi run run hé nụ cười nhỏ mà tôi nhận ra: nó không phải nụ cười của một người hay được chú ý, mà là nụ cười của một ai đó lần đầu tiên tìm được chỗ đứng cho mình.

Ngay lúc ấy, tiếng ồn ào vang lên.

''Ôi trời, Shino-chan chủ động kìa!?'' Honoka-san reo lên, giọng vừa ngạc nhiên vừa thích thú, rồi quay sang Yae.

Yae cười hồn nhiên, nhưng tôi nghe thấy sự bất ngờ rõ rệt trong giọng: ''Không thể tin được\dots\ Misora Shino-chan cuối cùng cũng mạnh dạn rồi nhỉ?''

À\dots\ Misora Shino-san. Thì ra đó là tên của cô. Tôi khẽ nhẩm lại trong đầu để ghi nhớ.

Tôi thấy nhiều bạn học liếc nhìn về phía chúng tôi, có người nhoẻn cười, có người thì thì thầm to nhỏ. Nhưng nổi bật nhất vẫn là ánh nhìn sắc lạnh của cô gái tên Nanase. Một cảm giác khó chịu thoáng lướt qua lồng ngực tôi.

Tôi khẽ siết tay cô ấy, rồi ngẩng lên, cúi đầu nhẹ về phía cả lớp: ''Hy vọng mọi người sẽ giúp đỡ ba chúng mình thật chu đáo. Mong được chỉ bảo.''

Bầu không khí xung quanh thoáng thay đổi. Vài tiếng cười khẽ, vài cái gật đầu đáp lại. Shino vẫn đứng cạnh tôi, đôi má hơi ửng đỏ, như thể không tin được mình vừa làm điều đó. Tôi bất giác cười với phản ứng của cô ấy, dường như Shino-san là một người rất hướng nội, vậy nên việc cô ấy kết thân với tôi không khác gì vừa đạt được một thành tựu lớn lao như thế.

Nhưng sâu trong lòng, tôi không thể quên ánh nhìn sắc lạnh ban nãy. Một linh cảm mơ hồ mách bảo rằng, từ khoảnh khắc này, mọi chuyện sẽ không còn đơn giản nữa.

Trong khoảnh khắc ấy, tôi tự hỏi: ''\textit{Mình vừa vô tình chạm vào một điều gì đó phức tạp hơn mình nghĩ chăng?}''

\begin{center}
  \uline{Vài ngày sau}
\end{center}

Những ngày sau đó trôi qua lặng lẽ nhưng cũng đầy mới mẻ. Hai anh em Saikyou-san và Kyuen-san tiếp tục ở nhờ nhà của tôi. Dần dà, khi tôi bắt đầu sống với họ, tôi mới biết được những mặt khác trong con người họn. Kyuen-san thì bận rộn với công việc riêng, nhưng nói trắng ra thì cậu là một người hướng nội, giống như tôi ngày trước, còn Saikyou-san thì khá thoải mái, thỉnh thoảng còn trò chuyện với tôi về những chuyện vụn vặt ở trường. Ngôi nhà gỗ đơn sơ của tôi tuy không lớn nhưng ấm cúng, và cái cảm giác có người đồng hành khiến tôi bớt thấy lạc lõng trong môi trường vừa quen vừa lạ.

Ở trường, tôi dần làm quen với nhịp học tập và không khí lớp. Shino-san thường hay nhìn tôi với ánh mắt lấp lánh, như thể chỉ cần có tôi ở đó là cô ấy đã có thêm can đảm. Thỉnh thoảng Kyuen-san và Saikyou-san cũng góp chuyện, khiến tôi cảm thấy việc hòa nhập không quá khó khăn.

Nanase-san thì vẫn giữ vẻ lạnh nhạt, nhưng tạm thời tôi gạt bỏ để không khiến bản thân quá căng thẳng. Dẫu vậy, tôi biết rõ một điều: tôi đang dần thích nghi, dần tìm thấy chỗ đứng của mình nơi đây.

Một buổi sáng nọ, khi tôi vừa đặt cặp xuống bàn thì bất chợt Shino lao đến. ''A-A-nô\dots''

Tôi ngẩng lên, mỉm cười nhẹ. ''Ồ, chào buổi sáng, Shino-san. Có chuyện gì sao? Cứ bình tĩnh mà nói.''

''À, ừ. Ờm\dots\ Reiko-san, cậu có chút thời gian rảnh không?'' giọng cô ấy run run.

''Cậu cứ việc. Ta vẫn còn ít nhất vài phút nữa trước giờ truy bài mà.''

Shino-san cười trừ, hai tay vặn vẹo như đang tự trách. ''Cảm ơn cậu. Mình xin lỗi vì đã đột ngột tới gần bạn như vậy. Thực ra thì, mình đã muốn nói chuyện với cậu từ lâu rồi.''

Tôi khẽ nghiêng đầu. ''Nghe các bạn trong lớp nói rằng cậu là một người nhút nhát, đúng chứ?''

Cô nàng mím môi, rồi gật đầu. ''Ừ-Ừ. Mình vốn dĩ\dots\ không tự tin vào bản thân mình lắm đâu. Nhưng mà nhé\dots\ mình không muốn cứ mãi như thế. Mình muốn thay đổi bản thân, nên\dots\ mình đã lấy hết can đảm để bắt chuyện với cậu.''

Tôi bật cười khẽ, cảm giác đồng cảm dâng lên. ''Nhìn cậu tự thoát khỏi vỏ bọc của mình cũng là đáng nể rồi. Mình cũng từng là một người nhút nhát, nên mình hiểu được cảm giác đó. Dần dần cậu cũng sẽ quen cả thôi.''

Ánh mắt cô ấy sáng rỡ lên, như thể vừa tìm thấy một người có cùng tâm tư.

''Vậy, có chuyện gì mà cậu muốn nói với mình sao?'' tôi gợi mở.

''À, chả là thế này\dots\ Nếu mà cậu không phiền thì\dots''

''Ối chà, Shino-san, trông cậu có vẻ nhiệt tình quá nhỉ.'' Một giọng nữ vang lên, ngọt ngào nhưng kèm theo sự mỉa mai nhẹ.

Tôi ngẩng lên, ngước sang người vừa tới bên cạnh: ''À, Furukawa-san.''

''Gọi tôi là Nanase được rồi,'' nụ cười nàng quý phái thoáng hiện trên môi, rồi cô quay sang Shino-san. ''Cậu quả là phấn khích thật đấy, nhưng tôi nghĩ cậu nên kiềm chế lại một chút đi.''

''M-Mình xin lỗi!'' Shino-san cúi đầu rối rít.

Tôi xua tay. ''Không sao, không sao. Nhưng mà Nanase-san nói cũng đúng đấy. Shino-san cũng cần bình tĩnh lại một chút.''

''À, không, mình không có ý\dots\ Chỉ là\dots'' Shino luống cuống, mặt đỏ bừng.

Nanase-san bật cười khẽ. ''Phản ứng của cậu trông đáng yêu ghê.'' Tôi cười trừ, gật gù với quan điểm của cô ấy. Nhìn cô ấy đỏ mặt vì ngượng làm tôi muốn véo má cô ấy ghê.

''Phải chăng cứ mỗi lần một ai đó bắt chuyện lần đầu là cậu cứ sồn sồn như vậy sao? Thật là\dots'' Nanase thở ra nhẹ, giọng như đang trách yêu.

''Thành thật thì, mình cũng từng trải nghiệm cảm giác đó, nên mình không thấy điều đó kỳ quặc đâu,'' tôi xen vào. ''Chỉ là, Shino-san cũng cần có thời gian để tĩnh tâm chứ.''

Cô nàng tóc xanh cười, rồi bất chợt quay sang tôi, như vừa nhớ ra điều gì quan trọng: ''À phải rồi, Reiko-san. Cảm nhận của cậu thế nào về bộ phim mà tôi đã giới thiệu cho cậu hôm nọ?''

Tôi suy nghĩ một thoáng, rồi đáp một cách ráo hoảnh: ''Mình thấy thật sự ấn tượng. Những năm 1940, nhất là giai đoạn sau Thế chiến II, dòng phim chính kịch phát triển mạnh, thể hiện rõ nỗi đau, sự mất mát và khát vọng tái thiết của con người. Mình cảm nhận được cái chất hiện thực sâu sắc của nó. Nó làm mình nhớ tới những bộ phim tâm lý xã hội mà mình đã từng xem trước đó, tuy bối cảnh khác, nhưng tinh thần thì giống hệt: phản ánh một thời đại đầy biến động.''

Con gái của doanh nhân hơi nhướn mày, rồi nở nụ cười mãn nguyện. ''Cậu có con mắt khá tinh tế đấy. Nếu cậu thực sự hào hứng, tôi có thể rủ cậu đi xem bộ phim ấy vào cuối tuần tới. Cậu có thể rủ cả Kyuen-san và Saikyou-san đi cùng.''

''Thế ư?'' tôi hơi ngạc nhiên.

''Thật mà,'' Nanase-san gật đầu, giọng ngọt lịm. ''Tôi cũng đã sắp xếp một bữa ăn tối sau khi ta xem phim nữa. Ta sẽ ăn ở một nơi ngoại thành, nơi mà tôi quen thuộc và hợp khẩu vị với tôi nhất. Tôi hy vọng là nó cũng sẽ hợp với sở thích của cậu.''

''Mình thì dễ tính mà. Chắc là cũng sẽ ổn thôi. Chỉ là\dots\ không biết liệu điều đó có hơi khoa trương không. Mình không quen lắm với mấy thứ sang trọng như thế này.''

Cô bạn tóc xanh khẽ cười với phản ứng ngập ngừng của tôi. ''Đúng là dễ thương thật.''

''Ờm, m-mình muốn hỏi cậu một chuyện\dots'' Shino-san bất chợt chen vào.

Nanase-san liếc qua người vừa lên tiếng, rồi giả vờ sực nhớ. ''Ồ, tôi thực sự rất xin lỗi. Tôi sơ suất quá, không để ý cậu ở đó.''

Cô ngắm Shino-san một hồi lâu, rồi nghiêng đầu. ''Cậu là\dots\ cái gì ấy nhỉ?''

Tôi đổ mồ hôi hột, cười gượng. Cậu thực sự không để ý tới mọi người sao\dots\ hay là cố tình vậy nhỉ?

''Mình là Misora Shino,'' Shino vội vã đáp. ''Misora (Mỹ Thiên) trong `bầu trời đẹp', và Shino với `thời gian' và--''

''Tên của cậu nghe đẹp thật đấy, Misora-san.'' Nanase khẽ cười. ''Có lẽ nào\dots\ cậu cũng có hứng thú với phim ảnh?''

''Tất nhiên! Mình thường chỉ đọc sách, chưa bao giờ tiếp xúc với phim ảnh bao giờ, và mình đã luôn muốn xem nó một lần trong đời! Vậy nên\dots''

Cô nàng tóc xanh chợt ngước ra chỗ khác, im bặt. Tôi thoáng cảm nhận có gì đó bất thường trong thái độ ấy, như một lớp mặt nạ vừa trượt xuống.

''\dots mình có thể tham với các cậu vào cuối tuần tới được không?'' Shino-san hỏi, ánh mắt đầy hy vọng.

Ngay khi cô ấy dứt câu, một tiếng thở dài từ Nanase-san thốt ra, kèm theo đó là một cái tặc lưỡi không vui. ''Phiền thật đấy.''

''Hể?'' cả tôi lẫn Shino cùng thốt lên.

Nhưng ngay sau đó, cô ấy trở lại với nụ cười quý phái thường ngày. ''Tôi thực sự xin lỗi nhé, Misora-san, nhưng tôi không có vé cho cậu đâu. E rằng cậu không thể đi cùng với bọn tôi được rồi.''

Shino-san trùng giọng, nét mặt hiện lên sự thật vọng rõ rệt: ''À, tiếc quá. Chán thật\dots''

''Có lẽ cậu có thể tham gia với bọn tôi vào một lúc khác, nhỉ?'' Nanase-san thêm vào, vẫn nhẹ nhàng như thể không có chuyện gì.

Tôi mỉm cười, khẽ vỗ vào vai cô bạn mới quen: ''Không sao đâu. Khi nào xem xong, mình sẽ kể lại toàn bộ nội dung cho cậu nghe, từng chi tiết một. Như vậy cậu vẫn sẽ cảm nhận được câu chuyện.''

Đôi mắt cô ấy bỗng chốc sáng lên, lấp lánh niềm vui. ''Thật ư? Vậy thì tuyệt quá! Cảm ơn cậu nhiều, Reiko-san!''

Nanase-san thì chỉ khẽ nhướng mày, nụ cười trên môi có chút gượng gạo.

Trong lòng tôi, một suy nghĩ lặng lẽ dấy lên: ''\textit{Nanase-san\dots\ có vẻ như cô ấy hơi xấu tính với Shino-san. Nụ cười ấy lịch sự là thế, nhưng mình không thể xóa bỏ cảm giác rằng ẩn dưới đó là một điều gì khác.}''

\pov{Haragen Saikyou}

Tôi ngồi ở cuối lớp, lặng lẽ quan sát khung cảnh phía trước. Cuộc trò chuyện giữa Nanase-san, Reiko-san và Shino-san nhìn qua thì chẳng có gì khác thường, nhưng trực giác con gái mách bảo tôi rằng ở đó tồn tại một thứ gì đó không ổn.

Furukawa Nanase. Với dáng vẻ quý phái, giọng nói luôn toát lên sự dịu dàng, trước nay vẫn tỏ ra thân thiện với hầu hết bạn học. Ngay cả khi trò chuyện với tôi hay Nii-chan, cô ấy cũng giữ một thái độ nhã nhặn đến mức chẳng ai bắt lỗi được. Ấy vậy mà, khi đối diện Shino-san, giọng nói ấy trở nên lạnh lùng hơn, mang theo thứ gì đó giống như sự dè bỉu được giấu khéo dưới nụ cười. Đặc biệt, tôi để ý thấy: bất cứ khi nào cuộc trò chuyện xoay quanh Reiko-san, ánh mắt của cô ta lại thoáng qua vẻ gì đó khó đoán, gần như là gai góc. Nó không đáng sợ như Shino-san ở hành lang, nhưng cũng đủ khiến cho đối phương phải đổ mồ hôi hột.

Tôi khẽ nheo mắt, rồi liếc sang Nii-chan đang khoanh tay tựa vào bàn, đôi mắt nửa như trầm ngâm, nửa như dè chừng.

''\dots Nii-chan, anh có thấy Nanase-san có gì đó lạ không?'' tôi thì thầm.

Anh ấy chỉ gật đầu, giọng nhỏ nhưng chắc chắn: ''Có. Với người khác thì bình thường, nhưng với Shino-san thì\dots\ rõ ràng là xấu tính hơn. Và khi liên quan đến Reiko-san, anh cảm thấy có chút gì đó\dots\ ghen tuông ở trong.''

Tôi thoáng ngạc nhiên. Ghen tuông? Quả thực, cái cách mà Nanase-san cắt ngang lời Shino-san, hay làm ngơ rồi quay sang Reiko bằng thái độ hoàn toàn khác, giống như đang cố tình dựng lên một ranh giới.

Trong khi ấy, Reiko lại tỏa ra một thứ gì đó thu hút. Không chỉ cô ta, mà ngay cả các bạn học khác cũng thường vô thức hướng mắt về cô. Điều đó khiến tôi nhớ đến những ngày ở thời hiện đại, khi Saki, Touka, Suzune, Hanabi và Inari từng chú ý đến hai anh em chúng tôi ngay từ lần đầu tiên gặp gỡ. Cảm giác ấy giống hệt,  một thứ sức hút kỳ lạ kéo ánh nhìn của mọi người về phía một người duy nhất.

Tôi mím môi, ánh mắt vẫn không rời con gái của nhà kinh doanh. Có điều gì đó ở cô gái ấy\dots\ một mặt nụ cười dịu dàng, mặt khác là sự khó đoán, như thể dưới lớp vỏ hoàn hảo kia còn ẩn chứa thứ gì đó không ai dễ dàng nhận ra.

''Chúng ta cứ quan sát đã,'' Nii-chan khẽ nói, như đoán được ý nghĩ trong đầu tôi.

Tôi gật nhẹ, thu ánh mắt về. Đúng lúc ấy, Yae-san và Honoka-san vui vẻ bước đến, kéo tôi và anh trai vào một cuộc trò chuyện sôi nổi về bài kiểm tra sắp tới. Cả lớp nhanh chóng trở lại với không khí quen thuộc: ồn ào, nhộn nhịp, giống như chưa từng có gì xảy ra.

Thế nhưng, trong lòng tôi vẫn còn một nỗi lấn cấn. Furukawa Nanase\dots\ cô ấy rốt cuộc là người như thế nào?

\begin{center}
  \uline{Ngày hôm sau.}
\end{center}

Trời lại tiếp tục đổ mưa.

Những hạt mưa lăn dài trên khung cửa sổ tàu hỏa, tạo thành những vệt dài uốn lượn. Reiko-san khẽ thở dài, vừa bước vào sân trường vừa lẩm bẩm: ''Lại mưa nữa rồi\dots\ Từ lúc mình nhập học đến giờ, ngày nào cũng thế cả.''

Tôi mỉm cười, cố gắng xua tan cái không khí u ám ấy. ''Có khi bọn tôi phải mua thêm một cái ô thật, chứ cứ dựa cậu thế này thì có hơi phiền.''

Cô bạn tóc hai bím bất giác quay sang, đôi mắt long lanh như ánh sáng xuyên qua màn mưa, rồi mỉm cười dịu dàng: ''Không cần đâu. Mình cũng chẳng ngại khi cả ba người đi chung ô. Với cả\dots\ ô của mình cũng đủ rộng cho ba chúng mình mà.''

Câu nói ấy khiến tôi khẽ cười theo. Nii-chan cũng vậy, ánh mắt anh thoáng qua một tia mềm mại. Bỗng chốc, cái rét se lạnh của cơn mưa chẳng còn làm tôi khó chịu nữa.

Không thể ngờ, Reiko-san lại tốt bụng như thế.

Nhưng rồi, khi chúng tôi bước vào lớp, bầu không khí lại đổi khác. Trên bàn của Misora Shino-san, một lọ hoa bỉ ngạn đỏ rực nằm chình ình, nổi bật đến mức khó lòng mà bỏ qua.

Cô nàng mắt hai màu giũ ô xuống, nhìn qua chiếc lọ ấy, vẻ mặt ngạc nhiên: ''Sao lại có thứ này ở đây\dots?''

Tôi và Nii-chan thì chết lặng. Tim tôi bất giác đập nhanh, những ký ức mơ hồ ập về: bông hoa đỏ rực ấy đã từng hiện ra trong giấc mơ quái dị ở trong giấc mơ, và một lần nữa, tôi còn thấy nó ở hành lang. Cả hai lần, đều chẳng phải điềm lành gì. Và ở đây, nó lại xuất hiện trên bàn của một người vẫn còn sống\dots\ Tôi biết quá rõ ý nghĩa của hành động ấy.

Reiko-san thoáng run tay, nhưng lập tức cắn môi rồi vội vã nhấc lọ hoa ra ngoài, như thể chỉ cần nhìn nó thêm một chút nữa thôi cũng khiến cô khó chịu. Trong khoảnh khắc ấy, tôi liếc sang phía Nanase-san.

Cô ta vẫn ngồi ở chỗ của mình, lưng thẳng tắp, dáng vẻ tao nhã như thường ngày. Nhưng ngay khi bắt gặp ánh mắt của tôi, tôi kịp nhận ra\dots\ cô vừa liếc ra phía sau. Chỉ trong một khoảnh khắc ngắn ngủi, rồi cô thản nhiên quay về bàn, như thể chưa từng quan tâm đến bất cứ điều gì.

Sự nghi ngờ trong lòng tôi và Nii-chan bỗng hướng thẳng về phía cô ấy.

Reiko quay lại chưa kịp ngồi xuống thì Shino-san bước vào lớp. Trông thấy cô ấy, cô nở nụ cười thân thiện, giọng trong trẻo vang lên: ''Chào buổi sáng, Reiko-san!''

Cô gái tóc hai bím khựng lại, trên gương mặt hiện rõ sự bối rối. Cô bạn ngồi cạnh bàn nghiêng đầu, lo lắng hỏi: ''Cậu sao thế? Trông cậu có vẻ\dots\ căng thẳng vậy?''

Ánh mắt Reiko-san vô thức hướng về phía chúng tôi. Tôi và Nii-chan lập tức ra hiệu bằng một cái lắc đầu nhẹ, bảo cô ấy đừng nói gì.

''À\dots\ không có gì đâu,'' cô ấy vội lảng đi, nặn ra một nụ cười gượng gạo. ''Chắc tại mình\dots\ thấy khó chịu với thời tiết dạo gần đây thôi.''

Tôi thoáng liếc sang Nanase-san. Gương mặt cô vẫn điềm tĩnh, nhưng ánh mắt thoáng tối lại. Tôi nhớ chính cô ta nói rằng cô ghét sự dối trá, và lời nói vừa rồi của Reiko-san chẳng khác gì một mũi gai.

Nhưng cô ta không nói gì. Chỉ im lặng, đặt tay lên bàn, đôi môi mím chặt trong nụ cười đầy lịch thiệp.

Một khoảng lặng lạ lùng phủ lên cả ba chúng tôi.

\ct

Giờ ra chơi đến, lớp học nhanh chóng tách thành từng nhóm nhỏ. Shino-san kéo Reiko-san ra một góc gần cửa sổ, ríu rít trò chuyện như thể đã quen biết từ lâu. Nanase-san thì lại thản nhiên tiến đến chỗ Nii-san, cầm sách vở hỏi han vài bài tập còn khúc mắc. Còn tôi thì vẫn ngồi tại bàn, sắp xếp lại sách vở chuẩn bị cho tiết tiếp theo. Mọi thứ, thoạt nhìn, vẫn diễn ra rất bình thường.

Thế nhưng, trong đầu tôi vẫn lẩn quẩn hình ảnh lọ hoa bỉ ngạn đỏ ban sáng. Dù bản thân chẳng muốn vội vàng kết luận, nhưng linh cảm mách bảo thủ phạm chỉ có thể là Nanase-san. Cái cách cô ấy thoáng liếc ra sau rồi quay phắt đi\dots\ rõ ràng không phải vô tình. Nhưng chứng cứ thì không có, mà bản thân cũng chỉ mới chuyển lớp ba ngày. Tốt hơn hết, trước mắt, tôi vẫn phải giữ im lặng và quan sát thêm.

Ngay lúc tôi còn đang trầm ngâm thì một giọng nói vang lên từ bàn phía trên: ''Saikyou-san, Shiina-san, hai cậu rảnh chứ? Lại đây một chút.''

Tôi ngước lên. Là Yae-san. Còn Shiina-san, vừa mới tỉnh dậy sau tiết Sử dài lê thê, uể oải dụi mắt, cất tiếng đáp: ''Yo, Yae.''

''Yae-san, có việc gì gọi chúng tôi sao?'' tôi nghiêng đầu hỏi lại.

''Cậu nói rằng `có việc này muốn nói', ý là sao thế?'' cô bạn Sư tử thêm vào, giọng vẫn ngái ngủ.

Yae khẽ gật đầu, trông hơi do dự, nhưng rồi bắt đầu:

''Ừ, đúng là tớ có chuyện muốn nói thật. Chuyện là thế này-- Ể-Ể!? Shiina-san! Tay cậu bị thương rồi kìa! Cậu ổn chứ? Chuyện gì thế này?''

Theo phản xạ, tôi lập tức nhìn xuống bàn tay phải của Shiina-san. Quả nhiên, cả bàn tay cô ấy được quấn băng, nhưng sơ sài và chắp vá đến nỗi trông có vẻ còn đau hơn là bảo vệ.

''Ồ, cái này có là gì đâu,'' cô bật cười hờ hững. ''Giao lưu võ thuật một chút với con mắm ở trường phía Bắc thôi ấy mà.''

''Ý cậu là cái đứa lúc nào cũng huênh hoang, tự cho mình là nhất đấy hả? Tôi là tôi muốn đấm bỏ mẹ nó rồi đấy,'' tôi khẽ nhíu mày.

''Nó chứ ai. Tự dưng muốn lên mặt đòi tỉ thí với tôi, và thế là\dots\ đấy. Tôi thắng rồi. Vài vết vớ vẩn cũng chẳng sao. Mà tại sao cậu muốn tẩn nó thế?''

Tôi đáp thẳng, giọng dứt khoát: ''Nó dám xúc phạm Nii-chan của tôi.''

Shiina-san im lặng trong giây lát, rồi nhếch môi: ''\dots Ồ. Cậu định đánh thật á?''

''Nếu không vì Nii-chan cản thì tôi đã hành cho ra bã rồi. Anh ấy bảo không để tâm, nhưng thử hỏi tôi có chịu để yên không?''

Ánh mắt cô ấy dịu lại, pha chút thấu hiểu: ''\dots Tôi hiểu cảm giác của cậu. Cậu yêu Nii-chan của mình lắm nhỉ?''

''Chứ sao! Nhưng mà cậu đã xử lý nó rồi thì\dots\ thôi, coi như xong.''

''Vấn đề không phải ở đó!'' Yae đột ngột chen ngang, giọng trách móc. ''Con gái con đứa sao lại phải dây dưa với mấy vụ bạo lực như thế?''

''Thì tôi có muốn thế đâu. Nhưng nếu không ai bảo vệ anh trai tôi thì ai sẽ làm?'' tôi phản bác.

Shiina-san nhướng mày, cười khẩy: ''Saikyou-san nói đúng. Sao cậu không dẹp bỏ mấy cái định kiến giới tính ấy đi, Yae? Thời nay, trai hay gái thì cũng phải biết tấn công.''

''Tớ thì không đồng ý. Một quý cô đúng nghĩa phải--''

''Trời ạ, lại bắt đầu rồi\dots'' Shiina-san chống cằm, chẹp miệng một tiếng chán nản.

Tôi bật cười, hỏi nhỏ: ''\hl{Yae-san lúc nào cũng như thế à?}''

Cô bạn gật gù, liếc Yae-san rồi đổi giọng châm chọc một cách đầy thẳng thắn: ''Lúc nào cũng `con gái thì phải nho nhã, nhẹ nhàng', rồi lại `con gái thì phải xinh đẹp, tài giỏi'. Ngày quái nào cũng thế.''

''Tớ nói đúng mà! Cậu lúc nào cũng để ngoài tai tớ hết!'' Yae-san phụng phịu.

''Thôi, dẹp mấy trò quý cô vô nghĩa này đi. Vào việc chính. Cậu gọi bọn tôi có việc gì?'' cô bạn chán nản cắt ngang.

Thấy vậy, cô gái thể thao hơi ngập ngừng, rồi mới chợt nhớ: ''À, đúng rồi. Nãy giờ tớ cứ để ý tới Nanase-cha-- à không, Nanase-san ấy. Hai cậu có thấy cô ấy\dots\ khác thường không?''

''Ý cậu là sao?'' Shiina-san nhướng mày.

''Tớ cảm thấy dạo gần đây cậu ấy\dots\ hơi khác so với mọi khi.''

''Khác? Tôi thấy cậu ta vẫn thế mà,'' cô nhún vai.

Tôi liếc sang Nanase-san, đang tươi cười trò chuyện với Reiko-san ở góc lớp, còn Shino-san thì đã quay về bàn với tâm trạng phấn khởi. Trong đầu, tôi nhớ lại cuộc nói chuyện hôm qua với cô gái tóc nâu, và cả lọ hoa đỏ lúc sáng.

Nhưng lời thoát ra khỏi miệng lại lạnh nhạt: ''\dots Ngoài việc cô ấy đang nói chuyện rất vui vẻ với mọi người thì tôi chẳng thấy có gì bất thường.''

''Ý tớ không phải là ngoại hình\dots'' Yae-san cúi đầu, do dự rất lâu, rồi mới hạ giọng. ''Mà là\dots\ Chuyện là hôm trước, tớ thấy Nanase-san cùng với một cậu bạn lớp bên, cái cậu mọt sách hay đeo kính ấy. Gần cả tuần nay cậu ấy nghỉ học, mọi người nói là bị ốm\dots\ Nhưng hôm tớ trực nhật muộn, tớ đã thấy cảnh tượng đó.''

Nói đến đây, cô ghé sát lại, thì thầm với chúng tôi. Giọng cô run run, như không muốn ai khác nghe thấy. Chỉ vài chữ thôi, nhưng đủ khiến cả hai chúng tôi rùng mình.

Ánh mắt Nanase-san mọi khi hiền hòa khi ấy\dots\ trống rỗng. Còn cậu bạn kia, mặt cắt không còn giọt máu. Và vị trí họ đứng\dots\ là con đường dẫn ra cánh rừng Aogami.

Shiina-san nhíu mày, bật thốt: ''Hả? Cậu ta đã làm chuyện đó sao?''

''Nghe kinh dị thật,'' tôi khẽ lẩm bẩm, lưng bất giác lạnh toát.

''Tớ không nói đùa đâu,'' Yae-san siết chặt tay, run rẩy. ''Ánh mắt ấy\dots\ đến giờ tớ vẫn còn ám ảnh. Mọi ngày cậu ấy có cư xử như thế đâu.''

Cô bạn Sư tử khoanh tay, trầm ngâm: ''\dots Được rồi. Tôi sẽ để mắt tới Nanase. Saikyou thì sao?''

''Tôi thì luôn quan sát mọi người rồi. Vậy nên, càng phải chú ý hơn tới cô ta. Nhưng hành động vội bây giờ\dots\ có lẽ không khôn ngoan đâu,'' tôi đáp, giọng nghiêm túc.

Yae gật đầu, thở dài: ''\dots Ừ, chắc vậy. Cảm ơn hai cậu.''

Chúng tôi trở về bàn đúng lúc chuông báo vào giờ vang lên. Nhưng trong lòng tôi thì không yên chút nào. Cảnh tượng cô bạn Thể lực vừa kể cứ như một mũi kim cắm sâu vào đầu.

Tôi khẽ nghiêng sang hỏi anh tôi: ''Nii-chan, Nanase-san vừa nãy đã hỏi anh chuyện gì thế?''

Anh bình thản đáp: ''À, chỉ vài chỗ trong bài lịch sử thôi.''

''Thái độ của cô ấy\dots\ có gì lạ không?''

Anh suy nghĩ một chút, rồi lắc đầu: ''Bình thường cả thôi. Không có gì kỳ lạ cả.''

Mình khẽ thở phào, nhưng trong lòng vẫn còn ngổn ngang. Nhất là khi hình ảnh bình hoa bỉ ngạn đỏ từ sáng sớm lại hiện về. Không thể lơ là.

Tốt hơn hết, tôi sẽ trông chừng Nanase-san thật cẩn thận.

\pov{Haragen Kyuen}

Giờ nghỉ trưa đến, lớp học rộ lên những tiếng trò chuyện và tiếng ghế xê dịch. Không khí lúc nào cũng náo nhiệt vào khoảng thời gian này.

Reiko-san và Shino-san ngồi cạnh cửa sổ, cùng nhau mở hộp cơm ra ăn. Tôi thoáng nhìn qua, thấy cả hai nói chuyện cởi mở hơn hẳn so với ngày hôm qua.

Ở phía bên kia lớp, Saikyou ngồi tụ tập cùng Yae-san, Shiina-san và Nanase-san. Nghe thì như đang vui đùa, nhưng tôi biết thừa rằng em gái mình đang ''cài cắm'' để tiện theo dõi thái độ của cô gái tóc xanh kia.

Còn tôi, như một thói quen, ở lại lớp để ăn trưa một mình. Trước đó, Saikyou cũng có rủ tôi đi cùng, nhưng tôi từ chối. Chỉ là bản thân tôi không quen chen vào mấy nhóm đông người.

Tôi lặng lẽ mở hộp cơm mà Reiko-san đã chuẩn bị từ tối hôm qua. Mùi thơm vừa tỏa ra thì một giọng nói vang lên từ bên cạnh: ''Này, Kyuen-san. Cậu có phiền khi tôi ngồi đây không?''

Ngẩng đầu lên, tôi thấy Sakura-san đang đứng đó với vẻ mặt điềm tĩnh. Tôi khẽ nhún vai: ''Cứ tự nhiên.''

Cô không nói thêm gì, chỉ kéo ghế ra và ngồi xuống, rồi rút trong cặp ra một cuốn sách dày cộp. Ngay khi tôi bắt đầu ăn, giọng cô lại vang lên, lần này mang theo chút trầm ngâm.

''Hừm\dots\ Xem nào, ở đây có lá bùa chống cúm này\dots''

Tôi ngước nhìn cuốn sách trên tay cô, liếc qua vài dòng chữ in chi chít.

''Rồi có cả lá bùa cầu nắng nữa\dots\ Trong này có nhiều loại lá bùa giúp ích thật.''

Cô gật nhẹ, tiếp tục lật sang trang khác. ''Lá bùa để làm một nồi súp miso thơm ngon\dots?''

Tôi sặc suýt nghẹn. ''\dots Được rồi, cái đó thì kỳ quặc đấy.''

Sakura-san đệm ra một tiếng thở dài, dường như tỏ vẻ không thoải mái. ''Không ổn rồi. Mấy cái này không phải là thứ tôi cần.''

Tôi gác đũa, phẩy tay: ''Ừ. Tôi cảm thấy chúng giống như mấy cái gì đó mê tín lắm ấy. Một vài trong số chúng thậm chí còn khiến tôi nhăn mặt này - có cả bùa để hái rau ngon, bùa để\dots\ ngất xỉu? Cái quái gì thế?''

''Tôi nghĩ cuốn sách này không có thứ ta cần đâu. Tôi cần một cái gì đó táo bạo hơn thế\dots''

Tôi ngả người ra sau, thở dài. Ánh mắt bất giác hướng lên trần lớp học. Hình ảnh lọ hoa bỉ ngạn đỏ lại chợt ùa về trong đầu: cái lọ hoa sáng nay chúng tôi phát hiện trên bàn Shino-san. Cơn mưa ở thị trấn Aogami này đã kéo dài từ lúc chúng tôi nhập học đến giờ, rồi cả câu chuyện Saikyou kể ở tiết vừa rồi: về cánh rừng mà Yae-san từng nhắc, cùng với những biểu hiện kỳ lạ của Nanase-san. Mảnh ghép nối tiếp mảnh ghép, khiến tôi càng thêm nặng lòng.

Khi tôi vừa đưa mắt về phía sau lớp, bất chợt bắt gặp Honoka-san đang rón rén bước vào từ cửa sau. Thấy tôi định mở miệng, cô vội đưa ngón tay lên môi ra dấu im lặng, rồi khẽ tạo khẩu hình: ''Để tớ tạo bất ngờ cho cậu ấy!'' Tôi khẽ gật đầu, trở lại nhìn Sakura-san, người vẫn đang mải miết lướt qua các trang sách.

''Dù gì thì mấy cái thứ dễ kiếm kiểu này cũng tầm phào thật. Cậu nghĩ sao, Kyuen-san?'' cô bạn chợt hỏi.

''Họ cũng chỉ sưu tầm được mấy loại bùa thường gặp thôi. Mấy cái cao siêu hơn thì có lẽ phải tìm ở một chỗ nào đó khác rồi.''

''Phải ha\dots'' cô trưng ra bộ mặt chán nản, rồi chợt khựng lại khi thấy một lá bùa vàng óng nổi bật trên trang giấy. ''Hửm? Một loại bùa cầu tiền sao\dots?''

Sakura-san trố mắt, dí cuốn sách sát vào mặt, lẩm bẩm: ''Giờ đúng là loại mà mình đang tìm đây này!''

Tôi khẽ cười chát: ''Trời ạ. Cứ nhìn thấy tiền là mắt lại sáng lên thế kia\dots''

Tôi liếc về phía sau. Honoka-san vẫn đang lấp ló, mắt sáng long lanh, đầy thích thú.

''Sakura-chan! Đọc cái gì mà chăm chú thế?'' tiếng cô bạn tóc vàng vang lên bất ngờ.

Sakura-san giật mình ngẩng lên, nheo mắt nhìn. ''\dots Cậu là ai?''

''Trời ạ, là tớ đây! Honoka đây! Đến bạn học trong lớp mà cậu còn không nhớ sao?'' cô vung tay múa chân, mắt trợn ra vì sốc.

''Tôi thực tình cũng chẳng để tâm mấy đâu,'' cô bạn tóc hồng nheo mắt nhìn chằm chằm một hồi, rồi búng tay một tiếng. ''À, nhớ rồi, kẻ ồn ào nhất trong lớp.''

''Đó là những gì mà cậu nhớ về tớ sao!?'' Honoka-san ngỡ ngàng. ''Cậu vẫn còn nhận ra Kyuen-kun đấy thôi!''

''Cả lớp 1-B này có mỗi cậu ta là con trai, làm gì mà tôi không nhận ra?''

Tôi đưa mắt nhìn quanh lớp, rồi gật gù: ''\dots Đúng thật.'' Từ lúc ba chúng tôi nhập học ở đây, chỉ có duy nhất mình tôi là thằng đực rựa.

''Vậy còn Saikyou-chan!?'' Honoka-san bật lại.

''Cô ấy là em gái song sinh của cậu ta, nên cũng dễ nhận ra nốt. Cậu thì là tên bày trò với cả hay tám chuyện nhất trong lớp mà, đúng không?''

''\dots Không sai!'' cô bạn tai cáo bật cười.

''Cậu ấy có khen cậu đâu mà hì hì cái nỗi gì chứ\dots'' tôi lắc đầu, chưng hửng nhìn cô gái nhiệt tình ấy.

''Thôi, mấy chuyện nhận bạn nhận bè thì để sau. Thế\~{} cậu đang đọc cái gì thế, Sakura-chan?''

Cô bạn ngồi cạn tôi giơ cuốn sách ra, mở đúng trang vừa nãy: ''Cuốn sách về các loài bùa ngải. Mượn được ở thư viện đấy.''

''Bùa ngải!?'' Honoka-san tròn mắt.

Cả tôi và Sakura-san cùng gật đầu.

''Nãy giờ bọn tôi cũng đã đọc lướt qua vài cái bùa rồi, nhưng mà trông cũng ọp ẹp lắm,'' tôi nhún vai.

Cô bạn tóc hồng nhăn mặt. Tôi vội đính chính: ''\dots Ngoại trừ cái bùa làm giàu ra.''

''Phải!'' cô gật đầu lia lịa, mắt lại sáng lên, nhưng rồi hắng giọng để lấy lại bình thản. ''Nhưng mà đó không phải là thứ mà tôi muốn nghiên cứu đến. Mà là một loại bùa khác.''

''Ái chà chà!'' cô bạn tóc vàng rạng rỡ. ''Bùa cơ à!? Có vẻ như cậu đang nhắm tới ai đó rồi đúng không?''

''Hả?'' tôi nhướng mày.

''Cậu đang nói cái quái gì vậy?'' Sakura-san đồng thời cũng nhíu mày.

''Tớ nghe đồn rằng lớp 1-A cạnh bên có một người dùng một loại bùa ngải để tiếp cận gần hơn tới người mà cậu ta thầm thương đấy. Cậu cũng đang định kiếm một cái bùa kiểu như thế hả?''

''Tôi không nghĩ Sakura-san sẽ thích nghiên cứu mấy cái sến súa như thế đâu\dots'' tôi chen ngang.

''Như Kyuen-san nói đấy.'' Sakura-san thở dài. ''Mấy cái thứ yêu đương yêu đủng đó không phải là cái mà tôi cần.''

''Ể\~{}!? Vậy mà cứ tưởng thế nào chứ\dots'' Honoka-san bĩu môi.

''Đừng có hóng hớt mấy cái chuyện như thế chứ, Honoka-san,'' tôi liếc nhẹ. ''Cậu định tự biến mình thành cái loa phường của trường chắc?''

''Bị phát hiện rồi!'' cô bạn cười trừ. ''Vậy nếu không phải là bùa yêu thì cậu định tìm cái gì ở trong đó?''

Sakura-san khép cuốn sách lại, gương mặt nghiêm hơn hẳn: ''Tôi muốn tìm hiểu các loại bùa mà có sức mạnh siêu nhiên và bí ẩn hơn như thế, kiểu như là ma thuật hay ảo ảnh, chẳng hạn!''

''\dots À, ừ, đúng rồi,'' Honoka-san gật gù, như vừa sực nhớ. ''Gia đình cậu sống là một ngôi đền mà, tớ quên mất. Vu nữ các cậu ở đó cũng phải học mấy thứ như vậy chứ nhỉ?''

''Không hẳn. Ngôi đền đó do bố mẹ tôi tiếp quản, còn tôi thì sẽ không dự định đi theo họ đâu.''

''Cậu không phải là vu nữ à?'' tôi hỏi.

''Có, nhưng chỉ khi nào có việc nhờ thì tôi mới làm,'' Sakura-san nhún vai. ''Mấy cái bùa ngải này chỉ là thấy gì đó thú vị nên mới tự tìm hiểu thôi chứ.''

''Hay nha, hay nha!'' cô bạn tóc vàng cười tủm tỉm. ''Trông cậu đang say mê thế kia thì chắc là cậu gặp phải có gì đó thú vụ nhỉ?'' Cô xán lại gần hơn. ''Thế cụ thể là cậu đang đọc cái gì vậy? Có gì thú vị không?''

''G-Gần quá, Honoka-san\dots\ Xa ra, xa ra chút!'' tôi lùi lại, giơ tay ngăn.

Cô nàng giật mình, vội lùi ra sau, cuống quýt: ''À\dots\ xin lỗi, tớ hơi hăng quá!''

''Thú vị à\dots'' Sakura-san trầm ngâm hồi lâu, rồi nói nhỏ, ánh mắt hơi nheo lại. ''Nếu cậu thực sự ổn với điều đó thì được thôi.''

Tôi ngả người ra sau một chút, vừa nhai phần cơm hộp, vừa nhìn cô ấy lật thêm vài trang sách. Cô khẽ chép miệng, rồi buông một câu nhỏ, giọng đều đều nhưng nghe ra có chút nặng nề: ''Có gì đó\dots\ bất thường. Khó tả lắm, nhưng giác quan vu nữ của tôi cứ thấy lạ. Như thể có một thứ gì đó âm u đang bao trùm quanh đây vậy.''

Tôi khựng lại, đôi đũa dừng giữa không trung. Thứ ''âm u bao trùm'' ấy\dots\ không phải tôi và Saikyou cũng đã mơ hồ cảm thấy từ khi bước chân vào ngôi trường này sao? Nhưng Sakura-san lại nói ra điều đó một cách rất bình thản, như thể chỉ đang kể chuyện thời tiết.

Honoka-san thì phá lên cười, chống hai tay vào hông, giọng trong trẻo vang khắp lớp: ''Trời ạ, làm gì có âm u âm iếc gì! Chẳng qua là mưa dầm nhiều ngày nên cậu mới thấy rờn rợn thôi. Ngồi trong lớp tối tối thì ai mà chả thấy lạ.''

Sakura-san vẫn không rời mắt khỏi cuốn sách, lắc đầu nhẹ: ''Không, tôi không nghĩ thế. Cái cảm giác này\dots\ không đơn thuần chỉ do thời tiết.''

Ngón tay cô lật đến một trang khác, ánh mắt dừng lại khá lâu. Rồi cô khẽ lẩm bẩm, gần như chỉ nói cho chính mình nghe: ''\hl{Cần một thứ gì đó\dots{} đủ mạnh để trấn yểm. Một cái gì đó có thể xua đi tà khí quanh đây.}''

Honoka-san suýt sặc, rồi ôm bụng cười khanh khách: ''Ủa, nghe ghê vậy! Cậu tính dựng cả ngôi đền trong lớp à!? Trấn yểm gì mà trấn yểm, nghe như mấy phim kinh dị ấy!''

Sakura-san chỉ liếc sang cô một cái, đôi mắt lạnh tanh, rồi lại cúi xuống, tiếp tục dán chặt vào cuốn sách. Không một nụ cười.

Tôi nhìn hai người họ, chỉ biết cười trừ. Nhưng trong lòng thì chẳng thể nào thấy nhẹ nhõm được. Cái cách mà cô ấy nói ''trấn yểm tà khí'' dẫu chỉ thoáng qua, lại khiến da gà tôi nổi lên từng đợt. Không phải ngẫu nhiên mà cô ấy cảm thấy điều đó.

Tôi khép hộp cơm lại, hít một hơi sâu, thầm ghi nhớ từng chữ Sakura-san vừa nói. Sẽ phải bàn chuyện này với Saikyou sau.

Nhưng giờ thì\dots\ trước mắt cứ ăn nốt bữa trưa đã.

\ct

Sau giờ nghỉ trưa, các tiết học tiếp diễn như thường lệ. Cả lớp 1-B vẫn ồn ào vừa đủ, với tiếng trò chuyện khe khẽ vang lên xen kẽ tiếng giảng bài trên bục. Nhưng với tôi, từng chi tiết nhỏ dường như trở nên rõ nét hơn hẳn.

Nanase-san từ đầu đến cuối vẫn giữ khuôn mặt bình thản, ghi chép cần mẫn, nhưng cứ mỗi khi ánh mắt thoáng lướt qua Shino-san, sắc lạnh lại hiện rõ. Cô gái có mái tóc nâu thì trái lại, vẫn giữ nụ cười dịu dàng, thi thoảng quay sang Reiko thì thầm đôi câu. Một hình ảnh bình thường thôi, vậy mà trong mắt tôi lại ẩn chứa điều gì đó khó lý giải.

Ngoài trời, cơn mưa kéo dài không ngớt, từng giọt đập lách tách vào khung cửa kính. Tôi nhớ lại buổi sáng nay, hình ảnh lọ hoa bỉ ngạn đỏ đặt trên bàn Shino-san, và ánh mắt thoáng liếc về phía sau của Nanase-san. Cộng thêm những lời Sakura-san lẩm bẩm về ''tà khí cần phải trấn yểm\dots'' tất cả như những mảnh ghép rời rạc, chưa thể xếp thành bức tranh hoàn chỉnh, nhưng trực giác mách bảo tôi rằng chúng không hề vô nghĩa.

Tan học. Ba chúng tôi - những người mới chuyển về đây - lại cùng nhau che chung một chiếc ô. Reiko-san khẽ than phiền: ''Trời lại mưa nữa\dots\ Cứ như thế này thì đến bao giờ mới tạnh đây\dots?''

''Chắc là đến khi tới Tết Công-gô quá,'' em gái tôi cất giọng trêu.

''Đ-Đừng có đùa như vậy chứ, Saikyou-san!'' cô gái có hai màu mắt nói, khẽ sốc nhẹ, nét mặt hiện rõ sự lo lắng.

Chúng tôi lặng lẽ bước trên con đường ướt mưa. Bầu không khí ẩm ướt nặng nề bao trùm khắp nơi, như thể đang đè nặng lên vai. Những vũng nước đọng phản chiếu bầu trời xám xịt, thỉnh thoảng lại gợn sóng khi có giọt mưa rơi xuống. Tiếng mưa rả rích đã trở thành âm thanh quen thuộc sau nhiều ngày, hòa cùng tiếng bước chân chậm rãi trên vỉa hè ẩm ướt. Không khí lạnh và ẩm thấp len lỏi qua từng kẽ hở của quần áo, khiến người ta không khỏi rùng mình. Những tán cây hai bên đường rũ xuống vì nặng nước, thỉnh thoảng lại nhỏ những giọt to xuống vai áo những người đi qua.

Về cơ bản, không khác gì từ lúc chúng tôi bị dịch chuyển tới đây.

Bỗng Saikyou nghiêng đầu nhìn tôi, ánh mắt sắc bén: ''Nii-chan, sao mặt anh lại trầm ngâm thế? Đừng bảo là không có gì nhé. Em nhìn là biết ngay mà.''

Tôi thở dài, gãi nhẹ đầu. Cô em gái này lúc nào cũng nhạy cảm đến mức chẳng giấu nổi gì. Ghé sát tai Saikyou, tôi thì thầm kể lại những gì Sakura lẩm bẩm trong giờ nghỉ trưa: về thứ âm u bao trùm, về chuyện cần có vật trấn yểm.

Đôi mắt cô bé hơi mở to, nhưng không ngắt lời. Cô chỉ khẽ gật đầu, biểu cảm căng thẳng hơn thường lệ.

Tôi nhân cơ hội hỏi tiếp: ''\hl{Nanase-san có gì bất thường không?}''

Saikyou lắc đầu, đáp gọn: ''\hl{Không có gì cả ạ. À, trừ việc cậu ta thỉnh thoảng than đau đầu nhẹ.}''

''\hl{Đau đầu\dots?}'' tôi nhíu mày.

''\hl{Vâng. Không nghiêm trọng lắm đâu, chỉ là thoáng qua vài lần. Hôm qua em còn thấy cô ấy nhăn mặt một chút, khi mà vừa nói chuyện với Shino-san.}''

Tôi im lặng giây lát rồi nói: ''\hl{Có lẽ do thời tiết thay đổi. Mưa nhiều, trời lạnh, dễ khiến người ta mệt mỏi hơn.}''

Nói vậy, nhưng trong lòng tôi chẳng thấy nhẹ nhõm chút nào. Cái gì đó mách bảo rằng đó không chỉ là cơn đau đầu lặt vặt. Tôi vừa nghĩ thế, Saikyou đã liếc sang, như thể đọc được suy nghĩ trong tôi. Không cần mất công tôi giải thích nhỉ.

Đúng lúc ấy, Reiko-san quay đầu lại, hỏi: ''Hai người thì thào cái gì thế? Nghe cứ như có bí mật ấy.''

Tôi và Saikyou thoáng giật mình, rồi đồng thanh: ''À\dots\ không có gì đâu. Việc riêng của hai bọn tôi thôi.''

Cô ấy khẽ nghiêng đầu, tỏ vẻ khó hiểu, nhưng không gặng hỏi thêm. Khi cô bước đi trước một chút, tôi nhỏ giọng: ''\hl{Có nên nói cho Reiko-san biết không?}''

Saikyou nhìn tôi, giọng chắc nịch: ''\hl{Chưa. Reiko-san chưa sẵn sàng đâu. Cứ để cô ấy an tâm thêm một thời gian nữa.}''

''\hl{Sớm muộn gì cũng phải nói thôi.}''

''\hl{Em biết chứ.}''

Tôi im lặng. Cơn mưa rả rích rơi trên ô, từng giọt nặng trĩu. Trong lòng tôi chỉ còn lại một mối bất an lơ lửng, chưa có lời giải đáp.

\begin{center}
  \uline{Ngày hôm sau.}
\end{center}

Sau chuỗi ngày mưa kéo dài, bầu trời hôm nay quang đãng hơn. Không hẳn là có nắng, nhưng ánh sáng chan hòa dịu nhẹ sau lớp mây mỏng cũng khiến không khí trong lành hơn hẳn, trái ngược với bầu không khí ảm đạm những ngày đầu chúng tôi bị kéo về đây.

Sau nhiều ngày mưa dầm dề, bầu trời thị trấn Aogami cuối cùng cũng quang đãng trở lại. Ánh nắng dịu nhẹ len lỏi qua những tán cây, tạo nên những vệt sáng lung linh trên những con đường lát đá còn đọng lại vài vũng nước nhỏ. Không khí trong lành và mát mẻ sau cơn mưa khiến cả thị trấn như được tắm gội, sạch sẽ và tươi mới.

Những dãy nhà gỗ hai bên đường hiện lên rõ nét dưới ánh nắng ban mai. Mái ngói xám đen lấp lánh những giọt sương còn sót lại, trong khi những chậu hoa treo trước hiên nhà đung đưa nhẹ trong làn gió, tô điểm thêm sắc màu cho khung cảnh yên bình. Đâu đó vang lên tiếng chim hót véo von, hòa quyện với âm thanh của những cửa hiệu bắt đầu mở cửa đón ngày mới.

Trên những con phố nhỏ, người dân địa phương bắt đầu xuất hiện đông đúc hơn. Không khí rộn ràng nhưng không ồn ào, tất cả đều mang một nhịp điệu nhẹ nhàng, thanh bình đặc trưng của một thị trấn nhỏ vùng nông thôn Nhật Bản.

Reiko-san khẽ dang tay hít một hơi thật sâu, gương mặt cô thư giãn thấy rõ: ''Hôm nay trời quang đãng ghê!''

Saikyou cũng thở phào, mái tóc đuôi ngựa của cô em đung đưa theo gió. Tôi nhìn cả hai cô gái, làm tôi cũng thấy nhẹ lòng theo. Dẫu vậy,tôi vẫn giữ thói quen đề phòng. Chiếc ô vẫn nằm gọn trong tay tôi và em gái, như một phản xạ chẳng thể bỏ được sau bao ngày mưa dầm.

Trên đường tới trường, chúng tôi bắt gặp Shino-san.

''Chào buổi sáng, Shino-san!'' Reiko-san vẫy tay rạng rỡ.

''A, chào mọi người!'' Shino-san đáp lại với nụ cười tươi tắn. ''Hôm nay thời tiết đẹp quá nhỉ?''

''Ừm, không còn mưa nữa rồi,'' Saikyou-san gật đầu. ''Shino-san trông có vẻ vui hơn mọi ngày đấy.''

''Thật sao?'' cô bạn khẽ đỏ mặt. ''Có lẽ vì\dots\ mình đã có thêm những người bạn tốt.''

''Cậu nói gì vậy,'' Reiko-san cười nhẹ, vỗ vai Shino-san. ''Chúng ta là bạn bè mà, đương nhiên phải vui vẻ cùng nhau rồi.''

''Nào, chúng ta nên đi thôi,'' tôi nhắc nhở. ''Sắp muộn giờ vào học rồi đấy.''

''Ừ nhỉ, đi thôi!'' Shino-san hớn hở nói.

Hôm nay, tâm trạng cô ấy khác hẳn: khuôn mặt tươi tắn, đôi mắt trong veo sáng rỡ hơn mọi ngày. Có lẽ nhờ thời tiết dịu lại, hoặc cũng có thể vì từ khi kết thân với Reiko-san nên cô ấy đã mở lòng hơn chăng. Nhìn cách Shino-san vui vẻ chào hỏi, tôi chợt thấy hình ảnh ấy thật bình yên, quá bình yên so với những thứ đã xảy ra một cách rời rạc trong đầu mình.

Liệu đây có phải là một sự yên bình trước một cơn bão sắp tới\dots? Tôi không quá bi quan, nhưng trong lòng tôi lại đang chực chào một thứ cảm giác như vậy.

Một thứ gì đó\dots\ giống như trong trận động đất định mệnh ấy.

\ct

Vào đến lớp, bầu không khí cũng tươi tỉnh hơn. Tiếng trò chuyện rộn ràng, chẳng còn cái u ám nặng nề như mấy hôm trước. Tôi dừng bước trước cửa lớp, quay sang nói với cả nhóm: ''Ba người vào lớp trước đi. Tôi cần ra ngoài có chút việc.''

''Sao thế?'' Shino-san nghiêng đầu. ''Cậu quên gì à?''

''À không, chỉ là\dots\ tôi định đi ra ngoài hóng gió một chút,'' tôi liền bịa ra một lý do, hy vọng không tiết lộ cho cô ấy.

''À\dots\ ra vậy,'' cô nàng tóc nâu bật cười nhẹ. Nhìn gương mặt rạng rỡ có chút đỏ hỏn vì hơi ngượng, tôi không khỏi nổi da gà.

Saikyou gật đầu, không hỏi thêm. Có lẽ em gái tôi đã đoán được tôi định làm gì. Những gì mà chúng tôi đã thống nhất, tất cả đều đã vạch ra từ đêm hôm qua rồi.

Tôi đã mong rằng hai người còn lại chỉ nghĩ rằng đó chỉ là một chuyện vụn vặt không đáng bàn. Vấn đề này\dots\ tôi không muốn ai can dự nào, cũng không muốn kéo bất cứ ai vào đây.

Nhưng không.

Reiko-san quay sang hai người kia và bảo: ''Hai cậu vào lớp trước đi. Tôi có chút chuyện muốn hỏi Kyuen-san.''

Hai người kia gật đầu, không tỏ vẻ nghi ngờ gì, rồi bước vào lớp, hòa chung với bầu không khí náo nhiệt bên trong.

Khi chỉ còn lại hai đứa, cô gái có hai đôi mắt đỏ và xanh kia nhìn thẳng vào mắt tôi: ''Cậu định tìm lại cái lọ hoa bỉ ngạn phải không?''

Tôi sững người. Rõ ràng tối qua, tôi và Saikyou đã đợi đến tận khuya, khi tiếng thở đều đặn của cô ấy vang lên mới dám thì thầm bàn bạc với nhau. Chúng tôi đã rất cẩn thận, thậm chí còn cố tình nói thật khẽ để không đánh thức cô ấy. Vậy mà\dots\ làm sao cô ấy lại biết được?

Cảm giác như bị đóng băng, tôi không thể cử động nổi. Mọi tính toán cẩn thận của tôi đều đổ sông đổ bể. Tôi không muốn kéo thêm ai vào chuyện này, nhất là Reiko-san - người đã cho chúng tôi một nơi nương tựa.

''\dots Sao cậu biết?''

''Tối qua\dots'' cô ngập ngừng một chút rồi nói tiếp, ''tôi\dots\ đã nghe được cuộc nói chuyện của cậu và Saikyou-san.''

Tôi thoáng sững lại. Đây vốn là chuyện riêng của anh em tôi, và tôi không muốn cô ấy bị kéo vào.

Làm sao có thể được chứ?

\pov{Haragen Reiko}

Tôi biết rõ việc nghe lén người khác là không nên. Nếu là chuyện gì bình thường, tôi chắc chắn đã quay đi ngay từ đầu. Nhưng vì cái tên Shino-san vang lên trong đoạn đối thoại ấy, tôi không thể làm ngơ chưa nghe thấy gì. Cô ấy là bạn thân nhất của tôi - người duy nhất mà tôi cảm thấy cần được bảo vệ. Nếu có điều gì liên quan tới cô ấy, tôi buộc phải nghe cho ra lẽ.

\begin{center}
  \uline{Hồi tưởng\\Đêm hôm qua.}
\end{center}

Đó là một buổi đêm nọ, lúc mà trời mưa tầm tã bên ngoài, cũng không khác gì những đêm trước đó.

''Ư\dots\ Buồn tè quá\dots''

Tôi lồm cồm ngồi dậy, cảm thấy hơi choáng váng vì giấc ngủ bị cắt ngang. Bóng tối dày đặc bao trùm căn phòng, chỉ có tiếng mưa rơi rả rích ngoài hiên làm nền. Tôi dò dẫm bước xuống giường, chân trần chạm vào sàn gỗ lạnh ngắt.

Tay tôi mò mẫm trong bóng tối, tìm đến chiếc bàn học gỗ quen thuộc. Chiếc đèn dầu vẫn nằm yên vị ở góc bàn từ tối qua. Tôi lần tìm hộp diêm, rút ra một que, quẹt nhẹ. Ánh lửa bùng lên trong khoảnh khắc, rồi tôi vội vã đưa nó đến tim đèn. Ngọn lửa nhỏ nhảy múa, tỏa ra thứ ánh sáng vàng ấm áp, đủ để soi đường trong đêm tối.

Tôi cầm đèn dầu, rón rén bước ra hành lang. Những tấm ván gỗ cũ kỹ kêu cót két dưới bàn chân trần của tôi. Ánh đèn dầu hắt lên những bức tường, tạo thành những cái bóng nhảy múa kỳ quái. Nhà vệ sinh nằm ở cuối hành lang, và tôi phải đi ngang qua phòng dành cho khách - nơi hai anh em sinh đôi Kyuen-san và Saikyou-san đang ở.

Nhưng rồi, khi đi ngang qua căn phòng ấy, tôi nhận ra một vệt sáng mờ hắt ra từ khe cửa. ''\textit{Sao giờ này họ vẫn còn thức?}'' Tôi tự hỏi, bất giác dừng bước. Tò mò lấn át cả cơn buồn tiểu, tôi nhẹ nhàng đặt đèn dầu xuống sàn, rồi áp tai vào khe cửa. Tiếng thì thầm của hai anh em vọng ra, đủ rõ để tôi có thể nghe thấy.

''\dots bình hoa bỉ ngạn đỏ sáng nay\dots\ không bình thường.''

''\dots giấc mơ hôm đó, anh vẫn nhớ chứ, Nii-chan? Bông hoa ấy cũng xuất hiện\dots''

''\dots Nanase-san, thái độ của cậu ấy. Khi nói chuyện thì khác, nhưng hễ nhắc đến Reiko-san, mà nhất là từ Shino-san\dots\ ánh mắt lại thay đổi.''

''\dots nghĩa là Shino-san có thể đang bị bắt nạt\dots?''

Tôi nín thở. Mỗi câu từ họ nói như lưỡi dao xoáy vào tim mình. Tôi hiểu Shino-san hơn bất kỳ ai: cô ấy nhút nhát, hay tự ti, chẳng bao giờ dám to tiếng với ai. Và giờ, chính hai anh em ấy lại nghi ngờ rằng Shino-san đang bị bắt nạt\dots\ hoặc thậm chí có người muốn hãm hại cô ấy.

Tôi muốn nghe tiếp, muốn biết rõ hơn những gì hai người họ đang bàn. Nhưng ngay lúc ấy, cơn mót tiểu ập tới, không cho tôi kiềm chế thêm. Tôi cắn môi, nuối tiếc bỏ đi, chạy vội vào nhà vệ sinh. Lúc tôi quay lại thì căn phòng đã tối om, không còn một tiếng thì thầm nào nữa.

Cả đêm ấy tôi không sao chợp mắt được. Những lời họ nói cứ vang vọng trong đầu tôi, trộn lẫn với lo lắng cho cô bạn nhút nhát của tôi. Và đến sáng nay, tôi đã quyết: phải hỏi thẳng họ cho ra nhẽ.

\begin{center}
  \uline{Kết thúc hồi tưởng.}
\end{center}

Giờ đây, khi Kyuen-san tìm lại lọ hoa ấy và bảo tôi nên quay vào lớp, tôi hiểu ngay cậu muốn đẩy tôi ra khỏi chuyện này. Nhưng càng bị đẩy ra, tôi càng muốn kéo vào. Có lẽ đó là bản tính cố chấp của tôi, hoặc có thể là trực giác mách bảo rằng tôi không thể bỏ qua chuyện này.

Tôi ngẩng mặt nhìn thẳng vào mắt cậu ấy, giọng kiên quyết: ''Cái kim trong bọc lâu ngày cũng phải lòi ra thôi. Nếu cậu không nói, sớm muộn gì tôi cũng sẽ tự tìm ra. Tôi không phải là người dễ dàng từ bỏ đâu, Kyuen-san.''

Cậu con trai thoáng khựng lại. Trong ánh mắt cậu hiện lên sự lưỡng lự, như thể đang cân nhắc có nên để tôi biết sự thật hay không. Cậu chắc chắn sợ rằng tôi sẽ sốc, hoặc không chịu nổi. Tôi biết chứ.

Hai anh em nhà Haragen ấy\dots\ họ đang giấu giếm gì từ tôi? Nó bây giờ không còn là sự tò mò giản đơn nữa, mà có lẽ\dots\ mọi chuyện còn phức tạp hơn như thế.

Tôi hít một hơi thật sâu, rồi nói tiếp với giọng chắc nịch: ''Dù là chuyện gì đi nữa, tôi sẽ chịu được. Miễn là nó liên quan đến Shino-san, tôi không thể làm ngơ. Cậu không thể ngăn cản tôi đâu.''

Một thoáng im lặng bao trùm. Kyuen-san khẽ gật đầu, thở dài nhún nhường nhưng vẫn ánh lên một sự cẩn trọng: ''Được rồi. Nhưng có một điều kiện: tuyệt đối không được để Shino-san biết. Cho dù là nửa chữ.''

Tôi gật đầu, đáp không chút do dự. Trong lòng tôi dấy lên một cảm giác kỳ lạ - nửa như hồi hộp, nửa như lo lắng, nhưng cũng có cả sự quyết tâm mãnh liệt. Tôi muốn tin rằng, lần này, sự cứng đầu của mình là đúng.

Thế là, thay vì tách nhau ra, tôi và cậu con trai ấy cùng bước đi, rời khỏi lớp, tìm lại chiếc lọ hoa bỉ ngạn ấy - thứ có lẽ sẽ mở ra những bí mật mà chúng tôi chưa từng tưởng tượng đến. Mỗi bước chân như càng khẳng định thêm quyết tâm của tôi: dù có chuyện gì xảy ra, tôi cũng sẽ không lùi bước.

Nhưng điều khiến tôi day dứt hơn cả\dots\ chính là khi nghe giọng nói trầm tĩnh của hai anh em trong căn phòng tối hôm đó, tôi lại thấy trong lòng mình nhói lên một cảm giác rất lạ.

Giống như\dots\ tôi đã từng trải qua cảnh này ở đâu đó.

Giống như\dots\ họ không phải là ''người ngoài'' đối với tôi.

Khi nghe thấy giọng hai anh em đêm hôm ấy trao đổi, tôi lại thấy trong lòng mình gợn lên một ký ức rất mơ hồ. Ký ức của một lần gặp gỡ nào đó\dots\ nơi mà một người con trai đã dẫn đường cho tôi, để rồi tôi được đoàn tụ với một người thân. Cảnh ấy mờ nhòe như giấc mơ, nhưng để lại dư âm rõ ràng đến mức trái tim tôi vẫn nhớ.

Để rồi, khi cùng Kyuen-san bước đi dọc hành lang, tôi cảm thấy tim mình dường như đang hòa cùng nhịp với họ, như thể ba chúng tôi đã cùng nhau bước đi từ rất lâu rồi. Càng gần bên họ, tôi càng thấy như mình vốn thuộc về nơi này. Giống như đã từng là một phần của gia đình họ, dù lý trí tôi không thể nào nhớ lại nổi.

Thậm chí, khi Saikyou thoáng liếc sang tôi bằng ánh mắt yên lặng, tôi có linh cảm cô ấy đã biết trước rằng tôi sẽ không chịu đứng ngoài. Và đúng vậy, tôi không thể.

Cứ cảm tưởng rằng\dots\ vấn đề của Kyuen-san và Saikyou-san cũng là vấn đề của tôi vậy.

\ct

Tôi dẫn cậu con trai đến căn phòng bỏ trống cạnh dãy hành lang, nơi tôi đã giấu lọ hoa bỉ ngạn đỏ sáng hôm ấy. Trên chiếc bàn gỗ bụi phủ, lọ hoa đỏ rực vẫn nằm đó, những cánh hoa lặng lẽ xòe ra, đẹp đến mức nếu ai đó hiểu được bản chất của bông hoa này, họ sẽ rợn người tới già.

''Cậu muốn xem lại nó sao?''''tôi lên tiếng.

Kyuen-san gật nhẹ, ánh mắt dán chặt vào lọ hoa như bị thôi miên. Tôi không kìm được, buộc phải hỏi: ''Tại sao loài hoa này lại khiến cậu mê hoặc đến vậy?''

Cậu im lặng một lát, rồi trả lời bằng giọng đều đều: ''Bởi vì tôi đã thấy nó\dots\ trong giấc mơ. Một lần, ở một lớp học trống, có lọ hoa đỏ đặt ngay ngắn trên bàn. Cảnh tượng ấy lặp đi lặp lại, giống hệt như sáng hôm qua khi chúng ta thấy nó trên bàn của Shino-san.''

''Giấc mơ?''

''Kể ra thì dài lắm. Nhưng mà nó cũng không dễ chịu gì cho cam cho quýt đâu.''

Tôi nuốt khan, không biết nên phản ứng ra sao.

''Cậu có biết loài hoa này mang ý nghĩa gì không?'' cậu bất chợt quay sang hỏi tôi.

Tôi rùng mình, ký ức về buổi sáng hôm ấy lại ùa về.

''Tôi biết chứ. Hoa bỉ ngạn đỏ tượng trưng cho chia ly, và\dots\ cho cái chết. Thế mà nó lại bị đặt trên bàn của một bạn học vẫn còn sống. Ai đó đã làm vậy, chẳng khác nào gửi đi một lời nguyền tàn độc\dots''

Tôi ngập ngừng, nhưng rồi lấy hết dũng khí để nói ra: ''Có khi nào\dots\ tất cả những điều này liên quan đến Shino-san?''

Kyuen-san nhìn thẳng vào tôi. Đôi mắt cậu ánh lên một tia lạnh lẽo. ''\dots Có thể. Và tôi cũng lờ mờ đoán được người đã làm việc đó. Nhưng vẫn còn quá sớm để kết luận. Tôi cần thêm chứng cứ.''

Một thoáng rùng mình chạy dọc sống lưng tôi. Tôi cảm nhận rõ rệt: càng đi sâu vào, tôi càng khó lòng quay lại. Chẳng khác nào bị loài hoa chết chóc ấy quyến rũ, nhưng lại là một sự quyến rũ đầy u tối.

\pov{Haragen Kyuen}

Tiếng chuông báo hiệu giờ học vang lên, cắt ngang dòng suy nghĩ của tôi. Tôi bảo Reiko-san trở lại lớp trước. Cô do dự vài giây, nhưng rồi gật đầu.

''Cậu đừng giữ tất cả cho riêng mình. Cứ bảo tôi nếu tìm được chi tiết nào đó khác thường'' cô buông lại một câu trước khi rời đi. Tôi chỉ đáp bằng cái gật khẽ, rồi quay lại đối diện với lọ hoa.

Một mình trong căn phòng tĩnh lặng, tôi nhìn chằm chằm vào sắc đỏ đến chói mắt của những cánh hoa. Hình ảnh Shino-san hiện về trong đầu: nét mặt bình thản ở hành lang khi nhìn thấy bông hoa này, như thể cô đã nhìn thấy nó không biết bao nhiêu lần.

Và cả ánh mắt Nanase-san, cái liếc thoáng qua đầy khó hiểu. Tất cả như những mảnh ghép chưa hoàn chỉnh, nhưng tôi biết chúng liên quan đến nhau.

Tôi cúi xuống, nhấc bông hoa ra khỏi lọ. Khi ấy, một vật gì đó lấp lánh ánh đen lóe lên bên trong. Tôi nheo mắt, rồi xóc nhẹ chiếc lọ. Một vật cứng rơi xuống bàn.

*CẠCH*

Tôi chết lặng. Đó là một sợi dây cáp đen, dài chừng ba gang tay. Đầu cắm hiện rõ ràng cổng Type-C.

Tôi nuốt khan, não bộ quay cuồng. Năm nay mới là 1946. Điện thoại di động đầu tiên ở Nhật còn phải đợi khoảng bốn mươi năm nữa, còn smartphone\dots\ chắc phải đến hàng chục năm sau nữa mới có. Thứ này không thể nào tồn tại ở đây.

''\dots Có thể\dots\ nó liên quan đến chiếc đồng hồ đeo tay?'' tôi lẩm bẩm, bàn tay vô thức siết chặt sợi dây cáp ấy.

Cả hai anh em tôi đã có mặt trước, mặt sau và hai dây đeo, từ đó lắp được một cái khung hoàn chỉnh của một chiếc đồng hồ đeo tay, nhưng phần bên trong vẫn chưa có. Tôi nghĩ đây chính là thứ được dùng để cấp điện cho nó.

Tôi chưa kịp đào sâu thêm thì tiếng chuông phụ vang lên, nhắc nhở tôi đã quá trễ. Tôi nhét vội sợi cáp vào túi, thu dọn mọi thứ rồi nhanh chóng quay trở về lớp.

Nhưng tâm tri tôi thì không thể yên. Lọ hoa, Shino, Nanase, và giờ đây là cả một vật thể từ tương lai, tất cả đang kéo tôi vào một mê cung đầy u ám, nơi câu trả lời chỉ khiến tôi đặt ra thêm nhiều câu hỏi.

Tạm thời, tôi phải để lại những nghi vấn ấy trong bóng tối.

\ct

Giờ nghỉ trưa đã tới.

Tôi thu dọn sách vở, chuẩn bị cho tiết học tiếp theo, nhưng tâm trí thì không tài nào yên ổn nổi. Lọ hoa bỉ ngạn sáng nay\dots\ không, chính xác hơn là thứ tôi đã lôi ra từ bên trong nó: sợi cáp đen bóng với đầu cắm Type-C. Thứ không thể tồn tại ở thập niên 40 này.

Càng nghĩ tôi càng thấy nghèn nghẹn trong ngực. Nếu nó liên quan đến Shino-san, thì quá bất hợp lý. Cô ấy không thể nào nắm giữ một thứ vượt ngoài thời đại như thế, và phiên bản hồn ma của cô cũng không dính dáng gì tới nó. Nhưng nếu không phải cho cô ấy\dots\ thì sao? Ý nghĩ lóe lên khiến sống lưng tôi lạnh đi. Có lẽ, thứ đó là để lại cho tôi và Saikyou. Một bí mật nào đó, được ai đó che giấu dưới lớp thời gian, mà giờ mới bắt đầu lộ diện.

Tôi nhíu mày lại, muốn nối kết các mảnh vụn này với nhau, thì một tiếng rên khe khẽ cắt ngang. Tôi ngẩng đầu lên, bắt gặp Shino-san đang cúi rạp xuống bàn học của mình, bàn tay run run lục lọi từng ngăn kéo nhỏ. Khuôn mặt cô nhăn nhó, xen lẫn hoảng loạn.

''Cậu\dots\ cậu đang tìm cái gì vậy?'' tôi đứng dậy, nhẹ nhàng bước tới.

Cô nàng giật mình ngước lên, mái tóc nâu của cô khẽ xõa xuống hai bên: ''Hộp bút của mình. Mình để trong ngăn bàn từ sáng, giờ không thấy đâu nữa. Cậu có thể\dots\ giúp mình tìm không?''

Một sự chần chừ hiện rõ trong giọng của cô ấy. Phải rồi, đúng là cô ấy đã kết thân với rất nhiều bạn của lớp, nhưng dường như chỉ có Reiko-san là cô tích cực nói chuyện nhất. Tôi và Saikyou là thứ hai, nhưng cũng không nhiều hơn các bạn khác là bao.

Tôi gật đầu, đảo mắt quanh lớp. Cô bạn thân nhất của Shino-san lại đang mải giải đáp bài tập từ các bạn khác, nên tôi không nghĩ cô ấy có thể giúp được.

Đúng lúc này, cánh cửa gỗ kêu cót két mở ra. Shiina-san bước vào. Cô liếc qua trong chốc lát, rồi quay cẳng tay vào phòng sau như định đi tiếp. Tôi gọi nhỏ: ''Shiina-san, chờ chút! Shino-san đang tìm hộp bút. Cậu giúp bọn này tìm nhé?''

Cô bạn tóc xám dừng lại, rướn mày nhìn tôi, rồi chậm rãi bước lại gần. ''Sao? Có việc gì?''

Shino-san lục tung ngăn bàn một lần nữa, giọng lắp bắp đầy lo lắng: ''Mình chắc chắn để trong ngăn bàn\dots\ nhưng bây giờ trống trơn rồi.''

Cô bạn bờm Sư tử đứng bên cạnh, mắt lướt nhanh qua mặt bàn đang bừa bộn những sách vở. ''Có khi cậu để quên ở phòng học trước chăng?''

''Tôi vừa thấy cậu ấy lấy bút từ hộp lúc tiết Toán mà,'' tôi lắc đầu, nhíu mày. ''Khó có khả năng bỏ quên lắm.''

Cô nàng tóc nâu siết chặt tay áo, giọng nhỏ dần như thể sợ chính suy nghĩ của mình: ''Hay là ai đó\dots\ lấy mất rồi?''

Shiina-san khoanh tay trước ngực, thở dài một tiếng nặng nề. Cô đảo mắt nhìn quanh lớp học vắng vẻ: ''Vậy thì chúng ta phải lục kỹ hơn. Tôi sẽ kiểm tra dãy bàn cuối.''

''Được,'' tôi gật đầu, quay sang chủ nhân của chiếc hộp bút ấy đang đứng run rẩy. ''Shino-san, cậu soát lại bàn mình lần nữa, để tôi xem quanh dãy cửa sổ.''

''Đ-Được rồi!''

Chúng tôi chia nhau đi tìm, không gian lớp học bỗng chốc nặng nề hẳn.

Trong lúc lục tìm, tôi không khỏi nghĩ ngợi. Cái cách chiếc hộp bút ''biến mất'' một cách đột ngột như thế này\dots\ quá đỗi kỳ lạ. Tôi liếc nhìn về phía bàn của Nanase-san - cô ta vừa rời khỏi lớp cách đây không lâu, với nụ cười khó hiểu trên môi. Trực giác mách bảo tôi rằng đây không phải là một sự trùng hợp ngẫu nhiên. Có ai đó đang cố tình chơi xỏ Shino-san, và tôi gần như chắc chắn về người đứng sau chuyện này.

\pov{Haragen Saikyou}

''Hầy\dots\ Quả thật Nii-chan vẫn luôn tốt bụng như thế nhỉ.''

Tôi liếc qua bên kia lớp, thấy Nii-chan cùng Shiina-san đang lom khom dưới dãy bàn, tìm hộp bút cho Shino-san. Ngay cả cô nàng tốc nâu cũng đang đỏ mặt, vội vàng xoay sang chỗ nọ chỗ kia như thể trong đó có cất giấu báu vật.

Tôi nhớ rõ, ở tiết trước hộp bút vẫn còn trên bàn cô ấy. Chính mắt tôi đã thấy Shino-san mở nó để rút cây bút chì. Thế mà khi tôi cúi đầu viết bài, chưa đến năm phút ngẩng lên thì hộp bút đã biến mất. Không một dấu vết. Chỉ tới khi chuông báo hết giờ, cô mới bàng hoàng nhận ra. Thứ đó biến mất theo cách chẳng hề tự nhiên chút nào.

Tôi thở dài, quay sang Nanase-san, người vừa mới bước vào. Cô nàng đứng khoanh tay, chẳng buồn nhúc nhích, ánh mắt đầy vẻ lạnh nhạt. ''Này Nanase-san, cậu có biết gì về cái hộp bút của Shino-san không?''

Cô ta chỉ khẽ liếc tôi, giọng đều đều: ''Tôi không để tâm đến mấy chuyện lặt vặt ấy.''

Tôi nhướng mày, cố giấu nụ cười mỉa. ''Vậy\dots\ cậu không định giúp cô ấy sao?''

''Đồ của ai thì người đó tự tìm. Tôi không có thói quen xen vào việc không liên quan.'' Cách cô nói không hẳn gắt, nhưng có một sự chua ngoa thoáng qua. Thái độ ấy chỉ khiến tôi càng thêm nghi ngờ rằng Nanase-san biết nhiều hơn những gì cô chịu nói.

Cũng không thể loại trừ khả năng\dots\ chính cô ta là người giấu chiếc hộp bút đó. Những gì liên quan tới Shino-san\dots\ quanh đi quẩn lại chỉ có thể là Reiko-san, hoặc Nanase-san. Cô bạn mắt hai màu sẽ không đời nào làm thế vì cô ấy ngồi cách rất xa, nhưng cô gái tóc xanh kia lại ngồi ngay cạnh bàn của cô ấy, nên càng có khả năng này hơn.

Tôi đang định buông lời chọc ngoáy thì một giọng lanh lảnh vang lên ngoài cửa. ''\textbf{Này, Nanase-chan ơi\~{}! Tớ có vài thứ muốn hỏi cậu đây\~{}\~{}!!}''

Honoka-san nhảy chân sáo bước vào, tay vẫy phấp phới như chào khán giả. Mắt cô sáng lên khi thấy tôi: ''Ồ, Saikyou-chan cũng ở đây à!''

Tôi chống cằm, cười nhạt: ''Cô bạn loa phường này lại có gì để mua vui chúng ta đây mà. Mồm bà này cứ mở ra là lại phun mấy lời ngớ ngẩn thôi.''

Cô bạn tóc vàbg bĩu môi: ''Tớ nghiêm túc mà! Muốn hỏi Nanase-chan một chuyện quan trọng.''

Nanase-san chỉnh lại gọng tay áo, vẫn giữ giọng bình thản: ''Có việc gì cậu muốn hỏi tôi vậy, Honoka-san?''

''Ờm, xem nào\dots'' cô bạn chống tay lên vai, nhắm mắt lại suy nghĩ, rồi bật ra: ''À, phải rồi! Gia đình cậu xây dựng nên ngôi trường Aogami này, đúng chứ?''

Tôi khoanh tay, hất cằm về phía Honoka-san: ''Cậu biết chuyện này từ đâu vậy? Lại bị gọi lên phòng hiệu trưởng vì làm gì đó lố lăng nữa à?''

Cô bạn bờm Cáo lập tức đỏ mặt, lườm tôi: ''Này! Tớ không có làm gì xấu hết! Chỉ là\dots\ họ gọi tớ lên vì chút việc thôi mà!''

''Nghi lắm\~{}''

Cô lảng tránh, rồi quay sang Nanase-san: ''Tóm lại thì, có đúng là gia đình cậu xây trường này không?''

Cô con gái doanh nhân gật đầu nhẹ. ''Đúng là vậy. Ông và bố tôi đã trực tiếp giám sát công trình này. Nhưng sao? Có vấn đề gì chăng?''

''Tớ muốn biết cụ thể hơn! Vùng đất quanh đây trước khi xây trường có `lịch sử' gì không? Như là\dots\ nơi này từng là mồ chôn tập thể, hay có người bỏ mạng ở đây chẳng hạn?''

Tôi thoáng cau mày. Mồ chôn tập thể? Người bỏ mạng? Ý tưởng nghe kỳ quặc, nhưng cũng không phải là không có khả năng. Một thoáng rợn rùng chạy dọc lưng tôi.

Nanase-san nhíu mày, đáp ngay: ''Tôi có thể khẳng định chắc chắn rằng không có chuyện đó. Ngôi trường này hoàn toàn sạch sẽ, không hề dính dáng đến những lời đồn thất thiệt.''

Tôi ngả người vào ghế, khẽ cười: ''Nghe mấy xì xào đó từ đâu vậy, Honoka-san?''

''Ừm\dots\ Mấy bạn lớp bên?''

''Đúng là bà chúa buôn chuyện,'' tôi chép miệng, rồi quay sang Nanase, hiếu kỳ hỏi tiếp. ''Này, nhà cậu xây trường này không sợ ma ám sao?''

Nanase chỉnh lại giọng, có phần trang trọng hơn: ''Tất cả các nghi thức trừ tà cần thiết đã được thực hiện trước khi công trình bắt đầu. Không có gì để lo ngại cả.''

Cô bạn tóc vàng tròn mắt: ''Cẩn thận đến mức thế cơ à?''

Tôi gật gù: ''Có thờ có kiêng, có thiêng có lành. Tôi đồng ý với cách làm đó. Cẩn thận vẫn còn hơn.''

Nanase-san cũng gật đầu: ''Chính xác. Thế nên, tôi khẳng định là ngôi trường này không bị thế lực nào vấy bẩn cả.''

Honoka-san mím môi trước những bằng chứng đanh thép từ chính con gái của vị doanh nhân đã dựng ngôi trường này. Cô quay về bên kia, lẩm bẩm như nói một mình: ''Vậy thì\dots\ lời đồn kia có khi lại đúng\dots''

Tôi nghiêng người, nheo mắt: ''Lời đồn kia?''

Cô nàng tóc xanh cũng chớp mắt: ''Cậu đang nói tới chuyện gì vậy?''

Cô nàng lăng xăng bật dậy, chỉ tay về phía trước như thể sắp công bố bí mật: ''Lời đồn về con ma ở phòng Mỹ thuật đó!''

Cả hai chúng tôi cùng đồng thanh: ''\dots Hả?''

Honoka-san chống hông, đắc ý: ''Tớ muốn xem nó có thật không nên đang đi thu thập thông tin. Mấy đứa lớp bên cũng nhắc đến nữa đấy!''

Nanase-san khoanh tay, thở ra một tiếng dài chán chường, như thể cô đã nghe những lời như vậy quá nhiều lần tới mức phát chán: ''Cậu mà cứ chạy theo mấy tin đồn vô căn cứ thì chẳng được ích lợi gì đâu.''

''Ể? Nhưng họ cũng kể rất chi tiết đấy, nên nghe cũng đáng tin hơn chứ nhỉ?''

Cô ấy bĩu môi nhẹ. ''Tôi cho rằng điều đó hoàn toàn không cần thiết. Những gì cô ấy nói\dots\ chỉ là chuyện hoang đường.''

''Ơ kìa Nanase-san, đừng nghiêm túc thế chứ,'' tôi cười nhẹ. ''Đôi khi những câu chuyện như vậy cũng thú vị mà.''

Honoka-san lại khoanh tay, im lặng vài giây, rồi bỗng lóe sáng ý tưởng. ''Vậy thì, sao không đích thân đi kiểm chứng? Một bài kiểm tra gan dạ chẳng hạn!''

Nanase-san nhíu mày, giọng nghiêm hẳn: ''Kiểm chứng thì kiểm chứng. Phòng Mỹ thuật đúng không? Vậy bây giờ để tôi--''

''Không không!'' cô bạn tóc vàng vội cắt ngang. ''Đi vào ban ngày thì có giá trị gì chứ? Người ta đồn là ban đêm mới xuất hiện kìa!''

''Ban đêm á?'' tôi nuốt khan, bỗng cảm thấy lạnh sống lưng.

Cô nàng tóc xanh nhìn cô bạn, ánh mắt dần hiểu ra ý định. ''Này\dots\ cậu không định\dots''

Honoka-san bật cười, giơ tay làm dấu thắng lợi: ''Đúng như cậu nghĩ đó! Bài kiểm tra gan dạ! Tớ sẽ báo cho cậu khi nào. Đừng có bám ti mẹ mà trốn nhé!''

''Cậu\dots!'' cô nàng gằn giọng, lao tới.

''\textbf{Á ha ha ha!}'' cô nàng bờm Cáo sa mạc vừa cười vừa chạy biến ra hành lang. Nanase-san nghiến răng đuổi theo.

Tôi nhìn cô ấy lao ra cửa, gương mặt đỏ bừng vì giận dữ, còn cô bạn nhí nhảnh thì vừa chạy vừa cười khanh khách, như thể cả ngôi trường này chỉ là sân chơi của cô ấy. Tiếng chân họ rầm rập khắp hành lang, xen lẫn tiếng ''Đứng lại đó cho tôi!'' và ''Bắt được thì bắt đi\~{}!'' vang vọng như một vở hài kịch bất đắc dĩ, nhưng cũng đủ để xua đi bầu không khí nặng nề trước đó.

Tôi khẽ lắc đầu, bật cười. Đúng là chẳng ai có thể đoán được cô bạn lắm trò sẽ lôi kéo một người quyền quý và nghiêm túc vào mấy trò trẻ con này.

Thế nhưng\dots\ khi nghe đến hai chữ ''thử thách gan dạ'', tôi bỗng khựng lại. Cái ý tưởng ngớ ngẩn ấy, không hiểu sao, lại khiến trái tim tôi rung lên một hồi chuông cảnh báo. Trong giấc mơ kia, tôi và Nii-chan từng bước vào phòng Mỹ thuật. Ở đó, tôi đã cảm nhận rõ ràng thứ gì đó bất thường, rùng rợn tới mức chỉ nhớ lại thôi cũng khiến da gà nổi khắp cánh tay.

Có lẽ\dots\ trò ''gan dạ'' mà Honoka-san vừa nghĩ ra sẽ không chỉ là trò đùa. Nó có thể mang tới câu trả lời cho những điều mà tôi và anh trai vẫn chưa thể giải thích.

Tôi tựa lưng vào ghế, dõi mắt nhìn theo hai bóng người vẫn đang rượt đuổi vòng quanh như trẻ con, nhưng trong lòng đã dấy lên một nỗi tò mò xen lẫn bất an.

Nếu cô ấy thực sự quyết định kéo chúng tôi vào cái ''thử thách'' kia\dots\ thì có lẽ tôi cũng sẽ phải tham gia.

\pov{Haragen Reiko}

''Và chỗ này thì cần liên kết lại với ý trên một chút\dots''

''Ồ, vậy là mình hiểu được phần nào rồi đấy! Cảm ơn cậu nhé, Reiko-san!''

''Không có gì đâu. Bạn bè trong lớp cũng phải giúp nhau mà. Mình thấy thế cũng bình thường thôi.''

Tôi vừa chỉ cho người bạn ngồi cạnh cách giải một bài tập khó trong môn sắp tới. Cậu ấy mừng rỡ gật gù cảm ơn rối rít trước khi đứng dậy đi về chỗ, nét mặt rạng ngời như thể vừa gỡ được nút thắt trong lòng. Tôi khẽ thở dài, rồi cũng bật cười nhẹ. Thấy bạn bè bớt căng thẳng vì bài vở, tôi cũng cảm thấy nhẹ nhõm phần nào.

Đưa mắt nhìn ra cửa sổ, tôi bắt gặp cảnh Nanase-san và Honoka-san đang rượt đuổi nhau vòng quanh hành lang, tiếng giày và tiếng cười vang vọng cả một góc trường, trông thật như một màn kịch tấu hài. Một tiếng cười nhẹ thoát ra khỏi môi, xua tan bầu không khí nặng nề thoáng chốc phủ xuống lớp học.

Nhưng rồi, trong khoảnh khắc ngắn ngủi ấy, hình ảnh lọ hoa bỉ ngạn đỏ mà tôi nhìn thấy cùng Kyuen vào buổi sáng lại hiện về. Kèm theo đó là những mẩu hội thoại tôi tình cờ nghe được đêm qua\dots\ tất cả như một chuỗi mảnh ghép rời rạc, ám ảnh tôi không thôi. Chúng gợi ra những dự cảm kỳ lạ, vừa mơ hồ vừa đáng sợ. Tôi thấy bồn chồn, như thể mình chỉ đứng sát mép một màn kịch khổng lồ, chỉ cần một làn gió nhẹ là sẽ bị cuốn thẳng vào trong.

Tại sao một loài hoa gắn với cái chết và sự chia ly lại xuất hiện ở nơi này? Và liệu nó có liên quan gì đến những điều bí ẩn quanh ngôi trường Aogami không? Tôi rùng mình khi nghĩ đến, nhưng lại không dám chắc.

Ngẩng mặt lên, tôi giật bắn mình khi thấy Sakura-san đã đứng ngay trước mặt từ lúc nào, đôi mắt chăm chú nhìn thẳng vào tôi. Cô ấy chăm chú quan sát gương mặt tôi với ánh mắt tò mò lạ lùng, khiến tôi không khỏi lúng túng.

''G-Gì vậy?'' tôi ấp úng hỏi.

''Reiko-san\dots\ Cậu quả thật có một cái gì đó bí ẩn thật đấy,'' cô ấy nói với giọng trầm ngâm.

''Ý cậu là sao chứ, Sakura-san?''

''Hm\~{} Nói sao nhỉ,'' Sakura-san nghiêng đầu, vẫn không rời mắt khỏi tôi. ''Cậu có một sức mạnh huyền bí nào đấy, như thể có thể thu hút mọi người về phía mình vậy.''

Tôi cười gượng, cảm thấy không thoải mái với nhận xét đó: ''M-Mình không nghĩ là mình có thứ sức mạnh ấy đâu.''

''Chả trách mà tại sao các bạn học trong lớp lại bị thu hút bởi cậu đấy?''

''Chắc là do mình có gì đó nổi bật về ngoại hình thôi chứ\dots'' tôi đưa tay lên chạm nhẹ vào khóe mắt, chỉ đến đôi mắt dị biệt một đỏ một xanh của mình.

Nhưng cô ấy lắc đầu, ánh mắt càng trở nên sắc sảo hơn. ''Không phải cái đó. Một cái gì đó\dots\ khác hẳn.''

''Cậu đang nói gì vậy?''

''Nhưng mà trông nó cũng mơ hồ lắm\dots\ Tôi nghĩ nó chính là--''

''Cậu đây rồi, Sakura-san!'' một giọng nói có phần hậm hực cắt ngang câu chuyện. Yae-san đứng chống hông với vẻ không vui, ''Tớ tìm cậu suốt từ nãy tới giờ đấy!''

Sakura-san nhìn cô bạn mới tới với vẻ mặt ngơ ngác. ''\dots Ờm, và cậu là\dots?''

''Hể!?'' cô bạn tóc xám trợn tròn mắt, quay sang tôi với vẻ không tin nổi. ''Reiko-san, cậu vẫn nhận ra tớ mà, đúng chứ?''

''Cậu là Shiromi Yae-san\dots\ đúng không nhỉ?'' tôi đáp, cố giữ giọng bình thản dù trong lòng không khỏi ngạc nhiên trước phản ứng điềm nhiên của Sakura-san.

Yae-san quay phắt lại phía cô bạn tóc hồng, giọng càng thêm bực bội: ''Đến cả cậu ấy còn nhớ mặt tớ kìa! Tớ phải nói với cậu bao nhiêu lần nữa đây, Sakura-san? Cậu thậm chí còn chẳng thèm để ý đến các bạn trong lớp cả!''

Tôi thầm nghĩ trong đầu, không khỏi băn khoăn: ''\textit{Cậu ấy thật sự không để tâm tới mọi người trong lớp sao?}, rồi lên tiếng hỏi: ''Sakura-san, cậu nhớ được đặc điểm gì từ Yae-san không?''

Sakura-san nheo mắt nhìn cô nàng từ đầu đến chân, rồi chợt sáng mắt lên. ''À đúng rồi, người có cái \emph{ấy} to nhất lớp.'' Cô nàng chỉ thẳng vào\dots\ ''bầu sữa'' khiến cho tôi không khỏi bật cười trong bất lực.

''Đó là tất cả những gì cậu nhớ về tớ sao?'' Yae-san thốt lên với vẻ ngạc nhiên lẫn bất bình. ''Cậu biết có nhiều thứ khác để cậu có thể nhận diện tớ mà! Cậu còn biết được Reiko-san đấy thôi!''

''Cậu ta là người duy nhất có mắt hai màu trong lớp này,'' Sakura-san đáp với vẻ thản nhiên. ''Chứ cơ thể của cậu thì lớp mình cũng có vài đứa như thế.''

Tôi không nhịn được một tiếng cười chát. ''Dường như Sakura-san lấy những điểm dễ nhìn nhất để nhận diện các bạn nhỉ?''

Yae-san nghe cô bạn giải thích, chỉ có thể thở dài chán chường, rồi chuyển sang chủ đề khác. ''Thôi, bỏ đi. Mà này, Sakura-san, tiết trước cậu không học thể dục à?''

Sakura-san vòng hai tay ra sau, tỏ vẻ không quan tâm: ''Không. Sao? Việc học môn đó hay không là việc của tôi, liên quan gì đến cậu?''

''Cậu không thể cứ chạy đi đâu đó bất cứ lúc nào có tiết thể dục được!'' cô bạn liền phản đối gay gắt.

''\dots Sakura-san đã trốn môn thể dục nhiều lần trước đó sao?'' tôi hỏi, không giấu được sự ngạc nhiên.

''Nhiều lần là đằng khác!'' cô đáp ngay.

Sakura-san liếc qua phía Kyuen-san đang cặm cụi tìm đồ cho Shino-san, rồi sán tới bàn của tôi. ''Ồ, dù gì thì môn thể dục thể dẹo kia cũng vô nghĩa thôi. Giống như Kyuen-san nói với tôi khi cậu ta học thể dục về hôm trước ấy. Cậu đồng ý với tôi chứ?''

Tôi cười trừ, cố tìm cách trả lời khéo léo. ''Quả thật không phải ai cũng thích học thể dục\dots\ Nhưng mà nó cũng không hẳn là vô dụng đâu.''

''Đúng đó!'' Yae-san nhiệt tình ủng hộ. ''Trước đây, người ta thường coi chuẩn mẫu người phụ nữ là sự khiêm tốn và e lệ, nhưng ở thời này, một người phụ nữ năng động và hoạt bát - như Saikyou-san hay Honoka-san - đang là hình mẫu mới đó.''

''Cậu quả thật là người am hiểu về chuẩn mẫu phụ nữ nhỉ, Yae-san,'' tôi nhận xét. ''Hình như cậu cũng từng nói tương tự như thế với Shiina-san hôm qua, đúng không?''

Sakura-san thở dài, vẻ mặt chán nản. ''\dots Nghe này, tôi thực sự chẳng thèm để tâm tới mấy cái chuẩn chỏe gì cả. Cậu cứ nhai đi nhai lại mấy thứ đó không thấy chán à?''

''Kể cả nếu cậu không quan tâm thì tớ vẫn thấy lo khi cậu cứ biến mất như thế!'' Yae-san nhấn mạnh, rồi hạ giọng xuống. ''Mà, thật đấy, cậu đã đi đâu vào cái giờ ấy vậy?''

Cô bạn Vu nữ ngước lên trần nhà, giọng bình thản đáp ngắn gọn: ''Tới phòng Mỹ thuật.''

''Phòng Mỹ thuật!?'' cô bạn giỏi thể lực trố mắt kinh ngạc.

''Ừ. Ngoài kia nóng tới mức muốn ngất xỉu luôn\~{}'' Sakura-san nói với vẻ thỏa mãn. ''Phòng ấy lại không có người, lại còn gần rừng, nên mới mát thế chứ\~{}''

''Mình nghe đồn là căn phòng ấy cũng bỏ hoang lâu rồi đấy chứ,'' tôi góp lời. ''Sáng nào mình đi qua đó cũng không có một bóng người ở đó luôn.''

''T-Tớ không nghĩ là cậu lại liều đến thế đấy\dots'' Yae-san lắp bắp, vẻ lo lắng hiện rõ trên gương mặt.

''Ý cậu là sao?'' cô bạn tóc hồng nhướn mày. ''Tôi nghỉ ở đó thì sao chứ?''

Yae-san nhìn quanh với vẻ bất an, rồi hạ giọng xuống. ''\dots Cậu đã từng nghe mọi người nói gì về căn phòng Mỹ thuật ấy chứ?''

''Hình như mình cũng từng nghe các bạn ở lớp bên nói rồi,'' tôi tò mò. ''Nó như thế nào vậy?''

''Đây nhé,'' cô bắt đầu với giọng kể chuyện ma. ''Chuyện kể rằng có một linh hồn của một cô gái đã ám ảnh căn phòng ấy\dots''

Hai chúng tôi cùng gật đầu, trong khi cô ấy ghé sát vào tai chúng tôi, chậm rãi kể tiếp: ''Đầu đuôi thế này: có một cô gái là một nghệ sĩ có tài hội họa xuất chúng, nhưng những kỳ vọng của mọi người xung quanh đè nặng lên cô, để rồi cô dần mất đi khả năng vẽ\dots''

''Nghe như thể vì áp lực quá mà mất hứng nhỉ,'' tôi nhận xét, cảm thấy có chút đồng cảm với nhân vật trong câu chuyện.

''Tớ cũng không biết,'' Yae-san nhún vai. ''Tiếp tục nhé: bất cứ khi nào đến trường, cô gái ấy sẽ trốn trong Phòng Mỹ thuật để vẽ tranh một mình.''

''Ừm-hứm,'' Sakura-san khẽ gật đầu, vẻ mặt vẫn bình thản.

\img{e1-6h}

Cô hạ giọng xuống thấp hơn nữa: ''Một đêm nọ, một người bảo vệ nhìn thấy căn phòng ấy vẫn còn sáng đèn. Thấy lạ, anh ta ló mắt vào trong để xem, chỉ để rồi nhìn thấy\dots''

Cả Sakura-san và tôi đều nín thở, chờ đợi phần tiếp theo của câu chuyện.

''\dots một cô gái đã chết với người bê bết máu, cùng với một nụ cười trên môi!''

Một cơn ớn lạnh chạy dọc sống lưng tôi khi hình dung ra cảnh tượng ghê rợn ấy. Trong đầu tôi hiện lên hình ảnh một cô gái nằm bất động, máu loang lổ khắp người, nhưng trên môi vẫn còn vương một nụ cười kỳ quái. Cảm giác sợ hãi và tò mò cùng lúc dâng lên trong lòng.

''T-Tại sao chuyện đó lại xảy ra chứ?'' tôi hỏi, giọng run run.

''Tớ nghe rằng\dots'' Yae-san tiếp tục với giọng trầm xuống, ''cô ấy đã quá tập trung cho bức tranh, muốn tạo ra một kiệt tác mà không ai làm được, đến nỗi dùng luôn máu của chính mình để làm màu vẽ.''

''Ghê thế\dots'' tôi thốt lên, rùng mình. ''Cực đoan đến mức đó luôn\dots''

''Rồi sao nữa?'' Sakura-san hỏi với vẻ bình thản đến kỳ lạ.

''Người ta nói rằng bức tranh cuối cùng của cô ấy trước khi ra đi vẫn được cất giữ ở đâu đó trong phòng Mỹ thuật,'' cô kể tiếp, giọng càng lúc càng nghiêm trọng. ''và những ai \emph{chẳng may hoặc tò mò nhìn bức tranh đó\dots\ sẽ bị nguyền rủa đến chết\dots}''

Yae-san thở dài nặng nề, vẻ lo lắng hiện rõ trên gương mặt. Tôi cũng thở dài, cảm thấy một nỗi buồn xen lẫn sợ sệt kỳ lạ dâng lên trong lòng khi nghĩ về số phận bi thảm của cô gái trong câu chuyện. Trong khi đó, Sakura-san vẫn giữ vẻ mặt bình tĩnh, như thể câu chuyện ma quái vừa rồi chẳng hề tác động đến cô ấy.

''\dots Đó là một câu chuyện hay đấy chứ?'' Sakura-san bình luận với giọng thản nhiên.

''Hể!?'' Yae-san trợn mắt, đập mạnh xuống bàn của tôi. ''Nó chẳng hay một chút nào đâu! Reiko-san thì sao?''

Tôi im lặng, không biết phải đáp lại thế nào. Câu chuyện vẫn còn văng vẳng trong đầu, khiến tôi không khỏi rùng mình mỗi khi nghĩ đến. Ý nghĩ về một bức tranh nhuốm máu được cất giấu đâu đó trong căn phòng tăm tối khiến tôi vừa muốn né tránh, vừa không kìm nổi sự tò mò. Hình ảnh ấy gợi lại những mảnh vụn trong giấc mơ từng ám ảnh tôi, về căn phòng đó, về một áp lực vô hình nặng nề đè lên lồng ngực. Tôi muốn tin chỉ là lời đồn, nhưng trái tim lại mách bảo rằng có điều gì đó thật sự đáng sợ ẩn khuất ở phía sau.

''\dots Reiko-san?'' cô gọi lại, giọng có chút lo lắng.

''À\dots\ xin lỗi. Câu chuyện đó thật sự\dots\ khiến mình lạnh sống lưng. Cứ như có gì đó rất thật vậy,'' tôi nở một nụ cười méo xệch.

Yae-san bĩu môi, thở dài một lần nữa, tưởng rằng tôi không tin vào những gì cô ấy nói, còn Sakura-san thì nhún vai, nhưng một nụ cười thích thú nhếch lên bên khóe môi. Cả hai cùng quay về chỗ chuẩn bị cho tiết học tiếp theo. Còn tôi, tay vẫn vô thức siết nhẹ mép bàn. Những chi tiết trong câu chuyện được đồn ở đây\dots\ dù không muốn tin là thật, nhưng chúng lại quá thực đến mức khiến tôi cảm thấy lạnh sống lưng.

Câu chuyện vừa kể vẫn luẩn quẩn trong đầu như một vệt khói ám. Câu chuyện về căn phòng Mỹ thuật\dots\ không thể chỉ bỏ qua như một tin đồn được.

Nhưng điều làm tôi bận tâm hơn cả\dots\ lại là lời từ Sakura-san từ trước đó. Cậu ấy nói rằng tôi có một thứ ''huyền bí'' nào đó. Ban đầu, tôi nghĩ rằng đó chỉ là trò đùa liên quan đến ngoại hình, với đôi mắt lệch màu vốn đã khiến tôi quen với việc bị người khác chú ý. Nhưng không. Khi nhớ lại cách cô ấy nói, tôi nhận ra ý cậu ấy hoàn toàn không nhắc đến đôi mắt.

Thứ mà cậu ấy ám chỉ\dots\ chính là cái họ của tôi: Haragen.

Cùng họ với Kyuen-san và Saikyou-san.

Một cảm giác mơ hồ trào lên. Thật sự, tôi chưa bao giờ nghĩ nhiều về chuyện này. Nhưng nếu Sakura-san nhận ra và gắn nó với một ''sức mạnh bí ẩn'', vậy nghĩa là\dots\ có một điều gì đó ẩn giấu phía sau cái họ ấy sao? Một điều mà ngay cả bản thân tôi cũng chưa từng được biết đến?

Nỗi bất an dấy lên, xen lẫn một chút tò mò không thể kìm nén. Có lẽ tôi nên kể lại chuyện này cho Kyuen-san nghe\dots\ và cả Saikyou-san nữa, nếu cậu ấy đã có linh cảm. Biết đâu chính từ những chi tiết tưởng chừng nhỏ bé này, chúng tôi sẽ tìm thấy một mảnh ghép quan trọng để giải mã mọi bí ẩn xoay quanh ngôi trường Aogami này.

Tôi hít sâu một hơi, cố gắng ổn định nhịp tim. Một dự cảm rất rõ ràng trong lòng tôi: con đường phía trước sẽ chẳng hề dễ dàng.

\pov{Haragen Kyuen}

Tôi và Shiina-san vẫn tiếp tục cuộc ''phiêu lưu'' tìm đồ cho Shiina-san. Chúng tôi đã tìm khắp các bàn học, ở dưới ngăn của chúng, và cả trên bàn giáo viên. Chúng tôi đã lục soát không dưới ba lần một cách cẩn trọng.

Nhưng rồi, thứ mà Shino-san cần tìm vẫn ''không cánh mà bay.'' Cuối cùng, tôi thở dài vì mệt mỏi, ngồi xuống bàn học của tôi, chống tay nhìn cô bạn Sư tử tìm kiếm. Shiina-san đang lom khom tìm ở những góc ít ai để ý. Cuối cùng, cô khựng lại trước mấy ngăn để giày ở cuối lớp, thứ mà chúng tôi không hề để tâm đến.

Bất chợt, cô ấy khẽ gọi: ''Shino, cái này có phải của cậu không?'' Trên tay cô là một chiếc hộp bút trông rất kỳ lạ, có phần cũ kỹ và xước xát, nhưng tên ở trên đó thì không lẫn vào đâu được.

Cô nàng tóc nâu vội bước tới, vừa nhìn thấy thì mắt cô sáng lên. ''Đúng\dots\ đúng rồi! Chính nó\dots!''

Gương mặt ngập tràn nhẹ nhõm, như thể cả tảng đá lớn vừa được nhấc khỏi ngực. Cả hai chúng tôi cũng cùng thở phào; cuối cùng thì cuộc tìm kiếm chẳng phải vô ích.

''\dots A-nô,'' Shino-san khẽ cúi đầu trước mặt cả hai, giọng run run nhưng chân thành, ''cảm ơn hai cậu nhiều lắm, Fukami-san, Kyuen-san. Nếu không có hai cậu, chắc mình chẳng thể tìm thấy chiếc hộp bút ấy đâu.''

Tôi cười nhạt, đưa tay xua nhẹ: ''Bạn bè trong lớp giúp nhau là chuyện bình thường thôi. Không cần phải khệ nệ kính cẩn quá đâu.''

Shiina-san cũng phụ họa, mỉm cười chống hông: ''Thật đấy. Không cần cảm ơn đâu. À, và cứ gọi tôi là Shiina là được rồi, giống như cách cậu gọi Kyuen-san ấy.''

Shino-san thoáng ngập ngừng, mắt liếc về phía tôi: ''\dots Đấy là vì lớp mình có ba người mang họ Haragen. Mình gọi thẳng tên họ như thế vì không muốn làm họ bối rối thôi.''

Cô bạn Sư tử nhíu mày, rồi bật cười khẽ: ''Đừng có lịch sự quá làm gì. Bạn bè với nhau mà cứ gọi họ trịnh trọng thì trông cũ kỹ lắm.''

''\dots Ờm. Vậy thì\dots\ Shiina-san, Kyuen-san,'' cô bạn mím môi, giọng nhỏ nhưng chắc. ''Mình muốn bày tỏ một chút lòng thành với hai cậu.''

''Đã bảo là không cần rồi mà,'' tôi khẽ lắc đầu, rồi liếc vào chiếc hộp bút trên tay Shino. Nó cũ đến mức trông như đã tồn tại hàng chục năm, nước sơn bong tróc lấm tấm. ''Shino-san, cậu quý trọng nó lắm nhỉ?''

Shino-san ôm hộp bút sát ngực, ánh mắt mềm đi. ''Ừ. Đây là kỷ vật bà mình để lại\dots\ trước khi bà mất.''

Shiina-san gật gù: ''Thế cơ à? Vậy thì tìm thấy nó cũng là may mắn đấy.''

''Đúng vậy!'' nàng tóc nâu cười rạng rỡ hiếm hoi. ''Mình\dots\ mình thực sự cảm tạ hai cậu.''

Trong thoáng chốc, tôi nhận ra: đây mới là Misora Shino thật sự. Hiền hòa, e thẹn, có phần yếu đuối nhưng không giấu giếm lòng biết ơn. Khác hẳn với một Misora Shino lạnh lùng đầy gai góc mà tôi từng bắt gặp ở hành lang. Một con người, hai thái cực đối lập.

Shiina-san khẽ ho một tiếng, rồi hạ giọng: ''\dots À, Shino này. Có điều này tôi muốn nói với cậu.''

''Chuyện gì vậy, Shiina-san?''

Ánh mắt của cô tóc xám lướt nhanh về phía Reiko-san đang chuẩn bị cho tiết tiếp theo. Tôi thoáng nhận ra chút do dự trong hơi thở của cô.

''Tốt nhất\dots\ cậu nên tránh xa cô ấy khi ở trong trường.''

''Ể!?'' Shino-san sững sờ, tròn mắt.

Tôi cau mày: ''Ý cậu là sao chứ?''

''Mình cũng thấy\dots\ ngạc nhiên đấy.'' cô run giọng, bàn tay siết chặt hộp bút. ''Bộ mình với Reiko-san có điều gì khiến các cậu không ưa sao?''

Shiina-san thở dài, giọng cứng rắn: ''Tôi làm vậy là để tốt cho cậu. Tôi không rõ Kyuen với Saikyou nghĩ thế nào, nhưng tôi chắc chắn một điều: cậu mà cứ kè kè đi với cô ấy thì chỉ chuốc lấy sự ghen ghét và phiền toái từ mọi người thôi.''

''Ý\dots\ ý cậu là\dots\ mình không được làm bạn với cậu ấy sao?''

''\dots Đúng vậy,'' Shiina-san gật đầu thẳng thừng.

Cô nàng mắt xanh bặm môi, mắt bắt đầu nhòe đi. ''V-Vậy à\dots''

Bầu không khí bỗng chốc nặng nề. Tôi có thể thấy vai Shino-san run lên, như thể lời mà cô bạn học sinh cá biệt vừa buông ra là một nhát dao lạnh lẽo. Tôi muốn xen vào, nhưng\dots\ đây chưa phải lúc. Có gì đó trong giọng của Shiina-san khiến tôi biết cô không nói bừa.

Cô nàng Sư tử liếc sang tôi, ánh mắt như hỏi ý kiến. Tôi khẽ nhún vai, thở dài: ''\hl{\dots nói lời gì để an ủi cô ấy đi.}''

Shiina xoay lại nhìn cô bạn đang rơm rớm nước mắt. ''Cậu biết đấy, Shino\dots''

''\dots Sao?''

''Nếu cậu có điều gì muốn nói, thì cứ tới chỗ bọn tôi.''

Tôi cũng gật theo, cố nghĩ ra một câu gì đó để khiến Shino-san nhẹ lòng hơn: ''Shiina-san trông hổ báo vậy thôi, nhưng thực ra tốt bụng lắm. Chắc hẳn có lý do gì khiến cậu ấy phải nói vậy. Và Shino-san\dots\ nếu có gì muốn trao đổi, cậu cũng có thể tìm tới bọn tôi.''

''\dots Ừ. Kyuen-san, Shiina-san\dots'' Shino cúi đầu thật thấp, nước mắt chực rơi, ''\dots cảm ơn hai cậu.''

Rồi cô lặng lẽ lủi về chỗ ngồi, vừa lau nước mắt, vừa ôm khư khư hộp bút như sợ nó biến mất lần nữa.

Tôi khẽ huých vai Shiina-san, gằn giọng nhỏ: ''\hl{Này. Sao cậu lại nói thế? Shino cuối cùng cũng có một người bạn để dựa vào, mà cậu lại--}''

''\hl{Không phải tôi muốn thế đâu,}'' cô ấy cắt ngang, mắt ánh lên sự căng thẳng. ''\hl{Nhưng có cái gì đó\dots{} buộc tôi phải nói. Như một linh cảm. Tôi không muốn Shino gặp nguy hiểm.}''

''\hl{\dots Nguy hiểm?}'' tôi nhíu mày. ''\hl{Cậu nói thế là sao?}''

Cô không đáp, chỉ lặng lẽ hất cằm về phía cuối lớp. Ở đó, Nanase-san đang cầm hai vai Honoka-san lắc tới lắc lui, sau màn rượt đuổi ồn ào khi nãy.

''\hl{Tôi phải canh chừng cô ta,}'' cô ấy thì thầm. ''\hl{Hôm qua Yae đã cảnh báo tôi. Và\dots{} tôi thấy tận mắt. Chính Nanase đã vứt hộp bút của Shino vào ô giày. Còn cái cặp dựng ở đó che lại, nên nó mới bị bỏ quên.}''

Tôi sững người, tim chợt đập mạnh. Nanase-sam\dots? Câu chuyện này đang rẽ sang một hướng u ám hơn hẳn.

Vậy những gì mà hai anh em tôi đặt giả thuyết trong hai ngày qua\dots\ xem chừng bắt đầu có lý.

Tôi lặng lẽ nhìn Shino-san, cô gái đang ngồi lau nước mắt ở dãy bàn bên kia. Một hình ảnh khác thoáng hiện trong đầu: Shino, nhưng trong bóng tối của hành lang, gương mặt lạnh lẽo và đáng sợ, như một hồn ma vất vưởng.

Phải chăng, chính những khoảnh khắc như thế này, khi cô bị tách biệt, bị tổn thương, bị ép buộc tránh xa người mình muốn tin tưởng, đã gieo nên bóng hình oan hồn ấy?

Ý nghĩ ấy khiến sống lưng tôi lạnh buốt.

\pov{Misora Shino}

Tôi không thể tập trung được vào bài học. Từng con chữ trên bảng đen như nhòe đi, chỉ còn văng vẳng trong đầu giọng nói của Shiina-san khi nãy: ''Tốt nhất cậu nên tránh xa cô ấy ra khi đang ở trong trường.''

Tại sao\dots\ tại sao Shiina-san lại nói như thế? Với tôi, Reiko-san đâu có làm điều gì sai. Từ lúc cô ấy bước vào lớp này, tôi đã cảm nhận được một điều gì đó mà không ai khác có. Một sự mạnh mẽ, nhưng không hề xa cách, một thứ ánh sáng tôi vừa sợ vừa muốn tiến gần.

Tôi nhớ lại buổi chiều trời mưa khi từ ngôi đền trở về. Lúc ấy, khi tôi còn ngập ngừng với lời nguyện cầu cho bản thân mà tôi đã cầu nhiều tới mức tôi không thể nhớ được, Reiko-san đã đi qua, mắt nhìn về phía trước, như thể chẳng có gì có thể khiến cô ấy chùn bước.

Ánh mắt ấy, quyết tâm ấy, là điều mà tôi muốn học hỏi\dots\ điều muốn nắm lấy để bước ra khỏi cái vỏ nhút nhát của tôi. Chính vì vậy, tôi đã tự hứa rằng sẽ ở cạnh cô ấy, để ít nhất tôi cũng có thể thay đổi.

Nhưng bây giờ thì sao?

Lời của Shiina-san như một lưỡi dao sắc, cứa vào trong lòng. Nếu tôi còn tiếp tục gần gũi với Reiko-san, liệu có khiến cô ấy phải chịu điều gì tệ hại hơn không? Liệu tôi có trở thành gánh nặng cho cô ấy không?

''Đến giờ về rồi đấy,'' Reiko-san khẽ gọi. Giọng nói của cô vang lên ngay bên tai, bình thản, quen thuộc\dots\ và ấm áp đến mức tôi muốn lập tức đáp lại, muốn mỉm cười với cô như mọi khi. Nhưng rồi\dots\ tôi lại tự nuốt lời xuống, như thể có một bức tường vô hình ngăn tôi lại. Nếu tôi cứ tiếp tục như cũ, chẳng phải là đang làm trái ý Shiina-san sao?

Tôi chậm rãi lắc đầu, cố nặn ra một nụ cười gượng: ''Ừm\dots\ Mình còn một vài việc phải làm ở trường. Reiko-san cứ về trước với Kyuen-san và Saikyou-san nhé.''

Cô bạn mắt hai màu thoáng cau mày, ánh nhìn nghi hoặc. ''Thật à? Cậu có chắc không?''

''Không sao đâu!'' tôi vội vàng gạt đi, giọng nghe hệt như một lời nài nỉ. ''Mình muốn tự xử lý, nên cậu không cần phải lo cho mình.''

Reiko-san còn định nói thêm điều gì đó, nhưng tôi đã quay mặt đi. Tôi biết rõ, nếu nhìn thẳng vào mắt cô ấy lúc này, tôi sẽ chẳng thể giữ được quyết định vừa rồi.

Khi cô ấy cùng hai anh em rời khỏi lớp, tôi ngồi phịch xuống ghế. Căn phòng bỗng trở nên quá rộng, quá tĩnh mịch. Trước đây, tôi vẫn hay cảm thấy thế này: nhút nhát, e dè, chỉ biết lủi thủi trong góc. Nhưng kể từ khi có một người đồng hành bên cạnh, tôi đã bắt đầu nghĩ rằng tôi có thể thay đổi.

Vậy mà giờ, tôi lại trở về vạch xuất phát. Một tôi, với bóng tối quen thuộc đang nuốt chửng lấy tôi lần nữa.

\dots Tôi đã làm gì nên nỗi chứ? Chẳng lẽ\dots\ Thần linh không muốn cho tôi một người bạn sao? Tại sao ông có thể làm thế với con chứ!?

\pov{Haragen Reiko}

Ánh hoàng hôn đổ dài trên những con phố ướt át của thị trấn Aogami, nhuộm vàng cam những vũng nước còn sót lại sau cơn mưa. Bầu trời đã tạnh hẳn, để lộ những dải mây bông trắng xốp đang trôi lững lờ trên nền trời trong vắt. Không khí mát mẻ, trong lành đến lạ - một sự thay đổi đột ngột so với những ngày mưa dầm ẩm ướt vừa qua.

Dọc hai bên đường, những cửa hiệu bắt đầu thắp đèn, ánh sáng vàng ấm áp hòa quyện với sắc hoàng hôn, tạo nên một khung cảnh thanh bình đến nao lòng. Tiếng chuông xe đạp của những học sinh tan trường, tiếng trò chuyện rộn rã của người qua đường, tất cả đều mang một sắc thái rạng rỡ hơn trong buổi chiều đẹp hiếm hoi này.

Thế nhưng trong lòng tôi lại đang cuộn trào những cảm xúc khó tả. Hình ảnh Shino-san lúc nãy - ánh mắt né tránh, nụ cười gượng gạo và cách cô ấy vội vã rời đi - cứ ám ảnh tâm trí tôi. Có điều gì đó không ổn, một bí mật nào đó đang được che giấu. Và dường như, không chỉ có mình tôi nhận ra điều đó.

Tôi ngước nhìn bầu trời đang chuyển dần sang màu tím nhạt. Những tia nắng cuối ngày như những sợi tơ vàng mỏng manh, phủ lên thị trấn một lớp ánh sáng dịu dàng mà man mác buồn. Giống như tâm trạng của tôi lúc này - một sự pha trộn giữa lo lắng và quyết tâm, giữa hoang mang và khao khát muốn tìm ra sự thật.

''Cậu ấy biết hết mọi thứ rồi\dots'' Kyuen-san khẽ nói, giọng thấp đến mức chỉ hai anh em có thể nghe thấy.

Saikyou-san thở dài, vẻ mặt hiện rõ sự cam chịu: ''\hl{Em đoán được mà. Từ lúc cậu ấy bảo em với Shino-san vào lớp trước là đã thấy sinh nghi rồi. Ai dè đâu\dots}''

Tôi bước chậm lại nửa nhịp. Thật sự, ánh mắt của Shino-san lúc ra về khi nãy\dots\ cứ như đang muốn giấu một điều gì đó. Có điều gì bất ổn đang xảy ra, nhưng tôi không sao chỉ ra được rõ ràng.

Trên đường về, chúng tôi kể lại cho nhau những gì mình đã trải qua trong giờ nghỉ.

Kyuen-san mở lời trước. Cậu kể rằng trong lúc tìm hộp bút cho Shino-san, cậu đã cùng Shiina-san đào bới từng ngăn bàn, cuối cùng cũng tìm ra. Thế nhưng chính sau đó, cô ấy lại khuyên Shino-san tránh xa tôi, ít nhất là ở trong trường. Giọng của cậu nặng trĩu, như thể trong lòng còn chưa hết bức xúc. Tôi nghe mà tim thắt lại. Shiina-san dù là một học sinh cá biệt, nhưng không phải kẻ xấu tính, vậy thì chắc hẳn cô ấy phải thấy điều gì đó nghiêm trọng mới buộc phải làm vậy.

Saikyou thì kể lại cuộc ''thử thách gan dạ'' mà Honoka-san bày ra, đồng thời cũng vừa gạ gẫm vừa chọc tức Nanase-san. Thoạt nghe thì buồn cười thật, thậm chí có phần điên rồ, nhưng ẩn trong đó là một cảm giác chẳng dễ chịu gì.

Đặc điểm chung của cả hai anh em nhà ấy, là họ đều có một sự nghi ngờ nhất định tới Nanase-san, khi những gì họ nói (mà họ thu thập được từ Shiina-san và bản thân Saikyou-san) đều cùng dẫn đến việc cô gái tóc xanh kia đã chơi xỏ Shino-san bằng việc giấu hộp bút đi, và đã thực hiện một cách rất tinh vi, bằng cách quẳng hộp bút của cô ấy vào ngăn giày, để rồi cô em gái của Kyuen-san tình cờ dựng lại chiếc cặp ở đó khiến chẳng ai tìm thấy được.

Câu chuyện ấy khiến tôi siết chặt nắm tay. Nanase-san\dots\ thật sự là đang làm gì vậy? Tôi cũng đã hiểu phần nào vì sao Shino-san lại cư xử khác thường, như thể muốn đẩy tôi ra khỏi mối quan hệ mới chớm nở ấy.

Có lẽ nào, đúng như hai anh em Haragen đặt giả thuyết\dots

\dots rằng Shino-san đang bị cô ấy bắt nạt? Nhưng kể cả thế, thì tại sao chứ? Chúng tôi cũng không có bất kỳ một bằng chứng nào để bắt tội, nhưng những gì mà con gái của vị doanh nhân ấy làm\dots\ thật sự rất xấu tính.

Còn tôi, tôi kể lại cuộc trò chuyện trong phòng Mỹ thuật với Sakura-san và Yae-san. Khi tôi hỏi đến ''bức tranh'' cũ kỹ bị che lấp trong căn phòng một cách kỹ lưỡng hơn, họ lại bỗng chốc phản ứng một cách rất kỳ lạ:

''Bức\dots\ bức tranh đó\dots'' cô em gái thở dốc, như thể chỉ cần nhớ đến thôi đã đủ khiến cô nghẹt thở.

Cậu anh trai cũng run rẩy không kém, ánh mắt lạc đi. ''Reiko-sam\dots\ cậu thực sự không muốn nhìn bức tranh ấy đâu. Tin tôi đi.''

Tôi lặng người. Bức tranh ấy ẩn giấu thứ gì, mà ngay cả hai người từng đối mặt cũng không dám thốt lên?

Câu trả lời ấy khiến tôi càng thôi thúc phải tìm ra sự thật.

Chúng tôi im lặng một quãng đường dài. Tôi không gặng hỏi thêm, vì rõ ràng cả hai đều chưa sẵn sàng. Nhưng tận trong lòng, tôi biết: muốn làm rõ chuyện đang xảy ra với Shino, tôi phải tìm ra chân tướng. Bất kể những bí mật mà họ đang giấu ấy khủng khiếp đến đâu, bất kể con đường phía trước ẩn chứa thứ gì.

Shino-san\dots\ tôi sẽ không để cậu phải chịu đựng một mình đâu.

\begin{center}
  \uline{Ngày hôm sau.}
\end{center}

Sáng nay, bầu trời như bị một tấm màn xám xịt che phủ. Không khí nặng trĩu, từng cơn gió âm ẩm thổi qua khiến làn da tôi gai lạnh. Tôi ngước nhìn, chỉ thấy mây vần vũ, dày đặc đến mức ánh sáng ban mai cũng trở nên nhợt nhạt.

Thật lạ\dots\ hôm qua, bầu trời còn trong xanh, rực rỡ. Tôi đã nghĩ, ít nhất trong khoảng thời gian này, mình có thể tận hưởng chút ấm áp bình yên. Thế nhưng, chỉ một đêm trôi qua, tất cả đã đổi khác. Thời tiết hôm nay giống như báo hiệu một điều gì đó chẳng lành.

Tôi khẽ lẩm bẩm than thở: ''Lạ thật\dots\ mới hôm qua còn nắng cơ mà.''

Bên cạnh, Kyuen-san thì xách cặp của mình với cô em gái, còn Saikyou-san thì siết chặt tay cầm chiếc ô mà vốn là của tôi, nhưng hai anh em cứ khăng khăng mang đi để đề phòng.

''Cẩn thận thì vẫn hơn,'' cô em gái mỉm cười. ''Tôi quên ô nhiều nên tôi biết mà. Cứ hôm nào quên thì kiểu trời lại mưa.''

Tôi cười gượng, chưa kịp đáp thì Kuien-san tiếp lời, giọng có chút xa xăm:
''Bầu trời kiểu này\dots\ thường là điềm chẳng lành. Nhất là khi hôm qua Shino-san bị Shiina-san ngăn cản như thế\dots''

Câu nói khiến tôi khựng lại. Saikyou cũng nghiêng đầu: ''Nii-chan, anh nói vậy nghe đáng sợ quá đấy. Nhưng mà\dots''

Cô ấy cũng thở dài một tiếng, ngước lên trời: ''\dots không phải là không có khả năng.''

''Ừ\dots'' tôi thở lấy một hơi chán nản, y như thời tiết bây giờ. ''Cả ngày hôm qua\dots\ mình cứ thấy bất an. Shino-san\dots\ trông cậu ấy không giống thường ngày nữa.''

Vừa nhắc tới, tôi chợt nhận ra một dáng người nhỏ bé ở phía trước con đường dẫn vào cổng trường. Mái tóc nâu dài, để xõa, bước chân vội vàng\dots

Chính là cô ấy. Người mà chúng tôi vừa nhắc tới.

''\textbf{Shino-san}!'' tôi gọi theo, hy vọng nụ cười của mình có thể xóa đi khoảng cách từ hôm qua.

Nhưng cô ấy quay lại, ánh mắt thoáng vẻ bối rối, rồi lập tức cúi đầu thật thấp, gần như né tránh. Cô ấy chỉ đáp lại một tiếng lắp bắp: ''\textbf{M-Mình xin lỗi!!}'' và rồi, trước khi tôi kịp bước nhanh hơn, Shino-san đã vụt chạy đi, chiếc cặp lắc lư bên hông.

Tôi chết lặng. Bước chân cô vội vã, dáng người nhỏ bé ấy dần khuất xa trong đám đông học sinh đang tiến về phía cổng trường. Tôi muốn chạy theo, muốn nắm lấy tay bạn mình mà hỏi ''Tại sao?'', nhưng đôi chân như bị ghim xuống đất.

Ngạc nhiên. Khó hiểu. Và trên hết\dots\ là cảm giác bị bỏ rơi.

Tôi đã gọi tên cô ấy bằng tất cả sự chân thành, nhưng đổi lại, chỉ là một cái bóng lưng vội vã. Một cơ hội để giải thích, để nối lại mối dây bạn bè\dots\ giờ đã vụt mất ngay trước mắt.

''Reiko-san\dots'' Saikyou-san khẽ gọi, giọng nhỏ nhẹ. ''Đừng vội buồn. Cậu sẽ có cơ hội để nói mà. Chỉ là\dots\ có lẽ bây giờ Shino chưa sẵn sàng thôi.''

Tôi cắn môi, cố gắng giữ bình tĩnh, nhưng lồng ngực nặng trĩu.

Kyuen-san thì đứng im một thoáng, đôi mắt nhìn về phía cổng trường, sâu thẳm và căng thẳng.

''Nii-chan? Sao thế?''

Đôi mắt cậu ấy vẫn ngước về phía trước. Thế rồi\dots

''Không\dots\ có gì đó không ổn. Với Shino-san\dots\ hôm nay sẽ xảy ra chuyện.''

Cậu nói dứt lời liền chạy vụt đi, bỏ lại hai cúng tôi với sự sững sờ.

''\textbf{Kyuen-san! Đợi đã!}'' Tôi gọi với, nhưng cậu không quay lại.

Chỉ còn cách nhìn nhau, tôi và Saikyou-san gật đầu, rồi vội đuổi theo cậu con trai đang lao thẳng về phía trường học, như thể phía trước thực sự có một cơn bão đang chờ đợi.

\pov{Misora Shino}

Tôi\dots\ tôi đã làm gì thế này?

Ngay khi nghe thấy tiếng gọi ấy vang lên phía sau - cái giọng nói quen thuộc, ấm áp mà tôi từng muốn bấu víu - tôi lại không đủ can đảm để quay đầu lại. Đáng lẽ ra, tôi phải mỉm cười, phải đáp lại, như trước giờ vẫn vậy. Nhưng rồi, lời của Shiina-san hôm qua chợt vang lên trong đầu, lạnh lẽo và cứng rắn đến mức khiến tôi chùn bước.

''Tốt nhất cậu nên tránh xa cô ấy ra khi đang ở trong trường.''

Vậy là, thay vì tiến về phía Reiko-san (mà vốn dĩ là đang ở bên ngoài), tôi đã bỏ chạy. Bỏ lại ánh mắt ngạc nhiên và hụt hẫng của cô ấy ở sau lưng mình. Trái tim tôi như bị ai bóp nghẹt. Một phần muốn dừng lại, quay lại giải thích. Nhưng một phần khác - sự yếu đuối mà tôi muốn chôn đi - bất chợt quay trở lại, khiến tôi sợ rằng nếu mình làm vậy, Reiko sẽ bị liên lụy.

Tôi chạy, như thể chỉ có vậy mới cứu được chính mình.

\ct

Khi tới cổng trường, tôi thở hổn hển, bàn tay run lên vì căng thẳng. Tôi nhanh chóng thay giày, rồi bước vội vào lớp học, hy vọng mọi thứ sẽ trôi qua yên bình.

Nhưng khi vừa đến gần chiếc bàn học ở cuối lớp quen thuộc ấy\dots\ tim tôi chợt khựng lại.

Ở đó có một chiếc lọ hoa. Một chiếc lọ chứa đầy những bông bỉ ngạn đỏ rực.

Tôi chết lặng. Đầu óc trống rỗng, cơ thể run rẩy.

Bỉ ngạn đỏ\dots\ loài hoa dành cho người đã khuất. Và tôi biết rõ nhất cái ý nghĩa khi nó bị đặt lên bàn một học sinh còn sống. Đó là cách ác ý nhất để nói rằng: ''Mày nên biến mất đi.''

Ai\dots\ ai lại có thể làm thế này với tôi?

Đầu gối tôi bỗng rã rời. Tôi ngồi phịch xuống ghế, mắt nhìn chằm chằm vào những cánh hoa đỏ như máu. Cái màu đỏ ấy xoáy sâu vào tim tôi, như thể đang chế nhạo rằng sự tồn tại của tôi chỉ là một cái gai trong mắt người khác.

Tại sao\dots\ tại sao lại là tôi?

Ký ức ngày xưa bỗng chốc ùa về, mỗi ký ức, mỗi trải nghiệm không mấy hay ho ấy như mỗi lưỡi dao cắt vào tâm trí tôi:

Những ngày đầu tiên đi học, bị bạn bè thì thầm chê cười vì giọng nói run rẩy.

Những lần bị gạt ra khỏi nhóm, lạc lõng trên sân chơi.

Những ánh mắt khó chịu khi tôi chỉ mới cất lời\dots ''Làm phiền quá.'' ''Đừng có bám theo.'

Những tiếng thì thầm ấy vẫn còn in hằn trong tai tôi.

Và bây giờ, tất cả lại lặp lại. Như một vòng tròn không có hồi kết.

Tôi cúi gằm mặt xuống bàn, nắm chặt vạt áo, cố giấu đi gương mặt đang méo mó vì sợ hãi. Trái tim đập loạn, vừa muốn bật khóc, vừa muốn biến mất ngay lập tức.

''Tôi\dots\ tôi đã làm gì đến nỗi này chứ?''

Câu hỏi bật ra trong tâm trí, lặp đi lặp lại, xoáy sâu, không một lời đáp. Tôi tự trách, tự ghét bỏ bản thân, tự thấy mình quá yếu đuối, quá vô dụng. Cứ thế, tôi co người lại, như muốn biến thành một cái bóng mờ nhạt để không ai còn nhìn thấy.

Nếu như\dots\ tôi chưa từng cố gắng bước ra khỏi cái vỏ của mình. Nếu như tôi chưa từng mở lòng với Reiko. Có lẽ\dots\ tôi đã chẳng phải chịu đựng điều này\dots

\pov{-}

Tiếng ồn ào của hành lang dần nhỏ lại khi Yae, Shiina, Honoka, Sakura và Nanase bước vào lớp. Ngay lập tức, tầm mắt của cả nhóm chạm đến hình ảnh Shino ngồi co rúm lại ở bàn mình, trước mặt là chiếc lọ thủy tinh chứa đầy bỉ ngạn đỏ rực.

Cả căn phòng như chùng xuống.

Cô nàng tóc ngắn khựng lại ngay tại chỗ, đôi mắt tròn xoe, miệng lẩm bẩm không thành tiếng: ''\dots Không thể nào\dots''

Honoka thì lấy tay che miệng, bước vội vài bước như định lao tới, gương mặt thoáng tái đi. Cô không rõ loài hoa kia mang ý nghĩa gì, nhưng cái việc một chiếc lọ đầy hoa tang tóc bị đặt trên bàn Shino\dots\ chỉ cần thế thôi cũng quá tàn nhẫn.

Shiina đứng chết lặng một lúc lâu. Cô không ngờ lại có người trong lớp chơi trò đê tiện đến vậy, biến lời cảnh báo của mình với Shino thành một ám ảnh thật sự. Trong lòng cô dâng lên một cảm giác nhói buốt. Ánh mắt Shiina thoáng dừng lại nơi hành lang, khi mà Reiko còn chưa đến, rồi cúi đầu trĩu nặng.

Nanase thì khác. Cô đứng sau cùng, bình thản như thể đã biết trước. Ánh mắt lạnh lẽo, khóe môi cong lên một nụ cười mờ nhạt: nửa giễu cợt, nửa thỏa mãn. Tất cả những hỗn loạn trước mặt như một vở kịch mà cô đã dàn dựng, chỉ chờ diễn ra đúng như mong muốn.

Còn Sakura, cô lặng lẽ đứng cạnh, hai tay đan lại trước bụng. Đôi mắt nghiêng về phía Nanase trong thoáng chốc, ánh nhìn sắc bén, biết rõ thủ phạm là ai. Nhưng cô không mở lời. Không phải thời điểm này.

Trong sự bấn loạn của lớp 1-B đó, khi mỗi con người đều thể hiện một thái độ riêng với cảnh tượng trước mắt, khi bầu không khí trong lớp học bỗng chốc nặng nề hơn hẳn, ba con người họ Haragen bước vào lớp.

Khung cảnh đầu tiên đập vào mắt họ là Shino đang ngồi bất động, như thể toàn bộ sức sống bị rút cạn bởi chiếc lọ hoa chết chóc trên bàn. Yae và Honoka quỳ xuống cạnh Shino, giọng run rẩy cố gắng gọi tên cô. Shiina thì đứng lặng bên cạnh, ánh mắt vừa lo âu, vừa áy náy với những gì cô nói ngày hôm qua.

Kyuen lập tức sững sờ, tay nắm chặt. Saikyou nuốt khan, còn Reiko thì cảm thấy tim mình rơi xuống một khoảng trống lạnh buốt.

''Lại là\dots\ bông hoa đó\dots'' cậu con trai lẩm bẩm.

''Như vậy là\dots\ không nghi ngờ gì nữa\dots'' cô em gái của cậu tiếp lời.

''Như thế là quá quắt rồi\dots'' cô gái có mắt hai màu lẩm bẩm.

Không ai dám bước lên ngay. Cả lớp im lặng, chỉ còn hơi thở run rẩy. Rồi Yae quay đầu nhìn người vừa dứt lời kia, đôi mắt xanh của cô hiện rõ sự hoảng loạn: ''Reiko-san\dots\ Cậu nói gì với Shino-san đi chứ? Tớ thấy hai cậu mọi khi hay thân nhau lắm mà\dots''

Honoka cũng gật mạnh, bàn tay vẫn giữ vai cô gái vẫn đang ủ rũ kia: ''Đúng đấy. Nếu là cậu\dots\ Shino-san sẽ an tâm hơn đó!''

Shiina lúc đầu định bước tới, nhưng rồi dừng lại, giọng khàn khàn cất lên: ''Reiko\dots\ Tôi nghĩ là Kyuen và Saikyou đã biết hết mọi chuyện rồi\dots\ Nhưng tôi nghĩ, chính cậu mới là người cần đến bên cô ấy để an ủi đấy.''

Câu nói ấy khiến Reiko khẽ chao đảo. Nhưng rồi, ánh mắt cô dần trở nên kiên định.

''\textit{Shino-san\dots\ Mình sẽ không để cậu đơn độc đâu.}''

Cô hít một hơi, bước đến gần Shino. Cô ngồi xuống, thấp giọng gọi: ''Shino-san\dots?''

Cô bạn khẽ run lên, nhưng không ngẩng đầu. Mái tóc che kín khuôn mặt, vai cô hơi co lại như sợ hãi.

''Là mình đây, Reiko-san.''

Cô nàng tóc nâu ngửa mặt lên, đôi mắt ứa lệ, trông bộ dạng thảm hại như vừa bị mất đi người thân. ''R-Reiko\dots san\dots?''

Giọng cô run lên, gọi lại tên của người bạn duy nhất của mình: ''Mình\dots\ mình thật sự vô dụng, Reiko à. Từ nhỏ đã luôn như thế, lúc nào cũng bị cười chê, bị xua đuổi\dots\ Cứ mở miệng là lại bị bảo `phiền phức'. Đến bây giờ cũng chẳng khác\dots\ ngay cả Shiina, người duy nhất mình nghĩ còn quan tâm, cũng bảo mình phải tránh xa cậu. Thế thì\dots\ mình còn có ai nữa đâu?'' Trong giọng nói ấy phảng phất nỗi tuyệt vọng của một người đã mất hết chỗ dựa.

Reiko lặng người, lòng nhói buốt. Hình ảnh của cô nàng lúc này khiến cô nhớ lại chính mình của những ngày cũ - cũng từng bị lời dèm pha vây quanh, cũng từng co mình lại để né tránh ánh mắt người khác. Cái cảm giác như bị cả thế giới bỏ rơi\dots\ cô hiểu rõ hơn ai hết.

Cô khẽ gật đầu, rồi chống hai tay lên bàn của cô bạn. Cô đưa tay ra, nhẹ nhàng đặt lên bàn tay Shino đang nắm chặt vạt áo. Giọng cô trầm ấm, chậm rãi nói: ''Xin lỗi\dots\ vì hôm nay cậu lại phải chịu thêm chuyện như thế này. Nhưng Shino-san, cậu không hề đơn độc. Nhìn xem\dots\ mọi người đều đang ở đây, vì cậu.''

Shino lắc đầu, giọng nghẹn ngào: ''Nhưng\dots\ tại sao lúc nào cũng là mình\dots? Mình\dots\ mình chẳng làm gì cả, vậy mà\dots''

Reiko mím môi, ánh mắt dứt khoát. ''Cậu có biết không, Shino-san? Chính sự dịu dàng của cậu, cái cách cậu quan tâm người khác, đã khiến mình học được rất nhiều. Mình chưa bao giờ nghĩ cậu là gánh nặng. Với mình\dots\ cậu là người mà mình muốn bảo vệ.''

Nước mắt của Shino chực trào, nhưng cô vẫn ngẩng đầu lên, đôi mắt đẫm lệ nhìn cô. ''Nhưng\dots\ nếu mình cứ kéo cậu vào những rắc rối này thì sao? Cậu sẽ bị tổn thương mất\dots''

Reiko lắc đầu, bàn tay siết nhẹ lấy tay Shino. ''Nếu là vì cậu, mình không hối hận. Mình thà cùng cậu đối mặt còn hơn là để cậu chịu đựng một mình.''

Nước mắt cô nàng tóc nâu rơi xuống, hòa vào những lời trong tiếng nấc: ''Nhưng\dots\ nếu rồi tất cả lại rời bỏ mình thì sao? Mình không chịu nổi lần nào nữa đâu\dots''

Reiko lắc đầu, siết tay cô chặt hơn: ''Thì mình sẽ là người ở lại. Cho dù thế giới có đổi thay thế nào.''

Lời nói ấy như một sợi dây tuy mảnh khảnh nhưng rắn rỏi níu Shino khỏi vực thẳm. Nỗi sợ vẫn còn, nhưng một hơi ấm bắt đầu lan dần trong lòng, nhắc nhở cô rằng mình không thật sự đơn độc.

Không khí lặng đi, chỉ còn tiếng nấc khẽ. Người con gái nhút nhát áp mặt xuống vai cô bạn, nước mắt thấm ướt vạt áo bạn mình. Trong vòng tay ấy, cô dần bình tĩnh lại, nỗi sợ tan bớt, chỉ còn một khoảng ấm áp len lỏi trong lồng ngực.

Yae, Honoka và Shiina cũng bước đến gần, mỗi người một lời động viên:

Cô nàng Cơ bắp dịu dàng đặt tay lên lưng Shino: ''Đừng nghĩ mình yếu đuối nữa. Cậu vẫn đang đứng vững đấy thôi. Với cả, cậu cũng đã có bọn tớ ở cạnh rồi.''

Cô bạn bờm Cáo thì nghiêng đầu, mỉm cười run run: ''Bất hạnh nào rồi cũng qua đi, Shino-chan à. Cậu sẽ không phải gánh nó một mình nữa.''

Còn cô bạn bờm Sư tử thở dài, giọng nhẹ như gió thoảng: ''\dots Xin lỗi vì đã khiến cậu thêm nặng nề. Nhưng tôi cũng mong cậu sẽ tin vào bản thân nhiều hơn.''

Shino chỉ biết nghẹn ngào gật đầu, cảm nhận rõ những vòng tay, những lời nói đang giữ chặt lấy mình. Cô an tâm hơn, nhưng trong đáy lòng, một vệt bóng tối vẫn âm ỉ, chưa tan biến hẳn.

Từ xa, Nanase đứng khoanh tay, nét mặt tối sầm lại. Nhìn thấy Shino được Reiko và các bạn an ủi, lòng cô bốc lên một cơn ghen ghét mơ hồ, như thể thứ gì vốn thuộc về cô đang bị cướp đi. Nụ cười khi nãy biến mất, thay vào đó là ánh mắt u ám.

Sakura liếc sang, thấy tất cả. Nhưng cô chỉ im lặng, đứng tựa vào bàn, trong đầu vang lên một dòng suy nghĩ:

''nase-san\dots\ rốt cuộc cậu định đi bao xa nữa đây?''

Ngay phía sau, Kyuen và Saikyou cũng không rời mắt khỏi cảnh tượng. Cả hai bắt đầu xâu chuỗi lại những gì đã diễn ra từ lúc trước khi họ bị chuyển dịch về thời kỳ này - khoảnh khắc Shino hầu như chỉ còn là một cái hồn đầy hận thù, mong muốn ''đưa người bạn học sinh mới chuyển đến đi về nhà''.

Cô em gái khẽ nghiến răng, thì thầm: ''\hl{Thì ra là\dots{} do bị dồn ép đến mức này.}''

Người anh thì mím môi, trong lòng dấy lên cảm giác bất lực lẫn tức giận. Một người vốn chỉ muốn sống yên ổn, vậy mà hết lần này đến lần khác bị giày vò\dots\ Giờ đây cậu mới hiểu rằng lý do Shino lảng tránh không phải do vô cớ, mà bởi vì cô đang cố gắng sống sót trong một vòng tròn bắt nạt vô hình.

Cả hai đồng thời hướng ánh nhìn về Reiko, người còn lại trong số ba bọn họ cũng là một ''học sinh mới chuyển đến.'' Mảnh ghép cuối cùng bỗng khớp lại, như một câu trả lời mơ hồ mà họ không muốn tin, nhưng lại không thể phủ nhận.

Họ sau đó\dots\ cũng nhìn thấy Nanase. Một cô gái với vẻ mặt lạnh như băng, không đoái hoài gì tới chuyện trước mặt.

Không khí trong lớp vẫn trĩu nặng, nhưng giữa nó, ánh sáng mong manh từ sự chân thành của cô gái có mắt hai màu đang xoa dịu nỗi đau của Shino.

\ct

\pov{Haragen Kyuen}

Giờ nghỉ giữa giờ đầu tiên trong ngày.

Tôi chống cằm nhìn trống rỗng về phía bàn học của Shino-san. Lọ hoa bỉ ngạn đỏ ấy vẫn còn in rõ trong đầu tôi như một vết loang khó gột rửa. Mới hôm kia nó xuất hiện, rồi sáng nay lại thêm một lần nữa\dots\ Có gì đó không ổn.

''Rõ ràng đây không phải trùng hợp,'' Saikyou hạ giọng, hai tay khoanh lại trước ngực. Em gái tôi nghiêm nghị đến mức làm tôi chột dạ. ''Một trò ác ý. Người đó không chỉ làm một lần, mà đến tận hai lần\dots\ Cá chắc là hắn ta hận Shino-san lắm đấy.''

Tôi thở dài, gõ nhẹ ngón tay xuống mặt bàn. ''Vấn đề người đấy là ai. Chứ nói suông thì chả ai nói được.''

Em tôi nheo mắt, đảo nhìn quanh lớp. Nanase-san không có ở chỗ ngồi, mới ra ngoài đâu đó. Trong đầu tôi và em gái cùng lúc hiện lên cái tên ấy. Cô ta từ hôm chúng tôi nhập học\dots\ đã chẳng bình thường. Thái độ của cô ta, nụ cười nhếch mép, dửng dưng đứng ngoài\dots\ và cả cách cô ta lảng tránh khi nói chuyện về Shino-san.

''Anh cũng nghĩ thế sao?'' Saikyou hỏi, ánh mắt sắc như muốn đọc suy nghĩ tôi.

Tôi gật đầu. ''Ừ. Cái cảm giác đó\dots\ nó giống như lúc ở phòng Mỹ thuật. Em còn nhớ không?''

Đôi mắt Saikyou chợt chùng xuống. Tôi biết trong đầu em ấy đang hiện lại cảnh tượng ấy: căn phòng trống hoắc chỉ có vài giá vẽ, lời đồn về cô gái đã chết, bức tranh nhuốm màu máu. Và cả giấc mơ đêm định mệnh lúc đó, khi chúng tôi nhìn thấy cái thứ quái dị kia - bức tranh trông như cô gái đang gào thét kia, đè nặng lên tâm trí. Nếu hôm ấy không có ''bàn tay'' nào đó xé toạc khung cảnh, có lẽ hai anh em tôi đã phát điên. Thậm chí\dots\ tôi còn không chắc mình là ai cho tới khi Saikyou gọi tôi bằng cái tên này.

Tôi đang định mở miệng thì một giọng nói xen ngang: ''Có vẻ như không chỉ mỗi hai cậu nghĩ thấy có điều gì đó kỳ lạ nhỉ.''

Tôi giật mình, quay sang.

Shinomya Sakura. Cô đứng đó, ánh mắt tĩnh lặng nhưng sâu hun hút, như thể nhìn xuyên qua lớp da thịt vào tận lõi của chúng tôi.

Cô ngắm tôi và Saikyou một lúc lâu, rồi khe khẽ mỉm cười. Không giống nụ cười chào xã giao, mà giống như cô vừa khẳng định một linh cảm mơ hồ.

''Tôi không rõ gọi nó là gì,'' Cô nàng tóc hồng nói, ''nhưng cả hai cậu\dots\ tỏa ra một thứ sức mạnh. Khác Reiko-san, nhưng cũng\dots\ không bình thường. Cứ như thể\dots''

Câu nói bỏ lửng, nhưng tim tôi thoáng loạn nhịp. Phải. Không phải chỉ vì chúng tôi đến từ năm 2018 trong khi đây là 1946. Nó lại là\dots\ cái cảm giác tôi không hoàn toàn là chính mình. Cái tên Kyuen này, nếu không phải nhờ cô em gái gọi, tôi còn chẳng biết để tự xưng là gì.

Thậm chí tôi còn chẳng nhớ bản thân mình là ai nữa. Tại sao tôi lại ở đây? Vì mục đích gì? Cũng không có một lời giải đáp nốt.

Saikyou thì vẫn thản nhiên, nhưng tôi thấy khóe môi em gái run nhẹ. Có lẽ cô cũng cảm nhận được phần nào điều bất thường trong chính bản thân mình.

Không khí lắng xuống trong một thoáng.  Sakura-san đột nhiên xoay người, đôi mắt cô lia ra cửa sổ. Tôi tò mò nhìn theo\dots\ và thấy Shino-san.

Cô ấy đứng đó, bất động. Ánh sáng xám xịt từ bầu trời âm u hắt qua, phủ lên gương mặt nhợt nhạt, đôi mắt trống rỗng. Hai bàn tay siết chặt lấy nhau, trông như đang níu giữ chút can đảm cuối cùng để không gục ngã.

''Này\dots'' Cô Vu nữ khẽ gọi, rồi bước đi thẳng tới.

Hai chúng tôi nhìn nhau. Cả hai cùng hiểu: có lẽ đã đến lúc phải lại gần Shino-san rồi.

''Shino-san? Sao thế? Cậu trông có vẻ thẫn thờ thế kia\dots'' Cô nghiêng đầu, giọng nhẹ nhưng đủ khiến người khác khó mà làm ngơ.

Cô nàng tóc nâu thoáng giật mình, quay sang, đôi mắt hơi hoảng hốt như vừa bị bắt gặp trong lúc mơ màng. ''\dots À, ờm\dots\ Shinomiya-san? Có phải là tên cậu không?''

''Đúng rồi.'' Sakura-san mỉm cười, điềm tĩnh. ''Tên tôi là Shinomiya Sakura. Với cả, gọi tôi là Sakura là được rồi. Thế rốt cuộc thì, sao cậu lại trông ủ rũ vậy?''

Shino-san bối rối, hai tay siết chặt vạt áo. ''\dots À, không có gì đâu. Chỉ là mình\dots\ nghĩ mấy thứ vẩn vơ thôi.''

Tôi và Saikyou liếc nhau. Rõ ràng cô ấy vẫn còn đang bị chấn động sau sự vụ sáng nay.

''\dots Phải rồi, Shino-san.'' Sakura chống tay lên bàn, giọng nhỏ đi nhưng nghiêm túc hơn. ''Tôi có điều này muốn hỏi cậu.''

''Gì thế?''

''Cậu có thấy gần đây xung quanh cậu xảy ra điều gì kỳ lạ không?''

''Điều\dots\ kỳ lạ?'' Shino nhíu mày, ngập ngừng.

Sakura-san gật nhẹ. ''Nó là một thứ gì đó rất mơ hồ và cũng không chắc chắn, nhưng\dots\ cái giác quan vu nữ của tôi đang nói lên gì đó. Thực sự có điều gì đó xảy ra ở đây.''

Shino-san thoáng mở to mắt, rồi lại cúi xuống. ''\dots À, gia đình Sakura-san đang điều hành một ngôi đền gần đó mà. Nếu thế thì\dots\ mình cũng không rõ. Cũng chẳng liên quan gì đến mình nữa\dots ''

Sakura-san nghiêng người. ''Cái gì không liên quan cơ?''

''\dots Không có gì đâu. Ít nhất thì, cũng không liên quan gì tới mấy cái bùa chú hay là tâm linh mà cậu đang nghĩ tới cả,'' cô gượng cười, nhưng giọng run rẩy.

Tôi khoanh tay, bước lên nửa bước. ''Này, Shino-san. Tôi thấy cậu đang giấu điều gì đó. Nếu cậu không nói, thì làm sao bọn tôi biết để giúp?''

Shino-san cắn môi, mắt chớp liên tục như cố kìm nước mắt. ''\dots Thực ra thì\dots\ mình cũng chẳng biết bắt đầu từ đâu. Chỉ là\dots\ gần đây mình cứ có cảm giác mình thật sự xui xẻo. Chính xác hơn là\dots\ bất hạnh.''

Saikyou khoanh tay, giọng trầm hẳn. ''Có ai đó đã làm gì cậu, đúng không? Không phải tự nhiên mà lọ hoa đó xuất hiện tận hai lần đâu.''

Shino sững lại, rồi cúi đầu. Một khoảng im lặng nặng nề. ''\dots Ừ. Có vài người\dots\ họ hay châm chọc mình. Ban đầu chỉ là mấy câu nói thôi, kiểu `cậu lạc loài quá', hay `trông cậu yếu ớt nhỉ'. Nhưng rồi\dots\ có lúc đồ đạc của mình biến mất. Thậm chí\dots\ hộp bút mình cũng bị ném đi. Lọ hoa sáng nay\dots\ mình nghĩ, chắc lại là trò của ai đó ghét mình.''

Giọng cô run rẩy. ''Mình\dots\ cảm thấy mình thật sự thất bại. Mình chẳng làm được gì cả, chỉ biết cúi đầu chịu đựng. Mà rồi khi Shiina-san\dots\ nói rằng có lẽ nên tránh xa Reiko-san\dots\ mình lại càng thấy lạc lõng. Thật sự\dots\ mình không biết phải làm gì nữa.''

Sakura-san khẽ cau mày, ánh mắt dịu xuống. ''Shino-san\dots''

Tôi hít một hơi, rồi nói chắc nịch: ''Nghe này. Bọn tôi cũng từng trải qua cảm giác đó tới đường cùng: bị người khác xem thường, bị dồn ép. Nhưng nếu cứ ôm lấy suy nghĩ `mình thất bại' thì cậu sẽ bị họ nuốt chửng mất.''

Saikyou gật đầu, giọng sắc bén: ''Đúng đấy. Cậu cứ tự cho mình yếu đuối như thế thì làm sao có thể tồn tại ở môi trường đầy rẫy sự độc hại đó được?''

Shino-san đan tay, bàn tay run run. Cô nàng Vu nữ liền chen vào, giọng nhẹ nhàng hơn: ''Tôi nghĩ Shino-san cần một chút an tâm. Nếu cậu thấy mình xui xẻo, tôi có thể mang cho cậu một cái bùa từ đền nhà tôi. Chúng thực sự rất tốt trong việc xua đuổi bất hạnh.''

Cô bạn tóc nâu ngẩng lên, hơi sững người. ''Một cái bùa\dots\ có thể xua đuổi bất hạnh sao?''

''Ừ. Nhiều người đã dùng và nói rằng nó hiệu quả. Tôi nghĩ cậu cũng nên có một cái.'' Sakura mỉm cười.

''\dots Cảm ơn cậu, Sakura-san. Có một cái bên cạnh chắc sẽ làm mình yên tâm hơn.'' Shino-san khẽ gật, môi nở nụ cười yếu ớt.

''Mai tôi sẽ mang tới. Cậu cứ chờ nhé.''

''Ừm\dots''

Tôi lặng nhìn cảnh tượng ấy, trong lòng xen lẫn nhiều thứ. Có một phần tôi thấy thương cho cô ấy, nhưng cũng có một phần đang rối loạn. Không chỉ vì những trò bắt nạt mà cô ấy đang chịu, mà còn vì\dots\ có cái gì đó mờ ám hơn thế đang vây quanh.

Saikyou khẽ nghiêng đầu, thì thầm chỉ đủ tôi nghe: ''Anh thấy rồi chứ? Đây không chỉ là chuyện bắt nạt bình thường. Có gì đó sâu hơn, và Reiko chắc chắn sẽ bị lôi vào.''

Tôi siết chặt nắm tay. ''Ừ. Cảm giác này\dots\ giống như khi chúng ta ở trong giấc mơ kia. Có thứ gì đó đang rình rập, chỉ chờ chúng ta sơ hở.''

Trong khoảnh khắc ấy, tôi thấy rõ: nỗi sợ hãi của cô ấy không chỉ đến từ những lời châm chọc. Nó còn đến từ bóng đen nào đó đang lan rộng trong lớp học này, và chúng tôi, bằng một cách nào đó, đã bị cuốn vào.

\pov{Haragen Reiko}

Từ chỗ ngồi của mình, tôi vẫn để mắt sang bên kia lớp. Shino-san đang được Sakura-san và hai anh em ấy vây quanh. Cả ba đều trò chuyện rất nhiệt tình, còn Shino-san thì thỉnh thoảng cúi đầu, gượng gạo đáp lại. Lẽ ra tôi nên thấy yên tâm vì Shino không còn đơn độc nữa\dots\ nhưng lạ thay, trong lòng tôi vẫn nhoi nhói bất an.

Đúng là tôi cùng mấy bạn học khác đã ra sức động viên cô ấy, nhưng nhìn cảnh ấy, trong lòng tôi cứ chộn rộn. Liệu mình có nên lại gần thêm không? Hay là chỉ làm Shino-san ngột ngạt hơn?

Trong lúc còn đang băn khoăn, một giọng nói hồ hởi vang lên ngay bên cạnh tôi.

''Này này!'' Honoka-san chống hai tay lên bàn tôi, mắt long lanh. ''Tớ có cái này muốn nói! Cực kỳ thú vị đấy! Muốn nghe không?''

Tôi giật mình. ''\dots Ờm, Honoka-san, mình--''

''Dĩ nhiên là cậu muốn nghe rồi!'' Cô ấy vẫy tay cắt ngang, cười tươi rói.

Tôi chỉ biết gượng cười, bất lực với cô bạn lém lỉnh ấy. ''\dots Thôi được rồi. Vậy, cái gì thế?''

''Tại sao chúng ta không lẻn vào trong trường để đi săn ma nhỉ?''

Tôi trợn mắt. ''\dots Gì cơ?''

''Tức là thế này,'' cô bạn tóc vàng hạ giọng, nhưng đôi mắt thì sáng rực như trẻ con sắp được quà. ''Cậu biết tin đồn về con ma trong phòng Mỹ thuật chứ? Cả lớp đang bàn tán ầm ầm luôn ấy!''

''Ừm\dots'' tôi chớp mắt, khẽ gật đầu. ''Kyuen-san và Saikyou-san cũng đã kể với mình hôm qua rồi.''

''Đây nhé, đây nhé!'' cô vẫy tay đầy phấn khích. ''Người ta nói trong căn phòng đó có một cô gái, vì quá dành nhiều tâm huyết vào kiệt tác của mình--''

''\dots mà cuối cùng cô ấy đã bỏ mạng ở đó rồi chứ gì?'' tôi thở ra, nhớ lại từng lời mà mọi người kể. ''Mình cũng đã nghe rồi. Thậm chí Yae-san còn nhấn mạnh chi tiết về\dots\ bức tranh nữa.'' Một cảm giác gai lạnh lướt dọc sống lưng khi nhớ tới ánh mắt nghiêm nghị của cô ấy lúc đó.

Honoka-san phồng má phụng phịu. ''Ể\~{}!? Cậu biết hết rồi sao? Mất vui thật\dots''

Tôi bật cười nhẹ. ''Xin lỗi nhé. Chắc do tin đồn lan nhanh quá, nên hầu như ai cũng biết rồi.''

''Thế càng tốt!'' cô chống nạnh, dõng dạc. ''Tớ định rủ cậu tham gia kiểm chứng xem liệu điều đó có thật hay không!''

''Mình còn chưa nói xem liệu mình có đồng ý hay không mà\dots''

''Ái chà chà\dots'' một giọng khác chen ngang. Shiina-san, không biết xuất hiện từ khi nào, đã đứng ngay sát chúng tôi, môi khẽ nhếch. ''Nghe có vẻ vui đấy. Tôi có thể tham gia không?''

''Được chứ!'' Honok-san reo lên. ''Càng đông càng vui mà! Reiko-san cũng sẽ đi, đúng không?''

''Ê, khoan đã--'' tôi vội lên tiếng, nhưng Honoka-san đã nhanh hơn: ''Đúng không, đúng không nào!?''

''Săn ma cơ à?'' cô nàng bờm Sư tử chắp tay sau lưng, nghiêng đầu đầy hứng thú. ''Dạng thử thách gan dạ vào ban đêm nhỉ?''

''Chính xác luôn!'' cô bạn tóc vàng gật mạnh.

Shiina-san mỉm cười nhàn nhạt: ''Nếu thế thì chuyện nhỏ như con thỏ thôi. Thế cậu đã có kịch bản gì chưa?''

''À thì\dots'' Honoka-san đưa tay ra như đang dẫn chuyện, thực hiện động tác lẻn vào bên trong: ''Đầu tiên, chúng ta sẽ lẻn vào trường vào đêm nay! Từ đó--''

''Khoan đã,'' tôi cắt lời, khẽ chau mày. ''Cậu không tính đến chuyện trường học sẽ bị khóa cửa sao?''

''Đúng vậy,'' cô bạn Cá biệt gật theo. ''Cậu có chìa khóa dự phòng không?''

Honoka-san chớp mắt, ngẩn ra vài giây, rồi trố mắt như vừa quên mất một chi tiết quan trọng. ''\dots\ À\dots\ ờ ha!''

Tôi và Shiina-san chỉ biết cùng thở dài. ''Trời ạ, cậu đúng là\dots'' / ''Cậu chẳng tính gì hết à\dots''

Cô bạn chống cằm, liếc sang tôi. ''Reiko, cậu có cao kiến gì không?''

''Hể? Mình\dots'' tôi lúng túng. ''Thật ra\dots\ mình chưa bao giờ ra khỏi nhà vào ban đêm kể từ lúc bắt đầu đi học, nên cũng chẳng biết nữa.''

Shiina-san nhướng mày. ''Một dạng gái ngoan điển hình nhỉ?''

''Ừ.'' Tôi cười ngượng. ''Nhưng mà\dots\ mình có một cách.''

''Cách gì, cách gì thế!? Chỉ tớ đi!'' Honoka-san bỗng chốc mắt sáng rực lên, sà thẩng vào bàn tôi.

Tôi khoanh tay lại, cố nghĩ ra một ý tưởng: ''Có thể thử trốn trong một lớp học trống sau giờ tan học, rồi chờ cho tới khi các bác lao công về hết.''

''Cũng được,'' Shiina-san khẽ gật. ''Mỗi tội\dots\ không hiệu quả với nhóm lớn. Sớm muộn cũng bị phát hiện thôi.''

''Đó chỉ là ý tưởng của mình thôi,'' tôi cúi đầu, nhớ lại tuổi thơ. ''Hồi bé mình từng trốn bố mẹ đi chơi đêm với lũ trẻ trong xóm. Có lần bố mẹ mình tìm cả đêm đến mức lạc cả giọng.''

Cô bạn bật cười khẽ: ''Xem ra Reiko cũng có quá khứ dữ dội phết nhỉ.''

Trong khi đó, Honoka-san vẫn loay hoay nghĩ ngợi. Rồi bất chợt cô ấy đập tay cái ''bốp'': ''À! Phải rồi! Tại sao chúng ta không để một người trốn trong lớp học nhỉ?''

''Cậu vừa nói lại cái mà Reiko nói đấy,'' Shiina-san thở dài. ''Không hiệu quả với nhóm lớn đâu.''

''Không, không!'' Honoka-san vẫy tay lia lịa. ''Ý tớ là, không cần tất cả mọi người cùng lúc. Chỉ cần một người thôi! Ví dụ như tớ chẳng hạn! Tớ sẽ trốn lại một mình, chờ lao công về, rồi lấy chìa khóa trường từ đâu đó, sau đó mở cửa từ bên trong. Xong! Thiên tài quá còn gì!''

Shiina-san im lặng vài giây, rồi nhún vai. ''\dots Ừ, cũng không phải không được.''

''Nhưng mà, Honoka-san\dots'' tôi nhìn cô ấy lo lắng. ''Cậu sẽ ổn chứ? Một mình như vậy\dots''

''Đừng lo! Tớ nắm rõ ngôi trường này như lòng bàn tay rồi. Chỗ nào mà các thầy cô hay đi qua, chỗ nào tối dễ trốn, tớ đều biết hết!'' cô bạn tóc vàng giơ ngón cái lên, khuôn mặt hiện rõ sự tự tin.

Shiina-san gật nhẹ. ''Thôi thì\dots\ tùy cậu. Nhưng đừng để bị bắt đấy.''

''Hai người cứ lo xa\~{}'' Honoka-san khoát tay, vẻ đầy hào hứng.

Tôi còn định hỏi thêm thì chợt buột miệng: ''\dots Vậy, cậu dự định khi nào?''

''Đêm nay luôn!'' cô nàng nói không cần suy nghĩ.

Tôi ngạc nhiên: ''Nhanh vậy sao?''

Shiina-san nhướng mày. ''Mới bàn kế hoạch năm phút mà đã triển khai ngay rồi hả?''

Honoka-san cười hề hề. ''Càng nhanh càng hồi hộp chứ sao!''

Cứ thế, hai người cứ bàn qua bàn lại như thể đã chốt hẳn kế hoạch. Tôi đứng bên cạnh, chẳng kịp chen vào, chỉ có thể cười trừ. Trong lòng, sự bất an càng dâng cao: không chỉ vì đây là lần đầu sau nhiều năm tôi mới ra khỏi nhà vào ban đêm, mà còn vì Shino. Liệu cậu ấy có tham gia cùng không? Hay sự xuất hiện của tôi trong cuộc thử thách này sẽ khiến Shino cảm thấy khó chịu hơn, nhất là sau chuyện xảy ra sáng nay?

Cô bạn Sư tử liếc tôi, nửa đùa nửa thật: ''Mà này, tôi nghe nói Reiko còn chưa đồng ý đấy nhể?''

''Ối dào, thấy cậu ấy cũng góp ý nhiệt tình thế thì coi như đồng ý rồi còn gì!'' cô bạn Cáo sa mạc cười toe, chẳng buồn để tôi thanh minh.

Tôi chỉ biết im lặng, cố giấu đi tiếng thở dài. Có lẽ\dots\ tôi sẽ phải hỏi ý kiến Kyuen-san và Saikyou-san sau giờ học.

\pov{-}

Giờ ra chơi dần trôi qua, trong lớp học chỉ còn lại tiếng xào xạc của vài quyển vở lật nhanh, xen lẫn những cuộc trò chuyện lẻ tẻ. Ấy vậy mà, không khí chẳng hề nhẹ nhõm. Ngược lại, thứ gì đó vô hình cứ lẩn quẩn trên đỉnh đầu mọi người, khiến cả căn phòng như nặng trĩu.

Yae ngồi ở bàn cuối, chống cằm nhìn ra ngoài cửa sổ. Từ sáng tới giờ, quá nhiều chuyện kỳ quặc đã xảy ra. Từ cái lọ hoa bỉ ngạn đỏ đặt trên bàn Shino, đến những lời đồn thổi trong phòng Mỹ thuật mà cô nghe thoáng qua\dots\ mọi thứ khiến cô chẳng thể tập trung nổi.

''\textit{Chẳng lẽ\dots\ ngôi trường này thật sự bị ám?}'' suy nghĩ ấy thoáng lướt qua đầu, dù chính cô cũng không tin lắm. Nhưng cái cách mà không khí lớp học lúc này trở nên nặng nề\dots\ lại khiến người ta khó phủ nhận.

Cô vừa thở dài thì *RẦM!*,  một tiếng động vang lên khiến cả lớp giật mình. Nanase sồn sồn bước tới bàn Yae, bàn tay đập mạnh xuống mặt bàn gỗ, khiến lọ mực trên bàn rung lên bần bật.

''Yae-san.'' Giọng Nanase gằn xuống, khuôn mặt cau có đến mức khó nhận ra.

''G-Gì vậy!?'' Yae giật thót, cảm tưởng như tim cô vừa rụng xuống.

''Tiết vừa rồi\dots\ cậu đang định chơi bài gì vậy?''

''Hể? Tiết vừa rồi? Tớ đâu có--''

''Cậu biết tôi đang nhắc tới ai mà,'' ánh mắt Nanase lóe lên, như một tia dao bén chĩa thẳng. Không cần nói ra, Yae cũng hiểu: cô đang nói tới cô bạn có đôi mắt hai màu.

Không khí trong lớp đông cứng lại. Vài bạn học khẽ thì thầm:

''Lại nữa rồi\dots''

''Nanase-chan sao thế nhỉ\dots?''

''Trông đáng sợ quá\dots''

Yae nuốt khan, cố giữ bình tĩnh: ''V-Về chuyện đó thì\dots\ lúc đó là tiết Thể dục, và tớ lại được phân công trực nhật\dots''

''Đừng ngụy biện nữa. Tôi còn chưa nói đến chuyện sáng nay cơ,'' Nanase liền cắt lời, nghiến răng, giọng ngày càng sắc. ''Kyuen-san và Saikyou-san thì thôi, nhưng cậu với Shiina-san và Honoka-san\dots\ các cậu cũng tham gia vào chuyện đó sao?''

Cả lớp xôn xao hơn. Rõ ràng cô ấy đang ám chỉ chuyện họ an ủi Shino sau sự cố lọ hoa hồi sáng.

Yae lặng người. Rồi, sau một hồi định thần, cô lấy hết can đảm để cất tiếng: ''Nanase-chan, mấy ngày hôm nay cậu bị làm sao vậy? Mọi ngày cậu đâu có cư xử thế này\dots''

Nanase nhìn xoáy vào cô, khuôn mặt tối sầm lại: ''Yae-san, tôi tưởng cậu là bạn thuở nhỏ của tôi. Có đúng không?''

''Tất nhiên là thế rồi!'' cô bạn gật mạnh, giọng lạc đi. ''Lúc nào tớ cũng là bạn của cậu, và vẫn sẽ luôn như thế!''

''Nếu vậy thì\dots\ cậu sẽ luôn giúp đỡ tôi, bất kể là điều gì, đúng chứ?''

''Cái đó thì\dots''

''\textbf{Đúng không!?}'' Nanase gằn giọng, đôi mắt rực lên sự bức bách.

Yae nín thở một hồi, rồi cuối cùng bật ra, gần như than thở: ''Nanase-chan\dots\ cậu thực sự đang cư xử rất kỳ đấy. Tớ thấy\dots\ cậu không còn là Nanase-chan mà tớ biết nữa. Chuyện gì xảy ra với cậu vậy\dots?''

''Kỳ lạ ư? Tôi á?'' Nanase bật cười nhạt, nhưng nụ cười ấy méo mó, khiến ai nhìn cũng phải rùng mình. ''Có mà cậu mới là người có vấn đề, Yae-san.''

Cô nàng tóc xám khẽ lắc đầu, giọng run rẩy: ''Không\dots\ không thể nào\dots\ Cậu\dots\ cậu đã trở thành một con người khác rồi. Tớ còn không nhận ra cậu nữa\dots''

Cô nàng tóc xanh nheo mắt, hơi cúi xuống, giọng lẩm bẩm: ''\dots Tôi hiểu rồi. Cậu cũng đang muốn chiếm lấy \emph{cô ấy} từ tôi nữa.''

''Cái\dots\ gì cơ?'' Yae sững sờ.

Ánh sáng trong mắt Nanase như dần tắt, để lại một khoảng trống tối mịt. Cô lẩm bẩm, tiếng nói rời rạc như vọng ra từ một người hoàn toàn khác: ''Bất cứ lúc nào tôi đi qua\dots\ tôi nhìn thấy ai cũng chỉ toàn ngáng đường\dots\ Tôi phải làm cho họ biến đi\dots\ Biến hết đi\dots\ Biến hết đi\dots''

Không khí trong lớp học đột nhiên trở nên lạnh lẽo và nặng nề.

Yae, người đang trực tiếp đối diện với Nanase, hoàn toàn sững sờ và kinh ngạc. Cô ấy chỉ có thể thốt lên ''Cái\dots\ gì cơ?'' với vẻ mặt không thể tin được vào những gì đang diễn ra trước mắt mình.

Các bạn học khác trong lớp chứng kiến cảnh tượng này với vẻ mặt hoảng sợ và bàng hoàng. Không ai dám lên tiếng hay can thiệp khi thấy cô bạn tóc xanh kia đang cư xử không khác gì một kẻ mất trí.

Sự im lặng bao trùm cả lớp học, chỉ còn tiếng lẩm bẩm đáng sợ của Nanase vang vọng trong không gian. Mọi người đều cảm thấy một luồng hơi lạnh chạy dọc sống lưng, không ai dám nhúc nhích hay lên tiếng, như thể họ đang chứng kiến một hiện tượng siêu nhiên đáng sợ đang diễn ra ngay trước mắt mình.

''Nanase-chan!?'' Yae hốt hoảng, vội sà tới nắm chặt vai cô bạn. ''Có gì đó không ổn với cậu rồi! Cậu\dots\ bị ốm sao!?''

Nanase chợt ngừng lại. Ánh mắt màu xanh đầy sao của cô dần trong trở lại, nét mặt mềm đi, như thể vừa tỉnh mộng.

''\dots Hả?'' cô thốt ra khẽ khàng. ''Không\dots\ tôi ổn mà. Chỉ là hơi\dots\ chóng mặt thôi.''

''Nhưng mà, kể cả thế thì--''

''Thật đấy. Tôi có thể tự đi được, đừng lo cho tôi. Nhưng\dots\ cảm ơn cậu nhé, Yae-san.''

Nói rồi, Nanase quay lưng, bước ra khỏi lớp với dáng đi lảo đảo một tay vẫn ôm lấy đầu.

Yae đứng đó, ngơ ngác nhìn theo, bàn tay vẫn run nhẹ. Trong lòng cô dấy lên một nỗi sợ mơ hồ: ''Đó\dots\ thực sự là Nanase-chan sao? Hay là\dots\ một ai khác đã chiếm lấy cậu ấy rồi?''

Trong lớp, mọi người lặng ngắt. Sự im ắng như có thể bóp nghẹt cổ họng, để lại một bầu không khí chẳng ai dám lên tiếng phá vỡ.

Kyuen khẽ liếc sang Saikyou và Reiko. Cả ba người đều hiểu rằng, để hiểu rõ Nanase đang thay đổi như thế nào, họ cần biết cô vốn là người ra sao.

''Nanase-san\dots\ vốn dĩ là người thế nào vậy?'' cô gái mắt hai màu cất tiếng hỏi, giọng điềm đạm nhưng dứt khoát. Câu hỏi khiến vài ánh mắt đổ dồn sang, thoáng ngập ngừng.

Một nữ sinh ngồi bàn gần đó lên tiếng trước, giọng dè dặt: ''Trước đây, Nanase-chan là người rất chính trực. Cậu ấy ghét sự bất công, lúc nào cũng đứng ra bảo vệ những người yếu thế. Dù xuất thân từ gia đình giàu có, nhưng cậu ấy chưa từng tỏ vẻ kiêu căng gì cả.''

Shiina tiếp lời, nhớ lại: ''Ừ\dots\ tôi còn nhớ hồi đầu năm, có lần một bạn nam bị bắt nạt ở lớp bên. Chính Nanase đã đứng ra ngăn chặn, thậm chí còn báo với giáo viên. Chúng tôi lúc đó ai cũng nể phục cậu ấy.''

Honoka gật gù, thêm vào: ''Thực ra, gia đình Nanase-chan cũng rất có thế lực. Bố cậu ấy là một doanh nhân nổi tiếng, còn ông nội cậu thì\dots\ là người đã xây dựng ngôi trường này đấy. Nhưng cậu ấy chưa bao giờ dùng điều đó để hạ thấp ai. Thậm chí, bọn tớ còn thấy dễ chịu khi chơi cùng, vì cậu ấy thật sự chân thành.''

Yae, vẫn còn hơi run sau cuộc cãi vã vừa rồi, siết chặt hai tay, rồi khẽ nói: ''Nanase-chan\dots\ ngày trước luôn là người bạn quan trọng nhất của tớ. Bọn tớ lớn lên cùng nhau. Nhưng\dots\ từ vài ngày gần đây, tớ thấy cậu ấy thay đổi. Ban đầu chỉ là dễ cáu gắt hơn một chút, nhưng dần dần\dots\ cậu ấy như thành một người khác vậy.''

''Thậm chí cậu ấy có phần xấu tính hơn nữa\dots\ Nhất là khi Nanase-san đối mặt với Shino-san,'' một cậu học sinh nam lên tiếng, giọng vẻ nghiêm trọng.

Cả lớp khẽ bàn tán phụ họa:

''Tớ cũng để ý thế\dots''

''Từ khi nào nhỉ?''

''Có lẽ\dots\ trùng lúc lớp mình có học sinh mới chuyển tới.''

Câu nói cuối cùng làm Kyuen thoáng khựng lại. Saikyou nhíu mày, còn Reiko thì ngẩng đầu lên, ánh mắt sắc lại.

Một nam sinh khác chêm vào, giọng thì thầm: ''Ừ\dots\ đúng rồi. Bệnh của Nanase-san chỉ thực sự nặng lên\dots\ kể từ lúc Kyuen-kun, Saikyou-san và cả Reiko-san đến lớp.''

Lời khẳng định ấy như một mũi kim chích vào bầu không khí vốn đã căng thẳng. Cậu con trai im lặng, cảm thấy ánh mắt vài người vô thức liếc về phía ba người họ. Không phải là buộc tội, nhưng rõ ràng có một sự liên kết mơ hồ nào đó mà ai cũng nhận ra.

Saikyou nghiêng đầu, khẽ siết lấy tay áo anh trai như muốn xác nhận rằng cô cũng đang thấy sự trùng hợp đáng sợ ấy. Reiko thì chỉ im lặng, song đôi mắt ánh lên vẻ nghi hoặc lẫn trăn trở.

Trong thoáng chốc, câu hỏi treo lơ lửng trong tâm trí cả ba người: ''Tại sao sự xuất hiện của chúng ta\dots\ lại gắn liền với việc Nanase trở nên như thế này?''

Honoka, vốn chẳng giấu nổi sự thẳng thắn của mình, khoanh tay nói: ''Thì rõ ràng là có sự trùng hợp mà. Nanase-chan càng ngày càng nóng tính, còn Shino-san thì\dots\ ừm, mọi người thấy đấy, tự dưng lại chịu mở lòng hơn. Mà cả hai chuyện này đều diễn ra sau khi Reiko-san chuyển trường tới.''

Yae vội xua tay: ''Ê, ê, đừng có nói quá như thế chứ. Reiko-san đâu có làm gì sai đâu\dots\ Chỉ là\dots\ đúng là hơi kỳ thật. Cứ như thể sự xuất hiện của cậu ấy\dots\ đã khuấy động mọi thứ vậy.''

Shiina nhướng mày, giọng nửa đùa nửa thật: ''Cái kiểu `yếu tố ngoại lai' ấy nhỉ? Một người mới đến làm thay đổi cả cục diện. Nghe thì hơi quá, nhưng nhìn Nanase và Shino bây giờ, cũng khó mà không nghĩ tới.''

Không ai trong số họ có ác ý cả, nhưng những lời bàn tán ấy khiến Reiko chỉ biết cắn môi, lòng chợt trào dâng một cảm giác khó gọi tên. Cô nhớ lại ánh mắt Nanase lúc nãy, thứ ánh mắt gần như mất đi sức sống, nhưng lại xoáy chặt vào mình.

Kyuen nhìn lướt qua cô. Cậu không nói gì, nhưng sự khó chịu hiện rõ trong mắt. Saikyou thì khẽ siết chặt nắm tay, như thể muốn phản bác lại tất cả những lời nghi ngờ ấy.

Trong khi đó, không ai nhận ra Sakura, đang đứng lặng ở cuối lớp, dõi theo tất cả. Trong ánh mắt cô thoáng qua một tia suy tư: sự thay đổi này\dots\ thực sự bắt nguồn từ Reiko, hay còn điều gì khác ẩn giấu phía sau?

\begin{center}
  \uline{Chiều hôm đó.}
\end{center}

Trời bỗng sụp xuống một màn mưa như trút nước. Những hạt mưa nện thẳng vào khung cửa kính, để lại những vệt dài chảy xuống, hòa cùng tiếng gió rít từng cơn. Lớp học vốn đã trĩu nặng sau chuỗi sự việc trong ngày nay lại càng ảm đạm, chẳng ai buồn lên tiếng.

Chính lúc ấy, Honoka vỗ mạnh tay xuống bàn, giọng đầy hứng khởi: ''\textbf{Được rồi! Để bầu không khí bớt nặng nề, tối nay chúng ta thử làm một cái thử thách gan dạ đi!}''

Cả lớp lập tức xôn xao. Shiina lập tức tiếp lời, ra vẻ đã bàn bạc trước: ''Thực ra tôi với Reiko cũng có nói chuyện rồi. Tôi thì đồng ý làm, còn cậu ấy thì\dots''

Reiko hơi khựng lại, ánh mắt thoáng dao động. Cô chưa nói gì, chỉ nhìn xuống bàn, bàn tay khẽ siết lại, như thể cô chưa thực sự đồng ý với ý tưởng có phần điên rồ từ cô bạn hoạt bát ấy.

Nanase vừa mới trở về từ phòng y tế, thẳng thừng cất lời: ''Mấy chuyện nhảm nhí đó thì chẳng có gì mà phải sợ. Ngôi trường này là do chính gia đình Furukawa chúng tôi xây dựng, lễ cúng trừ tà trước khi khởi công đã được làm đầy đủ. Không có ma quỷ nào ở đây cả.''

\img{e1-6i}

Honoka chống cằm, cười tinh nghịch: ''Hể\~{}? Thế có nghĩa là\dots\ cậu sợ à\~{}?''

Cô nàng tóc xanh liếc cô một cái sắc lẻm, rồi dõng dạc: ''Tôi sẽ đi. Tôi muốn tận mắt lật tẩy mấy cái tin đồn nhảm nhí này.''

Saikyou nheo mắt, hơi nghi ngờ: ''Liệu cậu có chắc đi nổi không? Trông cậu vẫn chưa khỏe hẳn kia kìa\dots''

Nanase khoát tay: ''Chỉ là đau đầu nhẹ thôi. Tôi vẫn đủ sức.''

Không khí ồn ào dần lên. Yae ôm chặt quyển sách trên tay, nhỏ giọng: ''Nhưng mà\dots\ con gái thì không nên ra ngoài chơi đêm đâu. Làm vậy\dots\ không đúng chút nào\dots''

Honoka bật cười ha hả, vung tay như gạt phăng đi: ''Thôi nào, Yae-chan cứng nhắc quá. Lúc nào cũng luật lệ với lễ nghi. Bớt nghiêm túc đi chút đi.''

Cô nàng Tóc xám mím môi, không chịu thua: ''Không phải cứng nhắc, mà là có chừng mực! Cậu thì lúc nào cũng tùy tiện quá mức thôi, Honoka-san!''

Kyuen ngồi khoanh tay, ánh mắt thoáng đăm chiêu. Cậu đã nghe Honoka nói về cách thức tổ chức hồi nãy, nhưng trong lòng vẫn còn nhiều do dự. Saikyou thì ngược lại, vốn là người hòa đồng, nhưng lần này cũng chẳng thấy an tâm chút nào. Cả hai chỉ gật đầu ngắn gọn: ''Nếu tối nay có thể sắp xếp thì chúng tôi sẽ đi.''

Shino nãy giờ ngồi yên, bỗng cất giọng nhẹ nhàng nhưng dứt khoát: ''Mình sẽ đi. Đây là cơ hội để cả lớp xích lại gần nhau hơn.''

Ánh mắt cô chuyển sang Reiko, mỉm cười khích lệ: ''Cậu đi cùng nhé, Reiko-san?''

Reiko thoáng chần chừ, nụ cười nhạt hiện trên môi nhưng đôi mắt lại đảo đi như đánh lạc hướng: ''\dots Mình\dots\ có lẽ sẽ xem xét.''

Không ai ép buộc thêm. Honoka đảo mắt nhìn quanh, rồi dừng lại ở Sakura, người đang đứng ở bên cửa, hầu như không tham gia cuộc trò chuyện này.

''Sakura-chan, còn cậu?''

Sakura chậm rãi ngẩng đầu, giọng bình thản: ''Tôi không có ý kiến phản đối.''

Cuộc trò chuyện tưởng chừng đã khép lại, thì Nanase bất ngờ liếc về phía Shino và Reiko, đôi mắt tối sầm: ''Misora-san\dots\ cậu ta có vẻ như lúc nào cũng kè kè bên Reiko-san nhỉ.''

\img{e1-6j}

Honoka lập tức chen ngang, cười hồn nhiên: ''Tớ thấy có vấn đề gì đâu. Phải nên mừng cho cậu ấy mới đúng chứ.''

Yae gật đầu theo, nói với sự chân thành: ''Đúng thế. Hồi trước Shino-san nhút nhát lắm, bây giờ mở lòng hơn rồi\dots\ Dù là chỉ với Reiko-san, nhưng cũng tốt còn gì.''

Nanase nhếch môi, giọng như nén lại: ''\dots Đúng là phiền phức thật.''

Shiina nhăn mặt, đáp thẳng: ''Thôi nào, mặc kệ người ta đi. Chẳng qua cậu đang ghen vì cậu cũng thích Reiko đấy thôi.''

''Im đi,'' cô gắt lên, bàn tay siết chặt. ''Tôi không khiến. Nhưng cứ nhìn Misora-san như vậy, tôi lại càng nhức đầu\dots\ lại càng ngứa mắt.''

Saikyou ở bên cạnh, liền nghiêm giọng: ''Nanase-san, cậu quá đáng với Shino-san rồi. Nếu thấy ngứa mắt thì đừng nhìn nữa.''

Nanase quay phắt sang, giọng lạnh lẽo: ''Cậu ngậm miệng vào cho tôi.''

Kyuen buông một tiếng thở dài, lắc đầu bất lực: ''Đúng là\dots\ chả khác nào một công chúa ương ngạch cả.''

Bầu không khí trong lớp căng thẳng một hồi, rồi lại trở về với guồng quay thường nhật của trường học. Những tiếng mưa ngoài trời tiếp tục rơi rả rích, như thể cũng đang hòa theo vào bản trường ca của sự thường ngày ấy.

Ở cuối lớp, Sakura lặng lẽ dõi theo. Ánh mắt cô toát lên một sự phức tạp: một thoáng lo âu khi nhìn Nanase - như thể cô bạn đang gánh chịu một cơn rối loạn tinh thần nào đó; một thoáng nghi hoặc khi dừng lại ở Kyuen và Saikyou - hai con người vừa xuất hiện mà đã bị cuốn vào tâm điểm xáo trộn. Trong đáy mắt Sakura, thứ ánh sáng lạnh lẽo ấy như thể đang dò xét một bí mật không thể nói ra.

\img{e1-6k}

\ct

Chuông báo hết tiết vang lên, kéo dài và lan ra khắp dãy hành lang. Cả lớp bắt đầu thu dọn sách vở. Ngoài trời, cơn mưa đã ngớt dần, chỉ còn lại những giọt lác đác rơi tí tách trên mái ngói và lan can, không giống như những trận mưa dai dẳng từ những ngày vừa qua.

Honoka vẫn giữ nguyên vẻ tươi cười đầy tinh nghịch, vỗ tay đánh bốp một cái: ``Thế nhé! Tối nay tớ mong chờ lắm đó. Nhất định sẽ vui phải biết!''

Một vài ánh mắt lập tức hướng về cô.

Yae khẽ chau mày, ôm chặt cặp xách: ``\dots Mình không thấy vui chút nào. Nhưng mà\dots\ mình sẽ đi. Ít nhất để còn canh chừng mọi người.''

Shiina lại cười rạng rỡ: ``Tất nhiên là mình đi! Lâu rồi mới có cơ hội thế này, háo hức luôn!''

Nanase đứng dậy, khoác cặp lên vai, giọng trầm thấp: ``Tôi chẳng muốn\dots\ nhưng bị lôi vào rồi thì tôi sẽ đi. Để lật tẩy mấy cái lời đồn vớ vẩn kia.''

Sakura ngồi ở bàn cuối, không biểu hiện gì, chỉ khẽ gật đầu.

Shino thì lại ngược lại, gần như là người duy nhất toát ra vẻ mong chờ thật sự. Gương mặt cô hồng lên, nụ cười ngượng ngập nhưng chân thành: ``Mình nghĩ\dots\ sẽ thú vị lắm. Mình cũng muốn thử.''

Ba người nhà Haragen vẫn chưa cho câu trả lời rõ ràng.

Kyuen chỉ hờ hững: ``Xem tình hình rồi sẽ tính.''

Saikyou mỉm cười xã giao, nhưng trong giọng lại thiếu đi sự chắc chắn thường thấy: ``Nếu sắp xếp được thì chúng tôi sẽ tham gia.''

Còn Reiko, như mọi khi, vẫn im lặng, chỉ đáp bằng một cái gật đầu mơ hồ.

Lần lượt từng người rời khỏi lớp học, để lại căn phòng trống trải dần, chỉ còn tiếng ghế kéo lạch cạch và tiếng mưa cuối buổi chiều đang dần tạnh.

Khi bóng dáng cuối cùng khuất khỏi hành lang, cô nàng mắt hai màu vẫn chưa rời đi. Cô bước chậm lại, tìm tới chỗ Kyuen và Saikyou, hai người vẫn đang đứng bên khung cửa sổ nhìn trời.

``Về cái thử thách tối nay\dots'' cô ngập ngừng, giọng nhỏ dần, ``hai người có định đi không? Liệu\dots\ nó có an toàn không?''

Kyuen liếc sang em gái, rồi quay lại nhìn Reiko. ``Không hẳn là không an toàn. Nhưng\dots\ có gì đó khiến tôi thấy cấn cấn.''

Saikyou gật đầu tán đồng, nụ cười tắt đi. ``Ngay khi nghe Honoka-san nhắc đến `thử thách gan dạ', tôi đã thấy\dots\ quen quen. Như thể đã từng nghe hoặc từng trải qua ở đâu đó.''

Reiko khẽ mở to mắt: ``Quen thuộc\dots\ nghĩa là sao?''

Hai anh em đồng loạt lắc đầu.

``Khó nói lắm,'' cậu con trai khẽ thở dài. ``Nó mơ hồ, không rõ ràng. Nhưng linh cảm cho tôi biết\dots\ nó không hề là trò đùa.''

``Thậm chí là\dots\ tôi còn có thể cảm thấy một sự chết chóc ở đây,'' cô em gái tiếp lời, giọng trầm hẳn xuống.

Một làn gió lạnh bất chợt lùa vào lớp qua khe cửa, làm cánh cửa sổ cũ kỹ lay động như những bóng ma vô hình. Reiko rùng mình, không chỉ vì cái lạnh. Những từ ``thử thách gan dạ'' và ``chết chóc'' cứ văng vẳng bên tai, khiến cô không khỏi liên tưởng đến những lời đồn về ``lời nguyền trong phòng Mỹ thuật'' mà các bạn trong lớp thường xì xào với nhau. Bầu không khí trong lớp chùng xuống, nặng trĩu bởi một thứ gì đó u ám và rợn người, như thể có một bóng đen vô hình đang len lỏi giữa những dãy bàn học.

Bầu không khí nặng nề ấy chỉ chấm dứt khi giọng một người vang lên phía sau.

``Hai cậu. Tôi có chuyện cần nói riêng.''

Sakura đứng đó từ lúc nào, ánh mắt bình tĩnh nhưng lạnh lẽo, chiếu thẳng vào hai anh em. Kyuen và Saikyou nhìn nhau, rồi đồng loạt gật đầu, bước theo cô nàng Vu nữ ra hành lang.

Reiko đứng chết lặng tại chỗ. Lý trí mách bảo cô nên về ngay, nhưng trái tim lại trĩu nặng băn khoăn. Cảm giác đó - cảm giác về sự khác thường nơi hai anh em Haragen - chợt ùa về, giống như cái ngày đầu tiên cô gặp họ, trước cả khi nhập học chính thức.

Cô lặng lẽ xách túi, bước thật khẽ để bám theo từ phía xa.

\ct

Reiko bước thật khẽ ra khỏi lớp, giữ một khoảng cách vừa đủ với ba bóng người đang đi trước. Kyuen và Saikyou đi cạnh nhau, còn Sakura dẫn đường, dáng vẻ vẫn bình thản nhưng bước chân dứt khoát. Tiếng giày của họ dội lên nền gỗ cũ, tiếng cọt kẹt vang vọng qua dãy hành lang vắng lặng.

Cô không dám tiến lại quá gần, chỉ men theo hàng cửa sổ. Thỉnh thoảng, khi ánh đèn dầu treo trên tường nhấp nháy, Reiko lại giật thót, sợ cái bóng của mình sẽ bị phát hiện.

Rồi bất chợt, Saikyou ngoảnh lại.

Reiko tim đập thình thịch, vội ép sát người vào bức tường nơi ô cửa. Chiếc rèm vải thô màu xám treo lửng che đi một nửa người cô, chỉ còn để lộ đôi mắt. Saikyou nhìn quanh, ánh mắt dò xét, rồi quay đi.

``Mình tưởng có ai đó đang đi theo chúng ta đấy chứ\dots'' cô nàng mắt xanh lẩm bẩm.

Cô nàng mắt hai màu thở ra khẽ khàng, mồ hôi lạnh rịn ra tay.

Hành lang dần mở ra không gian rộng hơn: hội trường lớn của trường Aogami. Nơi đây vốn được xây dựng để tổ chức lễ nghi, nền nhà lót gạch vuông đen trắng đã sờn, những tấm rèm nhung dày phủ bụi từ thời chiến. Trên trần cao, đèn treo thắp bằng khí đốt tỏa ánh vàng u ám, lay động theo từng luồng gió lùa. Tiếng mưa nhỏ giọt qua mái ngói hằn vết cũ nghe càng làm không gian thêm trống trải.

Reiko lặng lẽ bước theo mép tường, giấu bước chân trong tiếng mưa lách tách ngoài hiên. Từng nhịp tim của cô như đập cùng nhịp với tiếng gõ giày của ba người phía trước.

Cuối cùng, cả ba đi ra cửa sau hội trường, nơi dẫn ra khoảng sân hẹp sát khu rừng Aogami. Ngoài kia, cây cối ướt mưa, từng thân cây đen sẫm lấp loáng trong ánh chạng vạng. Gió mang theo mùi ẩm mốc của đất rừng và tiếng lá xào xạc.

Cô biết nếu bây giờ bước ra, chắc chắn sẽ bị lộ. Cô đảo mắt quanh, tìm chỗ ẩn. Ngay bên hành lang phụ có một tủ gỗ đựng dụng cụ thể dục, cửa hé mở. Cô len người vào, vừa đủ để nửa cánh tủ khép lại, để lại một khe nhỏ đủ nhìn ra ngoài. Bụi gỗ và mùi mốc ngai ngái xộc vào mũi, nhưng cô buộc phải nín thở.

Từ khe hở, cô thấy Sakura dừng lại dưới mái hiên, quay người đối diện với hai anh em. Ánh đèn vàng hắt lên khuôn mặt cô, sáng tối chập chờn. Giọng Sakura vang lên, trầm tĩnh nhưng chứa một sức nặng kỳ lạ:

``Trước khi nói tiếp\dots\ có một điều tôi muốn các cậu trả lời thành thật. Hai người\dots\ thực sự là đến từ đâu vậy?''

Reiko siết chặt mép tủ, trái tim đập loạn. Cô hiểu, khoảnh khắc này sẽ lật tung lớp vỏ bọc vốn khiến mình hoang mang từ lâu.

\ct

Sakura đứng đối diện với hai anh em nhà Haragen, ánh mắt cô bình thản nhưng trong đáy mắt lại lấp lánh một sự nghi ngờ đã được nung nấu từ lâu. ``Trước khi nói tiếp\dots\ có một điều tôi muốn các cậu trả lời thành thật. Hai người thực sự đến từ đâu vậy?''

Câu hỏi rơi xuống, im lặng bao trùm. Cả Saikyou lẫn Kyuen đều khẽ giật mình, ánh mắt họ lóe lên trong thoáng chốc, rồi nhanh chóng che giấu. Cô em gái hơi nghiêng đầu, cố tỏ vẻ bỡ ngỡ: ``Ý cậu là sao vậy, Sakura-san?''

Cậu anh trai cũng phụ họa, giọng trầm ổn như thể chính cậu cũng không hiểu vấn đề: ``Chúng tôi chẳng phải cũng giống như mọi người thôi sao? Là học sinh mới đến trường này, có gì lạ đâu?''

Sakura vẫn giữ im lặng trong vài giây, rồi chậm rãi tiến thêm một bước, giọng cô đều đặn: ``Lạ chứ. Cách nói chuyện của hai người có gì đó\dots\ khác biệt. Không phải giống kiểu nói của những bạn trong lớp. Nhiều lần nghe, tôi cảm giác như giọng điệu và cách chọn từ của hai người\dots\ \emph{không hề thuộc về nơi này}.''

Saikyou mỉm cười gượng: ``À\dots\ chúng tôi vốn từ một nơi xa đến. Ngôn ngữ, phong cách có chút lạ, chuyện đó cũng bình thường thôi.''

Nhưng cô nàng tóc hồng chưa dừng lại. Cô nghiêng đầu, hạ giọng: ``Và còn những từ ngữ mà hai người dùng. Tôi chưa từng nghe bao giờ, nhưng vẫn mơ hồ hiểu được nghĩa. Những từ đó\dots\ \emph{không phải của thời đại này}.''

Kyuen lập tức chen vào, vẻ điềm tĩnh nhưng giọng lại hơi cứng: ``Có thể chỉ là tiếng địa phương thôi. Vùng quê của chúng tôi cũng có nhiều cách nói khác lạ mà. Chẳng có gì to tát cả.''

Đôi mắt xanh lá mạ của Sakura nheo lại. Cô vẫn bình thản, song lời nói như mũi dao lạnh lẽo đâm thẳng: ``Nếu chỉ là `tiếng địa phương'\dots\ vậy cậu giải thích thế nào về cái này?''

Cô chỉ tay về phía túi áo của Kyuen. Bên mép vải lấp ló một vật dài mảnh màu đen. Sợi dây sạc mà hôm trước cậu tình cờ tìm thấy trong lọ hoa bỉ ngạn đỏ.

Cậu con trai bất giác khựng lại, bàn tay vô thức che lên túi áo. Trong khoảnh khắc đó, đồng tử đỏ của cậu co lại. ``S-Sao lại\dots''

Sakura không rời mắt: ``Tôi chưa từng thấy vật gì như vậy. Không ai trong trường hay trong thành phố này từng nhắc đến nó. Nhưng tôi nhớ\dots\ khi tìm thấy, cậu đã lẩm bẩm rằng `thứ đó không thể tồn tại ở thời đại này'.'' Ngay khi cô dứt câu, bầu không khí như đông cứng. Saikyou nín thở, mắt cau lại, còn Kyuen im lặng không đáp.

Cô nàng tóc anh đào bước thêm một bước nữa, giọng lạnh nhưng chắc: ``Và hôm trước, Saikyou-san\dots\ chính miệng cậu đã nói: `nơi này khác hẳn so với thời kỳ của mình'. Hai người nghĩ tôi không nghe thấy ư?''

Saikyou cứng người, gương mặt lộ rõ vẻ bất an. Kyuen thoáng quay sang cô em gái, như muốn cảnh báo cô giữ im lặng, nhưng đã quá muộn. Không còn đường lui, họ nhìn nhau, như trao đổi trong im lặng một quyết định cuối cùng.

Cuối cùng, chính Kyuen thở dài, giọng trầm xuống, không còn né tránh nữa: ``\dots Được rồi. Chúng tôi không thể giấu nữa. Sakura-san, cậu nói đúng.''

Saikyou cũng gật nhẹ, ánh mắt pha lẫn sự bất lực và thừa nhận: ``Chúng tôi\dots\ đến từ tương lai. Chính xác là từ năm 2018.''

Trong tủ gỗ gần đó, Reiko cũng nghe thấy hết mọi thứ. Tim cô đập thình thịch, hai tay nắm chặt vạt áo đồng phục. Những lời Sakura vừa vạch ra khiến cô kinh hãi: từ việc cô vạch ra được một vật dụng tới từ tương lai, tới việc nguồn gốc thật sự của hai người mà cô đã cho tạm trú. Cô bỗng nhớ lại cảm giác kỳ lạ ngay lần đầu gặp hai anh em: một thứ gì đó không thuộc về thế giới này. Giờ thì cô đã hiểu, sự linh cảm đó không hề sai.

Một khoảng lặng đè nặng xuống, chỉ còn tiếng gió rít qua tán cây.

Sakura sững sờ, đôi mắt mở to: ``2018? Nghĩa là\dots\ tận hơn bảy mươi năm sau ư?''

Trong giọng nói của cô có sự kinh ngạc rõ rệt, nhưng không phải hoàn toàn là sợ hãi. Nó giống như cú sốc khi một giả thuyết mơ hồ trong đầu bỗng nhiên được chứng thực.

Reiko, đang trốn, siết chặt bàn tay đến mức móng tay hằn lên da. Tâm trí cô quay cuồng - hai người bạn học bí ẩn này, hóa ra thật sự không thuộc về thời đại này. Câu trả lời quá đỗi phi lý\dots\ nhưng lại vang lên chân thật hơn bao giờ hết.

Kyuen nắm chặt bàn tay, cảm giác như gánh nặng giấu kín bấy lâu đã đổ ập xuống. Cậu không sợ bị phát hiện, nhưng lại sợ rằng sự thật ấy sẽ làm thay đổi tất cả. ``\textit{Liệu Sakura có tin không, hay sẽ xem mình là kẻ bịa chuyện điên rồ?}''

Saikyou thì khẽ cắn môi. Cô luôn là người nhạy cảm hơn, và giờ đây trái tim cô đập loạn nhịp, vừa lo sợ, vừa kỳ lạ nhẹ nhõm. ``\textit{Ít nhất\dots\ bọn mình không còn phải nói dối nữa.}''

Sakura thì lại khác. Trong đôi mắt cô, ngoài sự kinh ngạc còn ánh lên một tia sáng lạ thường, như thể trực giác lâu nay đã mách bảo rằng hai người này không hề bình thường. ``\textit{Hóa ra linh cảm của mình\dots\ đúng thật. Hai anh em Haragen này mang theo một sức mạnh, một dấu ấn từ thế giới khác.}''

Và rồi, cô khẽ liếc sang một góc khuất. Đôi môi cô nhếch thành một nụ cười nhẹ: ``\dots Và cả Reiko-san nữa.''

Cả Kyuen lẫn Saikyou đều giật mình, còn trong tủ gỗ, Reiko đông cứng. Khi Sakura chỉ thẳng vào nơi cô đang trốn, cô chỉ biết thót tim, chậm chạp bước ra ngoài với vẻ ngượng nghịu.

``C-cậu\dots\ biết từ khi nào vậy?'' người bước từ trong tủ ra lắp bắp.

``Ngay từ lúc rời lớp học.'' Sakura đáp, giọng nhẹ nhàng mà chắc nịch. ``Cách cậu bước chân không lẫn đi đâu được.''

Reiko cúi đầu, mặt đỏ ửng. Nhưng ánh mắt cô lại vô tình lướt qua hai anh em Haragen. Trong khoảnh khắc ấy, cô cũng cảm nhận được sự kỳ lạ, như một mạch lực vô hình đang kết nối cả ba người họ.

Sakura dừng lại, chậm rãi nói: ``Thú vị thật đấy. Ba người cùng chung một họ hiếm hoi, `Haragen'. Và giờ thì, cả ba đều tụ tập ở đây\dots\ hẳn không phải là trùng hợp.''

Bầu không khí như càng trở nên đặc quánh. Cô khoanh tay, thẳng thắn: ``Được rồi. Kyuen-san, Saikyou-san, tôi muốn nghe tất cả. Không giấu diếm nữa.''

Kyuen hít sâu, ánh mắt anh thoáng đượm nỗi buồn xen lẫn ám ảnh khi nhớ lại. Saikyou nhìn anh trai, rồi khẽ gật.

``Chúng tôi vốn dĩ là học sinh của trường Aogami, nhưng ở năm 2018,'' cậu bắt đầu, giọng đều đặn. ``Vào một ngày định mệnh, một trận động đất dữ dội đã khiến trường sụp đổ. Chúng tôi\dots\ may mắn sống sót, nhưng khi chạy đến hành lang, chúng tôi đã gặp một bóng dáng hồn ma.''

Reiko mở to mắt. Cơ thể cô run lên, chi tiết ấy trùng khớp với giấc mơ của cô.

``Người đó\dots\ đã kéo chúng tôi đến nơi này.'' Saikyou tiếp lời, khẽ siết chặt tay áo. ``Để `tìm hiểu quá khứ của cô ấy'. Chúng tôi không nói tên, nhưng chắc cậu cũng biết\dots\ người đó là ai.''

Trong im lặng, một hình ảnh hiện lên trong đầu cả ba: Shino, với nụ cười mơ hồ giữa hành lang tăm tối.

Kyuen tiếp tục: ``Trước cả khi nhập học, chúng tôi đã mơ. Giấc mơ đưa chúng tôi đến cả trường Aogami hiện đại lẫn trong quá khứ. Trong mơ\dots\ chúng tôi gặp lại hồn ma ấy. Và rồi, khi tỉnh dậy\dots\ mọi thứ bắt đầu.''

Không khí căng thẳng đến mức Reiko phải đặt tay lên ngực để giữ nhịp tim. Rồi như bị thôi thúc, cô cũng cất lời: ``\dots Mình cũng từng mơ. Một giấc mơ kỳ lạ. Ở đó, mình bị kẹt trong trường Aogami vào ban đêm, và\dots\ mình đã thấy một hồn ma, giống hệt Shino-san.''

Cả hai anh em thời hiện đại đồng loạt giật mình nhìn sang Reiko. Trong khoảnh khắc ấy, ba đôi mắt giao nha: ba con người, cùng một họ, cùng một giấc mơ, cùng một sợi dây định mệnh vô hình.

Sakura phá tan bầu không khí ngột ngạt bằng một nụ cười mơ hồ, giọng trêu chọc nhưng lại gợi ra sự bất an: ``Hóa ra\dots\ cả ba đều có cùng một giấc mơ, cùng nhìn thấy cùng một người? Thú vị thật. Chắc bởi vì cả ba cùng họ Haragen đấy mà\~{}''

Reiko thoáng đỏ mặt, còn Kyuen và Saikyou nhìn nhau với ánh mắt nghiêm trọng, như thể cô ấy đã vô tình nói trúng một bí mật mà họ chưa muốn nghĩ tới.

Sakura chợt rút ra từ túi áo hai món đồ lạ. Cô giơ lên, để ánh chiều tà hắt qua: một chiếc dock sạc hình tròn với đèn nhỏ ở mặt trước, bề mặt nổi lên hai khe cắm và một gờ bán ô-van làm điểm tựa; phía sau là một cổng Type-C. Vật kia là một củ sạc màu đen, quen thuộc với bất kỳ ai trong thế giới 2018.

``Đây,'' Sakura khẽ mỉm cười. ``Tôi tìm thấy trong phòng Mỹ thuật, lúc trốn tiết Thể dục hôm kia. Linh cảm mách rằng hai người sẽ cần đến nó.''

Kyuen và Saikyou sững sờ. Họ lập tức nhận ra những vật dụng vốn quen thuộc ở thời đại của mình. Cậu con trai cẩn thận nhận lấy dock sạc, còn cô em gái nâng củ sạc trong tay, ánh mắt rưng rưng một thoáng hoài niệm.

``Cảm ơn\dots'' Kyuen nói khẽ, giọng thành thật.

``Ừm\dots\ bọn tôi thật sự biết ơn\dots'' Saikyou thì thầm.

Reiko nghiêng đầu, nhìn hai món đồ lạ lẫm kia. Trong ánh mắt cô, vừa là sự bối rối, vừa là cảm giác mơ hồ.Chúng vừa quen thuộc, vừa hoàn toàn xa lạ.

Bầu không khí lặng im một lần nữa, nhưng lần này không còn căng thẳng. Nó như báo hiệu rằng từ giây phút này, cả bốn người đã cùng bước chân vào một vòng xoáy bí ẩn mới.

\ct

Mặt trời đỏ rực chậm rãi chìm dần xuống phía chân trời, nhuộm bầu trời Aogami thành một dải màu hoàng kim pha tím thẫm. Những tán cây cổ thụ quanh trường hắt bóng dài lên nền gạch ẩm ướt sau cơn mưa, tạo ra những vệt đen loang lổ trông như những cánh tay vươn ra, bám lấy tòa nhà cũ kỹ. Tiếng ve cuối mùa réo rắt xen lẫn tiếng gió thổi qua khung cửa sổ kẽo kẹt, khiến khung cảnh mang một vẻ đẹp vừa tĩnh lặng, vừa len lỏi đâu đó sự bất an.

Cả bốn người họ bước ra khỏi khu sân sau, men theo hành lang vắng. Trên đường đi, Sakura chợt cất tiếng:

``Vậy\dots\ tối nay, thử thách gan dạ. Ba người có tham gia không?''

Câu hỏi tưởng chừng nhẹ nhàng, nhưng lập tức làm bầu không khí chùng xuống.

Kyuen khẽ nhíu mày. Trong đầu cậu lại vang lên những mảnh vụn ký ức: mùi bụi vữa từ vụ động đất, ánh sáng nhập nhòe trong giấc mơ, và bóng ma kéo họ vào đây. Cảm giác ``chết chóc'' bủa vây từ khi nghe Honoka nhắc tới ``thử thách gan dạ'' vẫn còn nguyên. Cậu siết chặt tay, im lặng không trả lời ngay.

Saikyou thì nhìn xuống nền gạch loang lổ, môi khẽ mím lại. Cô vốn tò mò, muốn biết sự thật sau những lời đồn ma ám, nhưng linh cảm mách rằng nơi này ẩn chứa điều gì đó không thể xem nhẹ.

``\textit{Nếu đi, liệu chúng ta có còn nguyên vẹn để trở về?}''

Reiko thì lại khác. Trong tim cô, hình bóng Shino vẫn chiếm chỗ lớn nhất. Cô lo lắng cho bạn mình - người vừa mới chịu mở lòng - nhưng đồng thời, một cảm giác khó tả thôi thúc cô phải đứng ngoài. Một mặt, Reiko muốn Shino dần gắn kết với các bạn khác thay vì chỉ dựa dẫm vào mình. Mặt khác, chính sự hiện diện của Kyuen và Saikyou khiến cô bất giác thấy sợ hãi, như thể có một màn sương đen bao quanh hai người ấy, lạnh lẽo và nguy hiểm.

Ba người lặng lẽ đi trong ánh hoàng hôn, mỗi người chìm trong dòng suy nghĩ riêng. Sau một hồi đấu tranh, cuối cùng Kyuen lên tiếng, giọng trầm nhưng dứt khoát: ``\dots Được. Tôi sẽ đi. Nhưng chỉ để quan sát.''

Saikyou thở dài, rồi cũng gật đầu theo anh trai: ``Ừ. Tôi cũng vậy. Không thể ngó lơ\dots\ dù thấy chẳng an toàn chút nào.''

Ánh mắt họ lặng lẽ thoáng chạm nhau, trong đó vừa có tò mò, vừa có cảnh giác. Họ đều nhớ đến lời đồn về “con ma trong phòng Mỹ thuật”. Dù tò mò muốn tận mắt chứng kiến, nhưng cả hai lại có cảm giác như đang bước vào bẫy.

Reiko im lặng thêm một lúc lâu mới khe khẽ cất lời: ``Mình thì\dots\ sẽ không vào. Mình chỉ đứng bên ngoài trường thôi. Vừa để xem các cậu ra sao, vừa để Shino-san có cơ hội dựa vào những người khác ngoài mình. Và\dots'' cô ngập ngừng, rồi hạ giọng ``cũng bởi vì, ở cạnh mọi người, mình cũng cảm thấy\dots\ không cần thiết.''

Cậu con trai thoáng ngạc nhiên, nhưng không phản bác. Cô em gái chỉ cúi đầu, vẻ hơi chua chát, cảm tưởng như bị cô ấy phản bội lòng tin.

``M-Mình không có ý nói xấu hai cậu đâu,'' Reiko vội lên tiếng phản bác, múa máy quờ quạng. ``Nhưng mà ấy\dots''

``Nếu cậu cảm thấy không ổn thì cậu không cần phải đi đâu,'' Saikyou nói, vỗ vai cô nàng. ``Nếu có gì bất thường thì bọn tôi sẽ chạy ra ngay, được chứ?''

``\dots Mình biết rồi. Trăm sự nhờ các cậu.'

Sakura là người phá tan sự im lặng sau đó. Cô đưa ánh mắt sắc sảo nhìn cả ba, rồi khẽ mỉm cười: ``Tốt. Vậy thì quyết định thế đi. Gặp nhau ở cổng trường lúc chín giờ tối nay.''

Nói rồi, bốn người tách ra, mỗi người theo một hướng về nhà trọ của mình. Hoàng hôn tiếp tục buông xuống, nuốt chửng sân trường Aogami vào trong bóng tối.

\pov{Haragen Reiko}

Trời dần ngả tối khi chúng tôi rời khỏi trường. Những con đường đất ẩm ướt sau cơn mưa khiến bước chân cứ nặng nề, và gió thu lạnh lẽo len qua từng kẽ áo. Tôi đi cạnh Kyuen-san và Saikyou-san, nhưng trong lòng thì rối bời đến mức chẳng biết phải nghĩ gì nữa.

Một bên\dots\ là thử thách gan dạ tối nay. Tôi đã quyết rồi: sẽ đi, nhưng chỉ đứng ngoài cổng trường. Một phần để xem Shino-san có thể hòa nhập với các bạn khác hay không, một phần\dots\ vì tôi không thể phủ nhận sự đáng sợ toát ra từ cặp song sinh này.

Một bên khác\dots\ là bí mật kinh hoàng tôi vừa nghe lén được. Kyuen-san và Saikyou-san, những người chỉ mới đến lớp chúng tôi chưa lâu\dots\ lại thừa nhận rằng họ đến từ tương lai. Hơn bảy mươi năm sau. Tôi vẫn không thể nào tin nổi. Mọi thứ như một cơn ác mộng không có hồi kết.

Chúng tôi im lặng khá lâu, cho đến khi Kyuen bất ngờ lên tiếng, giọng trầm và khẽ khàng: ``\dots Xin lỗi. Vì đã giấu cậu.''

Tôi quay sang nhìn, bắt gặp ánh mắt nghiêm túc của cậu ấy. Saikyou cũng gật đầu, thêm vào: ``Chúng tôi không nói, vì\dots\ nếu thẳng thắn nói rằng mình từ tương lai đến, cậu sẽ tin sao? Hay cậu sẽ nghĩ bọn tôi bị điên?''

Tôi cảm thấy bối rối. Thực ra, nếu là vài ngày trước, có lẽ tôi đã gạt phăng đi như một trò đùa vớ vẩn, thậm chí còn cười xòa vì lý do nghe rất phi lý kia. Nhưng sau tất cả những gì đã xảy ra\dots\ tôi đâu còn chắc chắn vào điều gì nữa.

``Không sao đâu,'' tôi lên tiếng. ``Mình có lẽ cũng sẽ bấn loạn như hai cậu nếu mình cũng rơi vào tình cảnh đó mà.''

Tôi thở dài, mắt ngước lên trên bầu trời.

Mưa đã tạnh, để lại một bầu trời xám xịt trên thị trấn Aogami. Những tia nắng cuối ngày yếu ớt len lỏi qua những đám mây dày đặc, hắt xuống những mái ngói ướt át của những ngôi nhà cổ kính. Không khí ẩm ướt vẫn còn đọng lại, mang theo mùi của đất và cỏ sau cơn mưa.

Thị trấn như được gột rửa sạch sẽ, những con đường lát đá bóng loáng phản chiếu ánh hoàng hôn nhợt nhạt. Xa xa, những ngọn đồi phủ rừng rậm vẫn còn chìm trong màn sương mỏng, tạo nên một khung cảnh mờ ảo như tranh thủy mặc.

Ngoài kia, những đèn đường bắt đầu được thắp lên, từng chấm sáng nhỏ xuyên qua màn sương chiều, như những ngôi sao le lói trên mặt đất. Thị trấn Aogami dần chìm vào giấc ngủ, nhưng những câu hỏi trong lòng tôi vẫn còn thao thức, chưa tìm được lời giải đáp.

Nhưng trong lòng tôi lúc này lại đang cuộn trào bao cảm xúc khó tả. Sự xuất hiện đột ngột của hai anh em nhà Haragen khiến tâm trí tôi rối bời. Cái họ ``Haragen'' ấy cứ văng vẳng trong đầu như một giai điệu quen thuộc mà không sao nhớ ra được nguồn cội. Tôi tự hỏi liệu đây có phải là định mệnh, hay chỉ đơn thuần là một sự trùng hợp kỳ lạ.

Tôi thở dài, hơi thở phả lên mặt kính cửa sổ tạo thành một vệt mờ nhỏ. Ngón tay vô thức vẽ những đường ngoằn ngoèo trên đó, như thể đang cố gắng vẽ ra manh mối cho những câu hỏi đang quẩn quanh trong đầu. Giờ tôi đã biết họ là ai, đến từ đâu.

Nhưng trong lòng tôi vẫn còn thôi thúc một điều: Tại sao tôi lại cảm thấy gần gũi với họ đến thế?

``Mình\dots'' tôi hít sau, rồi chậm rãi hỏi, ``mình còn một chuyện. Hai người từng nói về dự cảm không lành… về thử thách gan dạ tối nay. Vậy rốt cuộc là gì?''

Hai anh em liếc nhìn nhau, ánh mắt thoáng lưỡng lự, trước khi Saikyou thì thầm: ``Không rõ ràng\dots\ chỉ là\dots\ chúng tôi từng nghe đâu đó trong giấc mơ. Về đêm tối trong trường học, và một thứ gì đó khủng khiếp chờ đợi.''

Tôi khựng lại. Giấc mơ. Từ đó vừa thốt ra đã khiến tim tôi đập mạnh. Phải, tôi cũng đã từng mơ\dots\ một giấc mơ mờ mịt và rùng rợn. Nhưng làm sao tôi có thể nói điều đó ra lúc này?

Chúng tôi lại rơi vào im lặng. Tiếng bước chân vang lên đơn độc trên con đường vắng, hòa vào tiếng gió u u thổi qua. Tôi nhìn bóng lưng hai người đi cạnh, cảm thấy giữa chúng tôi tồn tại một bức màn vô hình, vừa gần gũi vừa xa lạ.

``\textit{Mình đã bị cuốn vào chuyện gì thế này\dots?}''

Một thử thách gan dạ tưởng chừng vô hại\dots\ hay là một nghi thức để gọi lên điều gì đó đáng sợ hơn? Kyuen-san và Saikyou-san thực sự đến từ tương lai\dots\ nhưng vì sao lại là trường Aogami? Vì sao lại trùng khớp với giấc mơ của tôi? Và\dots\ Shino-san trong giấc mơ kia rốt cuộc là ai?

Hàng loạt câu hỏi xoáy sâu trong đầu, không có lời giải. Tôi chỉ biết siết chặt tay, lặng lẽ bước đi trong buổi chiều tà, khi bầu trời đã hoàn toàn nhuộm màu tím sẫm.

\end{document}